\chapter{Эллиптические кривые}\label{chap:elliptic-curves}
\label{chap:elliptic_curves}
%references http://infosec.pusan.ac.kr/wp-content/uploads/2019/09/Pairings-For-Beginners.pdf
Вообще говоря, эллиптические кривые — это геометрические объекты в проективных плоскостях (см. \secname{} \ref{sec:planes}) над данным полем, состоящие из точек, удовлетворяющих определённым уравнениям. Одна из их ключевых особенностей с точки зрения криптографии заключается в том, что если базовое поле имеет положительную характеристику, эллиптические кривые образуют конечные коммутативные группы (\secname{} \ref{sec:finite-groups}). В криптографических приложениях мы обычно работаем в большой циклической подгруппе такой группы кривой. Кроме того, считается, что в этом случае задача \concept{дискретного логарифма} (\secname{} \ref{def:DL-secure}) для многих групп эллиптических кривых является трудной, при условии что базовая характеристика достаточно велика.\footnote{Углублённое введение в эллиптические кривые дано, например, в \cite{silverman-1994}. Введение с криптографической точки зрения дано в \cite{hoffstein-2008}.}

Особый класс эллиптических кривых — так называемые кривые, дружественные к спариванию (pairing-friendly curves), для которых определено понятие группового спаривания (\secname{} \ref{pairing-map}), обладающего криптографически полезными свойствами.

В этой главе мы вводим эллиптические кривые так, как они используются в подходах к построению SNARK на основе спариваний. Рассматриваемые нами эллиптические кривые определены над простыми полями или расширениями простых полей, то есть мы в значительной степени опираемся на понятия и обозначения из \chaptname{} \ref{chap:algebra}. 

\section{\concept{Кривые в короткой форме Вейерштрасса}}
\label{sec:short_weierstrass_curve}
В этом разделе мы вводим \term{\concept{кривые в короткой форме Вейерштрасса} (short Weierstrass curves)}, которые являются наиболее общим типом кривых над конечными полями характеристики больше $3$ (см. \chaptname{} \ref{def:characteristic}), и начинаем с их так называемой \term{аффинной репрезентации (affine representation)}.\sme{do we still need the explanation of projective planes from the previous chapter? or can they be moved to this one?}\footnote{Введение эллиптических кривых в аффинной репрезентации, вероятно, не самый распространённый и концептуально наиболее чистый способ, но мы считаем, что такое введение делает эллиптические кривые более понятными для начинающих, поскольку эллиптическая кривая в аффинной репрезентации — это просто набор пар чисел, и поэтому она более доступна читателям, не знакомым с проективными координатами. Однако аффинная репрезентация имеет тот недостаток, что для описания групповой структуры кривой необходима специальная ``точка на бесконечности'', которая не является точкой на кривой.}

Затем мы вводим групповой закон эллиптической кривой и описываем скалярное умножение на эллиптической кривой, которое является конкретизацией экспоненциального отображения общих циклических групп \ref{scalarmultiplication}. После этого мы рассмотрим проективную репрезентацию эллиптических кривых, которая имеет то преимущество, что для представления точки на бесконечности не требуется специального символа.\footnote{Поскольку эта репрезентация концептуально более проста, именно так эллиптические кривые обычно вводятся на занятиях по математике. Мы считаем, что основной недостаток с точки зрения начинающего состоит в том, что в проективной репрезентации точки являются элементами проективных плоскостей, то есть классами чисел.}

Мы завершаем этот раздел явной эквивалентностью, которая преобразует аффинные репрезентации в проективные и наоборот.

\subsection{Аффинная \concept{короткая форма Вейерштрасса}} Вероятно, наименее абстрактный и наиболее прямой способ ввести эллиптические кривые для нематематиков и начинающих — это так называемая \term{аффинная репрезентация (affine representation)} \term{\concept{кривой в короткой форме Вейерштрасса} (short Weierstrass curve)}. Чтобы понять, что это такое, пусть $\F$ — конечное поле характеристики $q$ с $q>3$, и пусть $a,b\in \F$ — два элемента поля такие, что так называемый \term{дискриминант (discriminant)} $4a^3+ 27b^2$ не равен нулю. Тогда \term{\concept{эллиптическая кривая в короткой форме Вейерштрасса} (short Weierstrass elliptic curve)} $E_{a,b}(\F)$ над $\F$ в аффинной репрезентации — это множество всех пар элементов поля $(x,y)\in \F\times \F$, удовлетворяющих кубическому уравнению \concept{короткой формы Вейерштрасса} $y^2=x^3+a\cdot x+b$, вместе с выделенным символом $\Oinf$, называемым \term{точкой на бесконечности (point at infinity)}:

\begin{equation}
\label{def_short_weierstrass_curve}
E_{a,b}(\F) = \{(x,y)\in \F\times \F\;|\; y^2=x^3+a\cdot x+b \} \bigcup \{\Oinf\}
\end{equation}
Термин ``кривая'' используется здесь потому, что если эллиптическая кривая определена над полем характеристики нуль, например над полем $\Q$ рациональных чисел, множество всех точек $(x,y)\in \Q\times \Q$, удовлетворяющих $y^2 = x^3 +a\cdot x +b$, выглядит как кривая. Однако следует отметить, что визуализация эллиптических кривых над конечными полями как ``кривых'' имеет свои ограничения, и поэтому мы не будем слишком задерживаться на геометрической картине, а сосредоточимся на вычислительных свойствах. Чтобы понять визуальное различие, рассмотрим следующие две эллиптические кривые: 

\medskip

% let Sage draw some elliptic curve in R^2, but show only the picture
\begin{sagesilent}
E1 = EllipticCurve([-2,1])
C1 = E1.plot()
F = GF(9973)
E2 = EllipticCurve(F, [-2,1])
C2 = E2.plot()
\end{sagesilent}
\begin{minipage}{0.48\textwidth}
\sageplot[scale=.48]{C1} 
\end{minipage}
%
\begin{minipage}{0.48\textwidth}
\label{plot:elliptic_curve}
\sageplot[scale=.5]{C2}
\end{minipage}

Обе эллиптические кривые определены одним и тем же уравнением \concept{короткой формы Вейерштрасса} $y^2 = x^3-2x+1$, но первая кривая определена над рациональными числами $\Q$, то есть пара $(x,y)$ содержит рациональные числа, тогда как вторая определена над простым полем $\F_{9973}$, что означает, что обе координаты $x$ и $y$ принадлежат простому полю $\F_{9973}$. Каждая синяя точка представляет пару $(x,y)$, являющуюся решением $y^2 = x^3-2x+1$. Как мы видим, вторая кривая едва ли похожа на геометрическую структуру, которую можно было бы естественно назвать кривой. Это показывает, что наши геометрические интуиции из $\Q$ затуманиваются для кривых над конечными полями.

Уравнение $4a^3+ 27b^2\neq 0$ гарантирует, что кривая \term{несингулярна (non-singular)}, что в общих чертах означает, что кривая не имеет острых точек или самопересечений в геометрическом смысле, если рассматривать её как настоящую кривую. Как мы увидим в \ref{sec:affine_group_law}, острые точки и самопересечения сделали бы групповой закон потенциально неоднозначным.

На протяжении всей книги мы поощряли вас выполнять как можно больше вычислений вручную, поскольку это помогает глубже понять детали. Однако при работе с эллиптическими кривыми вычисления могут быстро стать громоздкими и утомительными, и мы можем потеряться в деталях. К счастью, Sage очень полезен при работе с эллиптическими кривыми. Следующий фрагмент показывает способ определения эллиптических кривых и работы с ними в Sage:

\begin{sagecommandline}
sage: F5 = GF(5) # определить базовое поле
sage: a = F5(2) # параметр a
sage: b = F5(4) # параметр b
sage: # проверить дискриминант
sage: F5(6)*(F5(4)*a^3+F5(27)*b^2) != F5(0)
sage: # кривая в короткой форме Вейерштрасса над полем F5
sage: E = EllipticCurve(F5,[a,b]) # y^2 == x^3 + ax +b 
sage: # точка на кривой
sage: P = E(0,2) # 2^2 == 0^3 + 2*0 + 4
sage: P.xy() # аффинные координаты
sage: INF = E(0) # точка на бесконечности
sage: try: 	# у точки на бесконечности нет аффинных координат
....:     INF.xy()
....: except ZeroDivisionError:
....:     pass
sage: P = E.plot() # создать графическое представление 
\end{sagecommandline}
Следующие три примера дают более практическое понимание того, что такое эллиптическая кривая и как мы можем её вычислять. Мы советуем вам внимательно прочитать эти примеры и, по возможности, выполнить вычисления самостоятельно. Мы неоднократно будем опираться на эти примеры в этой главе и использовать \examplename{} \ref{TJJ13} на протяжении всей книги.
\begin{example}\label{E1F5} Рассмотрим простое поле $\F_5$ из \examplename{} \ref{primfield_z_5}. Чтобы определить эллиптическую кривую над $\F_5$, мы должны выбрать два числа $a$ и $b$ из этого поля. Предположим, мы выбираем $a=1$ и $b=1$, тогда $\kongru{4a^3+ 27b^2}{1}{5}$. Это означает, что соответствующая эллиптическая кривая $E_{1,1}(\F_5)$ задаётся множеством всех пар $(x,y)$ из $\F_5$, удовлетворяющих уравнению $y^2=x^3+x+1$, вместе со специальным символом $\Oinf$, который представляет ``точку на бесконечности''. 

Чтобы лучше понять эту кривую, заметим, что если мы произвольно выберем для проверки пару $(x,y)=(1,1)$, то увидим, что $1^2 \neq 1^3+1 + 1$, и, следовательно, $(1,1)$ не является точкой на кривой $E_{1,1}(\F_5)$. С другой стороны, если мы выберем для проверки пару $(x,y)=(2,1)$, то увидим, что $1^2 = 2^3 + 2 + 1$, и, следовательно, пара $(2,1)$ является точкой на кривой $E_{1,1}(\F_5)$ (Помните, что все вычисления выполняются в арифметике по модулю $5$.)

Поскольку множество $\F_5\times \F_5$ всех пар $(x,y)$ из $\F_5$ содержит всего $5\cdot 5=25$ пар, мы можем вычислить кривую, просто подставив каждую возможную пару $(x,y)$ в уравнение \concept{короткой формы Вейерштрасса} $y^2 = x^3 + x +1$. Если уравнение выполняется, пара является точкой кривой. Если нет, это означает, что точка не лежит на кривой. Объединение результата этого вычисления с точкой на бесконечности даёт кривую следующим образом:
$$
E_{1,1}(\F_5) = \{\Oinf, (0,1),(2,1),(3,1),(4,2),(4,3),(0,4),(2,4),(3,4)\}
$$
Это означает, что эллиптическая кривая — это множество из $9$ элементов, $8$ из которых — пары элементов из $\F_5$, а один — специальный символ $\Oinf$ (точка на бесконечности). Визуализация $E_{1,1}(\F_5)$ даёт следующий график:\sme{should $\Oinf$ also be shown as a point?}
\begin{sagesilent}
F5 = GF(5)
E1 = EllipticCurve(F5, [1,1])
C1 = E1.plot()
\end{sagesilent}
\begin{center} 
\sageplot[scale=.5]{C1}
\end{center}
% sage: AffinePoints = [P.xy() for P in E1.points() if P.order > 1]
\end{example}
При разработке SNARK иногда необходимо выполнять криптографию на эллиптических кривых ``в схеме'', что в основном означает, что эллиптическая кривая должна быть реализована определённым образом, удобным для SNARK. Мы рассмотрим, что это означает, в главе \ref{chap:circuit-compilers}. Чтобы делать это эффективно, желательно иметь кривые со специальными свойствами. Следующий пример — это версия такой кривой для работы вручную, называемая Tiny-\comms{jubjub}. Мы спроектировали эту кривую специально так, чтобы она напоминала хорошо известную криптографически стойкую кривую, называемую \curvename{Baby-jubjub}, последняя из которых широко используется в реальных SNARK (см. спецификацию \href{https://github.com/iden3/iden3-docs/blob/master/source/iden3_repos/research/publications/zkproof-standards-workshop-2/baby-jubjub/baby-jubjub.rst}{Baby-JubJub}).\footnote{В литературе \term{кривая Baby-jubjub} обычно вводится как так называемая \concept{кривая в крученой форме Эдвардса (twisted Edwards curve)}, которую мы рассмотрим в \ref{sec:edwards}. Однако, как мы увидим в \ref{sec:edwards}, каждая \concept{кривая в крученой форме Эдвардса} эквивалентна \concept{кривой в короткой форме Вейерштрасса}, и поэтому мы начинаем с введения Tiny-Jubjub в её реализации \concept{короткой формы Вейерштрасса}.}
Мы советуем заинтересованному читателю внимательно изучить этот пример, поскольку мы будем использовать его и опираться на него в различных местах на протяжении всей книги. 

\begin{example}[Кривая \curvename{Tiny-jubjub}]\label{TJJ13} Рассмотрим простое поле $\F_{13}$ из \exercisename{} \ref{prime_field_F13}. Если мы выберем $a=8$ и $b=8$, то $\kongru{4a^3+ 27b^2}{6}{13}$, и соответствующая эллиптическая кривая задаётся всеми парами $(x,y)$ из $\F_{13}$ такими, что выполняется $y^2=x^3+8x+8$. Мы называем эту кривую \term{кривой \curvename{Tiny-jubjub}} (в аффинной \concept{короткой форме Вейерштрасса}), или сокращённо \TJJ{}.

Поскольку множество $\F_{13}\times \F_{13}$ всех пар $(x,y)$ из $\F_{13}$ содержит всего $13\cdot 13=169$ пар, мы можем вычислить кривую, просто подставив каждую возможную пару $(x,y)$ в уравнение \concept{короткой формы Вейерштрасса} $y^2 = x^3 +8x +8$. Мы получаем следующий результат:
\begin{multline}\label{eq:TJJ13-weierstrass}
%\begin{split}
\mathit{TJJ\_13} = \{\Oinf, (1, 2), (1, 11), (4, 0), (5, 2), (5, 11), (6, 5), (6, 8), (7,2), (7, 11), \\ (8, 5), (8, 8), (9, 4), (9, 9), (10, 3), (10,10), (11, 6), (11, 7), (12, 5), (12, 8)\}
%\end{split}
\end{multline}
Как мы видим, кривая состоит из $20$ точек; $19$ пар элементов из $\F_{13}$ и точки на бесконечности. Чтобы получить визуальное представление кривой \TJJ{}, мы можем изобразить все её точки (кроме точки на бесконечности): 
\begin{sagesilent}
F13 = GF(13)
TJJ_13 = EllipticCurve(F13, [8,8])
CTJJ_13 = TJJ_13.plot()
\end{sagesilent}
\begin{center} 
\sageplot[scale=.5]{CTJJ_13}
\end{center}
Как мы увидим далее, эта кривая довольно особенна, поскольку её можно представить в двух альтернативных формах, называемых \term{формой Монтгомери (Montgomery form)} и \term{\concept{крученой формой Эдвардса} (twisted Edwards form)} (см. разделы \ref{sec:montgomery} и \ref{sec:edwards} соответственно).
\end{example}
Теперь, когда мы увидели две удобные для работы вручную эллиптические кривые, давайте рассмотрим кривую, которая используется в реальной криптографии. Криптографически стойкие эллиптические кривые \hilight{качественно} не отличаются от кривых, которые мы рассмотрели до сих пор, но простое число — модуль их простого поля — значительно больше. Типичные примеры используют простые числа, двоичное представление которых имеет величину более чем вдвое превышающую желаемый \uterm{уровень безопасности}. Например, если требуется безопасность $128$ бит, выбирается простой модуль двоичного размера $\geq 256$. Следующий пример предоставляет такую кривую. 

\begin{example}[Кривая \curvename{secp256k1} Bitcoin]\label{secp256k1}
Чтобы привести пример реальной криптографически стойкой кривой, рассмотрим кривую \curvename{secp256k1}, которая известна тем, что используется в криптографии с открытым ключом \href{https://bitcoin.org/}{Bitcoin}. Простое поле $\F_p$ кривой \curvename{secp256k1} определяется следующим простым числом:
$$
p = \scriptstyle 115792089237316195423570985008687907853269984665640564039457584007908834671663
$$
 
 Двоичное представление этого числа требует $256$ бит, что означает, что простое поле $\F_p$ содержит примерно $2^{256}$ элементов, что считается довольно большим. Чтобы лучше представить, насколько велико базовое поле: число $2^{256}$ примерно того же порядка величины, что и оценочное число атомов в наблюдаемой вселенной. 

Кривая \curvename{secp256k1} определяется параметрами $a,b\in \F_p$ с $a=0$ и $b=7$. Поскольку $\Zmod{4\cdot 0^3 + 27\cdot 7^2}{p}=1323$, эти параметры действительно определяют эллиптическую кривую, заданную следующим образом:
$$
\mathit{secp256k1} = \{(x,y)\in \F_p\times \F_p \;|\; y^2 = x^3 +7\;\} 
$$
Очевидно, кривая \curvename{secp256k1} слишком велика, чтобы быть полезной для вычислений вручную, поскольку можно показать, что число её элементов — простое число $r$, которое также имеет двоичное представление длиной $256$ бит:
$$
r = \scriptstyle 11579208923731619542357098500868790785283756427907490438260516
3141518161494337
$$
Криптографически стойкие эллиптические кривые поэтому не подходят для вычислений вручную, но, к счастью, Sage эффективно обрабатывает большие кривые:
\begin{sagecommandline}
sage: p = 115792089237316195423570985008687907853269984665640564039457584007908834671663
sage: # Шестнадцатеричное представление
sage: p.str(16)
sage: p.is_prime()
sage: p.nbits()
sage: Fp = GF(p)
sage: secp256k1 = EllipticCurve(Fp,[0,7])
sage: r = secp256k1.order() # число элементов
sage: r.str(16)
sage: r.is_prime()
sage: r.nbits()
\end{sagecommandline}
%\seqsplit{115792089237316195423570985008687907853269984665640564039457584007908834671663}
\end{example}
\begin{example}[Кривая \curvename{alt\_bn128} Ethereum]\label{BN128}
Чтобы привести пример реальной криптографически стойкой кривой, которую мы будем использовать в наших реализациях на \lgname{circom}, рассмотрим кривую \curvename{alt\_bn128}, определённую в \href{https://github.com/ethereum/EIPs/blob/master/EIPS/eip-197.md}{EIP-197}. Эта кривая используется для верификации zk-SNARK в блокчейне Ethereum. Простое поле $\F_p$ кривой \curvename{alt\_bn128} определяется следующим простым числом:
$$
p = \scriptstyle 21888242871839275222246405745257275088696311157297823662689037894645226208583
$$
Двоичное представление этого числа требует $254$ бита, что означает, что простое поле $\F_p$ содержит примерно $2^{254}$ элементов. 

\curvename{alt\_bn128} — это кривая в короткой форме Вейерштрасса, определённая параметрами $a,b\in \F_p$ с $a=0$ и $b=3$. Поскольку $\Zmod{4\cdot 0^3 + 27\cdot 3^2}{p}=243$, эти параметры действительно определяют эллиптическую кривую, заданную следующим образом:
$$
\mathit{alt\_bn128} = \{(x,y)\in \F_p\times \F_p \;|\; y^2 = x^3 +3\;\} 
$$
Число точек на эллиптической кривой \curvename{alt\_bn128} — простое число $r$, которое также имеет двоичное представление длиной $254$ бита:
$$
r = \scriptstyle 21888242871839275222246405745257275088548364400416034343698204186575808495617
$$
Мы обозначаем $\F_{bn128}:=F_r$ соответствующее простое поле. Чтобы использовать эту кривую в наших примерах на circom, мы реализуем эту эллиптическую кривую в короткой форме Вейерштрасса в Sage:
\begin{sagecommandline}
sage: p = 21888242871839275222246405745257275088696311157297823662689037894645226208583
sage: p.is_prime()
sage: p.nbits()
sage: Fbn128base = GF(p)
sage: bn128 = EllipticCurve(Fbn128base,[0,3])
sage: r = bn128.order() # число элементов
sage: r.is_prime()
sage: r.nbits()
\end{sagecommandline}
\end{example}

\begin{exercise} Рассмотрите кривую $E_{1,1}(\F_5)$ из \examplename{} \ref{E1F5} и вычислите множество всех точек кривой $(x,y)\in E_{1,1}(\F_5)$.
\end{exercise}
\begin{exercise} Рассмотрите кривую $TJJ\_13$ из \examplename{} \ref{TJJ13} и вычислите множество всех точек кривой $(x,y)\in TJJ\_13$.
\end{exercise}
\begin{exercise}
Найдите определение кривой \curvename{BLS12-381}, реализуйте её в Sage и вычислите число всех точек кривой.
\end{exercise}
\subsubsection{Изоморфные аффинные \concept{кривые в короткой форме Вейерштрасса}}
\label{sec:isomorphic_curves} Как объяснялось ранее в этой главе, эллиптические кривые определяются парами параметров $(a, b)\in \F\times \F$ для некоторого поля $\F$. Важным вопросом при классификации эллиптических кривых является определение того, какие пары параметров $(a,b)$ и $(a',b')$ задают эквивалентные кривые в том смысле, что между множествами точек кривых существует взаимно однозначное соответствие.

To be more precise, let $\F$ be a field, and let $(a,b)$ and $(a',b')$ be two pairs of parameters such that there is an invertible field element $c\in \F^*$ such that $a' = a\cdot c^4$ and $b' = b\cdot c^6$ hold. Then the elliptic curves $E_{a,b}(\F)$ and $E_{a',b'}(\F)$ are \term{isomorphic}, and there is a map that maps the curve points of $E_{a,b}(\F)$ onto the curve points of $E_{a',b'}(\F)$:
\begin{equation}
\label{eq:curve_isomorphism}
I: E_{a,b}(\F) \to E_{a',b'}(\F):\; \begin{cases}(x,y)\\\Oinf\end{cases}
\mapsto \begin{cases}(c^2\cdot x,c^3\cdot y)\\\Oinf\end{cases}
\end{equation}
Это отображение является взаимно однозначным соответствием, и его обратное отображение задаётся переводом точки на бесконечности в точку на бесконечности и каждой точки кривой $(x,y)$ в точку кривой $(c^{-2}x,c^{-4}y)$.
\begin{example}
\label{ex:isomorphic_E1F5}
 Рассмотрим \concept{эллиптическую кривую в короткой форме Вейерштрасса} $E_{1,1}(\F_5)$ из \examplename{} \ref{E1F5} и следующую эллиптическую кривую:
\begin{equation}
E_{1,4}(\F_5):= \{(x,y)\in \F_5\times \F_5 \; | \; y^2 = x^3 + x + 4 \}
\end{equation}

Если мы подставим все пары элементов $(x,y)\in \F_5\times \F_5$ в уравнение \concept{короткой формы Вейерштрасса} $y^2 = x^3 + x + 4$ кривой $E_{1,4}(\F_5)$, мы получим следующее множество точек:
\begin{equation}
E_{1,4}(\F_5) = \{
\Oinf, (0,2), (0,3), (1,1), (1,4), (2,2), (2,3), (3,2), (3,3)\}
\end{equation}

Как мы видим, обе кривые имеют одинаковый порядок. Поскольку $2$ — обратимый элемент из $\F_5$ с $1 = 2^4\cdot 1$ и $4=2^6\cdot 1$, кривые $E_{1,4}(\F_5)$ и $E_{1,1}(\F_5)$ изоморфны: отображение $I: E_{1,1}(\F_5)\to E_{1,4}(\F_5): (x,y)\mapsto (4x,3y)$ из \ref{eq:curve_isomorphism} определяет взаимно однозначное соответствие. Например, точка $(4,3)\in E_{1,1}(\F_5)$ отображается в точку $I(4,3)=(4\cdot 4, 3\cdot 3) = (1,4)\in E_{1,4}(\F)$.
\end{example}
  
\begin{exercise}
Пусть $\F$ — конечное поле, пусть $(a,b)$ и $(a',b')$ — две пары параметров, и пусть $c\in \F^*$ — обратимый элемент поля такой, что выполняются равенства $a' = a\cdot c^4$ и $b' = b\cdot c^6$. Покажите, что функция $I$ из \eqref{eq:curve_isomorphism} отображает точки кривой в точки кривой.
\end{exercise}

\begin{exercise}
\label{ex:isomorphic_TJJ13} Рассмотрите кривую \curvename{Tiny-jubjub} из \examplename{} \ref{TJJ13} и эллиптическую кривую $E_{7,5}(\F_{13})$, определённую следующим образом:
\begin{equation}
E_{7,5}(\F_{13}) = \{ (x,y)\in \F_{13}\times \F_{13}\;|\; y^2 = x^3 + 7x +5\}
\end{equation}
Покажите, что $TJJ\_13$ и $E_{7,5}(\F_{13})$ изоморфны. Затем вычислите множество всех точек $E_{7,5}(\F_{13})$, постройте $I$ и отобразите все точки $TJJ\_13$ в $E_{7,5}(\F_{13})$.
\end{exercise}

\subsubsection{Аффинное сжатое представление}
\label{sec:affine_point_compression}
Как мы видели в \examplename{} \ref{secp256k1}, криптографически стойкие эллиптические кривые определены над большими простыми полями, где элементы этих полей обычно требуют более $255$ бит памяти на компьютере. Поскольку точки эллиптических кривых состоят из пар таких элементов поля, они требуют удвоенного объёма памяти.

Однако мы можем уменьшить объём памяти, необходимый для представления точки кривой, используя технику, называемую \term{сжатием точек (point compression)}. Чтобы понять это, заметим, что для каждого данного $x\in\F$ существует только $2$ возможных значения $y\in \F$ таких, что пара $(x,y)$ является точкой на аффинной \concept{короткой форме Вейерштрасса}, поскольку $x$ и $y$ должны удовлетворять уравнению $y^2 = x^3 + a\cdot x + b$. Отсюда следует, что $y$ можно вычислить из $x$, поскольку оно является элементом из множества квадратных корней $\sqrt{x^3 + a\cdot x +b}$ (см. \ref{ded:square_root}), которое содержит ровно два элемента для $x^3 + a\cdot x +b\neq 0$ и ровно один элемент для $x^3 + a\cdot x +b=0$. 

Это означает, что мы можем представить точку кривой в \term{сжатой форме (compressed form)}, просто сохранив координату $x$ вместе с одним битом, называемым \term{битом знака (sign bit)}, который детерминистически определяет, какой из двух корней выбрать. Одно из соглашений может состоять в том, чтобы всегда выбирать корень, ближе к $0$, когда бит знака равен $0$, и корень, ближе к порядку $\F$, когда бит знака равен $1$. В случае, если координата $y$ равна нулю, оба бита знака дают одинаковый результат.

\begin{example}[\curvename{Tiny-jubjub}] Чтобы немного лучше понять концепцию сжатых точек кривой, снова рассмотрим кривую \TJJ{} из \examplename{} \ref{TJJ13}. Поскольку эта кривая определена над простым полем $\F_{13}$, а числа между $0$ и $13$ требуют примерно $4$ бит для представления, каждая точка \TJJ{} на этой кривой требует $8$ бит памяти в несжатой форме. Следующее множество представляет несжатую форму точек на этой кривой:
\begin{multline}
\mathit{TJJ\_13} = \{\Oinf, (1, 2), (1, 11), (4, 0), (5, 2), (5, 11), (6, 5), (6, 8), (7,2), (7, 11), \\ (8, 5), (8, 8), (9, 4), (9, 9), (10, 3), (10,
10), (11, 6), (11, 7), (12, 5), (12, 8)\}
\end{multline}
Используя технику сжатия точек, мы можем уменьшить биты, необходимые для представления точек на этой кривой, до $5$ на точку. Чтобы достичь этого, мы можем заменить координату $y$ в каждой паре $(x,y)$ битом знака, указывающим, ближе ли $y$ к $0$ или к $13$. В результате значения $y$ в диапазоне $[0,\ldots,6]$ будут иметь бит знака $0$, тогда как значения $y$ в диапазоне $[7,\ldots,12]$ будут иметь бит знака $1$. Применение этого к точкам \TJJ{} даёт сжатое представление следующим образом:
\begin{multline}
\mathit{TJJ\_13} = \{\Oinf, (1, 0), (1, 1), (4, 0), (5, 0), (5, 1), (6, 0), (6, 1), (7,0), (7, 1), \\ (8, 0), (8, 1), (9, 0), (9, 1), (10, 0), (10,1), (11, 0), (11, 1), (12, 0), (12, 1)\}
\end{multline} 
Заметьте, что числа $7,\ldots, 12$ являются отрицательными (аддитивно обратными) числам $1,\ldots, 6$ в арифметике по модулю $13$, и что $-0=0$.

Чтобы восстановить несжатый аналог, скажем, сжатой точки $(5,1)$, мы подставляем координату $x=5$ в уравнение \concept{короткой формы Вейерштрасса} и получаем $y^2 = 5^3 + 8\cdot 5 +8 = 4$. Как и ожидалось, $4$ является квадратичным вычетом в $\F_{13}$ с корнями $\sqrt{4}= \{2,11\}$. Поскольку бит знака точки равен $1$, мы должны выбрать корень, ближе к модулю $13$, то есть $11$. Несжатая точка поэтому равна $(5,11)$. 
\end{example}
\begin{comment}
Если посмотреть на предыдущие примеры, коэффициент сжатия не выглядит очень впечатляющим. Однако рассмотрение реального примера кривой \curvename{secp256k1} показывает, что сжатие имеет значительные практические преимущества.
\begin{example}
Снова рассмотрим кривую \curvename{secp256k1} из \examplename{} \ref{secp256k1}. Следующий код использует Sage для генерации случайной аффинной точки кривой, а затем применяет к ней наш метод сжатия:
\begin{sagecommandline}
sage: P = secp256k1.random_point().xy()
sage: P
sage: # размер несжатой аффинной точки
sage: ZZ(P[0]).nbits()+ZZ(P[1]).nbits()
sage: # вычислить сжатие
sage: if P[1] > Fp(-1)/Fp(2):
....:     PARITY = 1
....: else:
....:     PARITY = 0
sage: PCOMPRESSED = [P[0],PARITY]
sage: PCOMPRESSED
sage: # размер сжатой аффинной точки
sage: ZZ(PCOMPRESSED[0]).nbits()+ZZ(PCOMPRESSED[1]).nbits()
\end{sagecommandline}
\end{example}\sme{add explanation of how this shows what we claim}
\end{comment}

\subsection{\concept{Аффинный групповой закон}}
\label{sec:affine_group_law}
%group law
% http://wwwmayr.informatik.tu-muenchen.de/konferenzen/Jass07/courses/1/Lukyanenko/Lukyanenk_Paper.pdf
Одно из ключевых свойств эллиптической кривой заключается в том, что на множестве её точек можно определить групповой закон так, что точка на бесконечности служит нейтральным элементом, а обратные элементы являются отражениями относительно оси $x$. Происхождение этого закона можно понять на геометрической картине, и он известен как \term{правило хорд и касательных (chord-and-tangent rule)}. В аффинной репрезентации \concept{короткой формы Вейерштрасса} правило можно описать следующим образом, используя символ $\oplus$ для группового закона:

\begin{definition}[\deftitle{Правило хорд и касательных: геометрическое определение}]\label{def:chord-and-tangent}

\defvsep{}

\begin{itemize}
\item (Точка на бесконечности) Мы определяем точку на бесконечности $\Oinf$ как нейтральный элемент сложения, то есть определяем $P\oplus\Oinf = P$ для всех точек $P\in E(\F)$.
\item (Правило хорды) Пусть $P$ и $Q$ — две различные точки на эллиптической кривой, ни одна из которых не является точкой на бесконечности: $P, Q\in E(\F)\backslash \{\Oinf\}$ и $P\neq Q$\\
Сумма $P$ и $Q$ определяется следующим образом:\\
Рассмотрим прямую $l$, которая пересекает кривую в точках $P$ и $Q$. Если $l$ пересекает эллиптическую кривую в третьей точке $R'$, определим сумму $P$ и $Q$ как отражение $R'$ относительно оси $x$: $R=P\oplus Q$. Если прямая $l$ не пересекает кривую в третьей точке, определим сумму как точку на бесконечности $\Oinf$. Такая прямая называется \term{хордой (chord)}; можно показать, что никакая хорда не пересечёт кривую более чем в трёх точках. Это означает, что сложение не является неоднозначным.
\item (Правило касательной) Пусть $P$ — точка на эллиптической кривой, не являющаяся точкой на бесконечности: $P \in E(\F)\backslash \{\Oinf\}$\\
 Сумма $P$ с самой собой (удвоение $P$) определяется следующим образом:\\
Рассмотрим прямую, касательную к эллиптической кривой в точке $P$, в том смысле, что она ``просто касается'' кривой в этой точке. Если эта прямая пересекает эллиптическую кривую во второй точке $R'$, сумма $P\oplus P$ является отражением $R'$ относительно оси $x$. Если она не пересекает кривую в третьей точке, определим сумму как точку на бесконечности $\Oinf$. Такая прямая называется \term{касательной (tangent)}; можно показать, что никакая такая касательная не пересечёт кривую более чем в двух точках. Это означает, что удвоение не является неоднозначным.
\end{itemize}
\end{definition}

Можно показать, что точки эллиптической кривой образуют коммутативную группу относительно вышеуказанного правила хорд и касательных так, что $\Oinf$ выступает нейтральным элементом, а обратным элементом для любого элемента $P\in E(\F)$ является отражение $P$ относительно оси $x$. 

Правило хорд и касательных геометрически определяет групповой закон эллиптической кривой, и мы только что неформально сформулировали его как интуицию выше. Чтобы применить эти правила на компьютере, мы должны перевести их в алгебраические уравнения. Для этого сначала заметим, что для любых двух данных точек кривой $(x_1,y_1), (x_2,y_2)\in E(\F)$ тождество $x_1=x_2$ влечёт $y_2=\pm y_1$, как объяснялось в \secname{} \ref{sec:affine_point_compression}. Это показывает, что следующие правила являются полным описанием группы эллиптической кривой $(E(\F),\oplus)$:

\begin{definition}[\deftitle{Правило хорд и касательных: алгебраическое определение}]\label{def:chord-tangent-algebra}
\defvsep{}
\begin{itemize}
\item (Нейтральный элемент) Точка на бесконечности $\Oinf$ является нейтральным элементом.
\item (Обратный элемент) Обратным для $\mathcal{O}$ является $\mathcal{O}$. Для любой другой точки кривой $(x,y) \in E(\F)\backslash \{\mathcal{O}\}$ обратный элемент задаётся $(x,-y)$.
\item (Групповой закон) Для любых двух точек кривой $P, Q \in E(\F)$ групповой закон определяется одним из следующих случаев:
\begin{enumerate}
\item (Нейтральный элемент) Если $Q=\Oinf$, то групповой закон определяется как $P\oplus Q=P$.
\item (Обратные элементы) Если $P=(x,y)$ и $Q=(x,-y)$, групповой закон определяется как $P\oplus Q=\Oinf$.
\item (Правило касательной) Если $P=(x,y)$ с $y\neq 0$, групповой закон $P\oplus P=(x',y')$ определяется следующим образом:
$$
\begin{array}{llr}
x' = \left(\frac{3x^2+a}{2y}\right)^2 -2x &,&
y' = \left(\frac{3x^2+a}{2y}\right)\left(x-x'\right) - y
\end{array} 
$$
\item (Правило хорды) Если $P=(x_1,y_1)$ и $Q=(x_2,y_2)$ такие, что $x_1 \neq x_2$, групповой закон $R=P\oplus Q$ с $R=(x_3,y_3)$ определяется следующим образом:
$$
\begin{array}{llr}
x_3 = \left(\frac{y_2-y_1}{x_2-x_1}\right)^2 -x_1-x_2 &, &
y_3 = \left(\frac{y_2-y_1}{x_2-x_1} \right)\left(x_1-x_3\right) - y_1
\end{array} 
$$
\end{enumerate}
\end{itemize}
\end{definition}
\begin{notation}
Пусть $\F$ — поле и $E(\F)$ — эллиптическая кривая над $\F$. Мы пишем $\oplus$ для группового закона на $E(\F)$, $(E(\F),\oplus)$ для коммутативной группы точек эллиптической кривой, и используем аддитивную нотацию (\notationname{} \ref{def:additive_notation}) в этой группе. Если $P$ — точка на \concept{короткой форме Вейерштрасса} с $P=(x,0)$, то $P$ называется \term{самообратной (self-inverse)}.
\end{notation}
Как мы видим, вычисление обратных элементов на эллиптических кривых очень эффективно. Однако вычисление сложения точек эллиптической кривой в аффинной репрезентации требует рассмотрения многих случаев и включает обширные деления в конечном поле. Как мы увидим в \ref{sec:projective_group_law}, закон сложения упрощается в проективных координатах.

Давайте рассмотрим несколько практических примеров того, как вычисляется групповой закон на эллиптической кривой.

\begin{example}\label{ex:01+42}
Снова рассмотрим эллиптическую кривую $E_{1,1}(\F_5)$ из \examplename{} \ref{E1F5}. Как мы видели, кривая состоит из следующих $9$ элементов:
\begin{equation}
E_{1,1}(\F_5) = \{\Oinf, (0,1),(2,1),(3,1),(4,2),(4,3),(0,4),(2,4),(3,4)\}
\end{equation}
Мы знаем, что это множество определяет группу, поэтому мы можем выполнять сложение любых двух элементов из $E_{1,1}(\F_5)$, чтобы получить третий элемент этой группы. 

Для примера рассмотрим элементы $(0,1)$ и $(4,2)$. Ни один из этих элементов не является нейтральным элементом $\Oinf$, и поскольку координата $x$ точки $(0,1)$ отличается от координаты $x$ точки $(4,2)$, мы знаем, что должны использовать правило хорды из определения \ref{def:chord-tangent-algebra} для вычисления суммы $(0,1)\oplus (4,2)$:
%\begin{tabular}{lr}
\begin{align*}
x_3  & = \left(\frac{y_2-y_1}{x_2-x_1}\right)^2 -x_1-x_2 & \text{\# insert points}\\
     & = \left(\frac{2-1}{4-0}\right)^2 -0-4  & \text{\# simplify in } \F_5\\
     & = \left(\frac{1}{4}\right)^2 +1
       = 4^2 +1
       = 1 +1 
       = 2
\\
\\
y_3  & = \left(\frac{y_2-y_1}{x_2-x_1} \right)\left(x_1-x_3\right) - y_1  & \text{\# insert points}\\     
     & = \left(\frac{2-1}{4-1} \right)\left(0-2\right) - 1   & \text{\# simplify in } \F_5\\    
     & = \left(\frac{1}{4} \right)\cdot 3 + 4   
       = 4\cdot 3 + 4
       = 2 + 4
       = 1          
\end{align*} 
%\end{tabular}
Итак, на эллиптической кривой $E_{1,1}(\F_5)$ мы получаем $(0,1)\oplus (4,2) =(2,1)$, и действительно, пара $(2,1)$ является элементом $E_{1,1}(\F_5)$, как и ожидалось. С другой стороны, $(0,1)\oplus (0,4) =\Oinf$, поскольку обе точки имеют равные координаты $x$ и обратные координаты $y$, делая их обратными друг другу. Складывая точку $(4,2)$ с самой собой, мы должны использовать правило касательной из определения \ref{def:chord-tangent-algebra}:
\begin{align*}
x'  & = \left(\frac{3x^2+a}{2y}\right)^2 -2x   & \text{\# insert points}\\
    & = \left(\frac{3\cdot 4^2+1}{2\cdot 2}\right)^2 -2\cdot 4 & \text{\# simplify in } \F_5 \\
    & = \left(\frac{3\cdot 1+1}{4}\right)^2 +3\cdot 4
      = \left(\frac{4}{4}\right)^2 +2
      = 1 +2 
      = 3
\\
\\
y'  & = \left(\frac{3x^2+a}{2y}\right)^2\left(x-x'\right) - y  & \text{\# insert points} \\
    & = \left(\frac{3\cdot 4^2+1}{2\cdot 2}\right)^2\left(4-3\right) - 2 & \text{\# simplify in } \F_5\\
    & = 1\cdot 1 + 3
      = 4
\end{align*}
Итак, на эллиптической кривой $E_{1,1}(\F_5)$ мы получаем удвоение точки $(4,2)$, то есть $(4,2)\oplus (4,2) =(3,4)$, и действительно, пара $(3,4)$ является элементом $E_{1,1}(\F_5)$, как и ожидалось. Группа $E_{1,1}(\F_5)$ не имеет самообратных точек, кроме нейтрального элемента $\Oinf$, поскольку ни одна точка не имеет $0$ в качестве координаты $y$. Мы можем использовать Sage для перепроверки вычислений. 
\begin{sagecommandline}
sage: F5 = GF(5)
sage: E1 = EllipticCurve(F5,[1,1])
sage: INF = E1(0) # точка на бесконечности
sage: P1 = E1(0,1)
sage: P2 = E1(4,2)
sage: P3 = E1(0,4)
sage: R1 = E1(2,1)
sage: R2 = E1(3,4)
sage: R1 == P1+P2
sage: INF == P1+P3
sage: R2 == P2+P2
sage: R2 == 2*P2
sage: P3 == P3 + INF
\end{sagecommandline}
\end{example}
\begin{example}[\curvename{Tiny-jubjub}]\label{ex:TJJ13-self-inverse} Снова рассмотрим кривую \TJJ{} из \examplename{} \ref{TJJ13}, и вспомним, что её группа точек задаётся следующим образом:
\begin{multline}
\mathit{TJJ\_13} = \{\Oinf, (1, 2), (1, 11), (4, 0), (5, 2), (5, 11), (6, 5), (6, 8), (7,2), (7, 11), \\ (8, 5), (8, 8), (9, 4), (9, 9), (10, 3), (10,
10), (11, 6), (11, 7), (12, 5), (12, 8)\}
\end{multline}
В отличие от группы из предыдущего примера, эта группа содержит самообратную точку, отличную от нейтрального элемента, определяемую $(4,0)$. Чтобы понять, что это означает, заметим, что мы не можем сложить $(4,0)$ с самой собой, используя правило касательной из \defname{} \ref{def:chord-tangent-algebra}, поскольку координата $y$ равна нулю. Вместо этого мы должны использовать правило для аддитивных обратных, поскольку $0=-0$. Мы получаем $(4,0)\oplus (4,0)=\Oinf$ в \TJJ{}, что показывает, что точка $(4,0)$ является обратной самой себе, потому что сложение её с самой собой даёт нейтральный элемент. Мы проверим наше вычисление с помощью Sage:

\begin{sagecommandline}
sage: F13 = GF(13)
sage: TJJ = EllipticCurve(F13,[8,8])
sage: P = TJJ(4,0)
sage: INF = TJJ(0) # Точка на бесконечности
sage: INF == P+P
sage: INF == 2*P
\end{sagecommandline}
\end{example}
\begin{example}
Снова рассмотрим кривую \curvename{secp256k1} из \examplename{} \ref{secp256k1}. Следующий код использует Sage для генерации двух случайных аффинных точек кривой, а затем складывает эти точки вместе: 
\begin{sagecommandline}
sage: P = secp256k1.random_point()
sage: Q = secp256k1.random_point()
sage: R = P + Q
sage: P.xy()
sage: Q.xy()
sage: R.xy()
\end{sagecommandline}
\end{example}
\begin{exercise}
Рассмотрите коммутативную группу $(\mathit{TJJ\_13},\oplus)$ кривой \curvename{Tiny-jubjub} из \examplename{} \ref{TJJ13}. 
\begin{enumerate}
\item Вычислите обратные элементы для $(10,10)$, $\Oinf$, $(4,0)$ и $(1,2)$.
\item Решите уравнение $x \oplus (9,4) = (5,2) $ для некоторого $x\in \mathit{TJJ\_13}$.
\end{enumerate}
\end{exercise}

\subsubsection{Скалярное умножение}
Как мы видели в предыдущем разделе, эллиптические кривые $E(\F)$ имеют структуру коммутативной группы. Можно показать, что эта группа конечна всякий раз, когда базовое поле $\F$ конечно, но она не обязательно циклическая. Как мы знаем из \eqref{scalarmultiplication}, это означает, что существует понятие скалярного умножения, ассоциированное с любой циклической подгруппой эллиптической кривой над конечными полями.

Чтобы понять это скалярное умножение, вспомним из \defname{} \ref{cyclic-groups}, что каждая конечная циклическая группа порядка $n$ имеет образующий $g$ и ассоциированное экспоненциальное отображение $g^{(\cdot)}: \Z_n \to \G$, где $g^x$ — $x$-кратное произведение $g$ с самим собой. Для любой точки $P$ конечного порядка $n$ на эллиптической кривой подгруппа $\langle P\rangle$ циклическая, и отображение $[\,\cdot\,]P:\Z_n\to \langle P\rangle$ определяется аналогично.

Скалярное умножение на эллиптической кривой — это экспоненциальное отображение, записанное в аддитивной нотации. Более точно, пусть $\F$ — конечное поле, $E(\F)$ — эллиптическая кривая, и $P\in E(\F)$ — точка порядка $n$. Тогда \term{скалярное умножение на эллиптической кривой (elliptic curve scalar multiplication)} с основанием $P$ определяется следующим образом (где $[0]P = \Oinf$ и $[m]P = P+P+\ldots + P$ — $m$-кратная сумма $P$ с самим собой):
$$
[\cdot]P: \Z_n \to \langle P\rangle\;;\; m \mapsto [m]P
$$
Таким образом, скалярное умножение на эллиптической кривой является конкретизацией общего экспоненциального отображения с использованием аддитивной нотации вместо мультипликативной.

\subsubsection{Логарифмический порядок}
\label{def:logarithmic_ordering}
Как объяснялось в \eqref{logarithm_map}, обратное экспоненциальное отображение существует, и оно обычно называется \term{отображением дискретного логарифма на эллиптической кривой (elliptic curve discrete logarithm map)}. Однако мы не знаем никакого эффективного способа фактически вычислить это отображение, и это одна из причин, почему некоторые эллиптические кривые считаются DL-стойкими (см. \defname{} \ref{def:DL-secure}).

Одно полезное свойство экспоненциального отображения в отношении примеров в этой книге заключается в том, что его можно использовать для значительного упрощения вычислений вручную. Как мы видели в \examplename{} \ref{ex:01+42}, вычисление закона сложения на эллиптической кривой требует немало усилий при выполнении без компьютера. Однако когда $g$ — образующий малой группы эллиптической кривой порядка $n$, удобной для работы вручную, мы можем использовать экспоненциальное отображение для записи элементов группы следующим образом, который мы называем их \term{логарифмическим порядком (logarithmic order)} относительно образующего $g$:
\begin{equation}\label{def:logarithmic_order}
\G = \{[1]g\to [2]g \to [3]g\to\cdots\to [n-1]g\to \Oinf\}
\end{equation} 
Для малых групп, удобных для работы вручную, логарифмический порядок значительно упрощает сложное сложение на эллиптической кривой до гораздо более простого случая арифметики по модулю $n$. Чтобы сложить две точки кривой $P$ и $Q$, нам нужно только найти их дискретно-логарифмические соотношения с образующим $P=[l]g$ и $Q=[m]g$, и вычислить групповой закон как $P\oplus Q = [l+m]g$, где $l+m$ — сложение в арифметике по модулю $n$. 

Однако читатель должен помнить, что многие эллиптические кривые считаются DL-стойкими (\defname{} \ref{def:DL-secure}), что означает, что для этих кривых логарифмический порядок не может быть вычислен эффективно.  

В следующем примере мы рассмотрим некоторые следствия из того факта, что эта конкретная группа эллиптической кривой конечна и циклическая, и применим логарифмический порядок.
\begin{example}\label{ex:G1G2-subgroups} Рассмотрим группу эллиптической кривой $E_{1,1}(\F_5)$ из \examplename{} \ref{E1F5}. Поскольку это конечная циклическая группа порядка $9$, а разложение на простые множители $9$ равно $3\cdot 3$, мы можем использовать основную теорему о конечных циклических группах (\defname{} \ref{def:fundamental_theorem_groups}) для рассуждений обо всех её подгруппах. Действительно, поскольку единственные делители $9$ — это $1$, $3$ и $9$, мы знаем, что $E_{1,1}(\F_5)$ имеет следующие подгруппы:
\begin{itemize}
\item $E_{1,1}(\F_5)[9]$ — подгруппа порядка $9$. По определению, любая группа является подгруппой самой себя.
\item $E_{1,1}(\F_5)[3] = \{(2,1),(2,4),\Oinf\}$ — подгруппа порядка $3$. Это подгруппа, ассоциированная с простым множителем $3$.
\item $E_{1,1}(\F_5)[1] = \{\Oinf\}$ — подгруппа порядка $1$. Это тривиальная подгруппа.
\end{itemize}
Более того, поскольку $E_{1,1}(\F_5)$ и все её подгруппы циклические, мы знаем из \defname{} \ref{cyclic-groups}, что они должны иметь образующие. Например, точка кривой $(2,1)$ является образующим подгруппы порядка $3$ $E_{1,1}(\F_5)[3]$, поскольку каждый элемент $E_{1,1}(\F_5)[3]$ может быть получен повторным сложением $(2,1)$ с самой собой: 
\begin{align*}
[1](2,1) & = (2,1) \\
[2](2,1) & = (2,4) \\
[3](2,1) & = \Oinf
\end{align*}
Поскольку $(2,1)$ — образующий, мы знаем из \eqref{exponentialmap}, что он порождает экспоненциальное отображение из конечного поля $\F_3$ на $\G_2$, определяемое скалярным умножением:
\begin{equation}
[\cdot](2,1): \F_3 \to E_{1,1}(\F_5)[3]\; : \; x\mapsto [x](2,1) 
\end{equation}
Чтобы привести пример образующего, который порождает всю группу $E_{1,1}(\F_5)$, рассмотрим точку $(0,1)$. Применяя правило касательной многократно, мы вычисляем следующим образом:
\begin{equation}
\begin{array}{lccr}
{}\begin{array}{lcl}
{}[0](0,1) &=& \Oinf \\
{}[2](0,1) &=& (4, 2) \\ 
{}[4](0,1) &=& (3, 4) \\ 
{}[6](0,1) &=& (2, 4) \\ 
{}[8](0,1) &=& (0, 4) \\ 
\end{array} & & &
{}\begin{array}{lcl}
{}[1](0,1) &=& (0, 1) \\
{}[3](0,1) &=& (2, 1) \\
{}[5](0,1) &=& (3, 1) \\
{}[7](0,1) &=& (4, 3) \\
{}[9](0,1) &=& \Oinf
\end{array}
\end{array}
\end{equation}
Снова, поскольку $(0,1)$ — образующий, мы знаем из \eqref{exponentialmap}, что он порождает экспоненциальное отображение. Однако, поскольку порядок группы не является простым числом, экспоненциальное отображение отображает не из поля, а из кольца $\Z_9$ арифметики по модулю $9$:
\begin{equation}
[\cdot](0,1): \Z_9 \to E_{1,1}(\F_5)\; : \; x\mapsto [x](0,1) 
\end{equation}
Используя образующий $(0,1)$ и ассоциированное с ним экспоненциальное отображение, мы можем записать $E(\F_1)$ в логарифмическом порядке относительно $(0,1)$, как объяснялось в \defname{} \ref{def:logarithmic_ordering}. Мы получаем следующее:
\begin{equation}
E_{1,1}(\F_5) = \{(0, 1)\to (4, 2)\to (2, 1)\to (3, 4)\to (3, 1)\to (2, 4)\to (4, 3)\to (0, 4)\to \Oinf \}
\end{equation}
Это указывает на то, что первый элемент — образующий, а $n$-й элемент — скалярное произведение $n$ и образующего. Чтобы увидеть, как логарифмические порядки подобного рода упрощают вычисления в малых группах эллиптических кривых, снова рассмотрим \examplename{} \ref{ex:01+42}. В том примере мы использовали правило хорд и касательных для вычисления $(0,1)\oplus (4,2)$. Теперь, в логарифмическом порядке $E_{1,1}(\F_5)$, мы можем вычислить эту сумму гораздо проще, поскольку мы можем непосредственно видеть, что $(0,1)=[1](0,1)$ и $(4,2)=[2](0,1)$. Тогда мы можем вывести $(0,1)\oplus (4,2)= (2,1)$ следующим образом:

\begin{equation}
(0,1)\oplus (4,2)=[1](0,1)\oplus [2](0,1)= [3](0,1)=(2,1)
\end{equation}
Чтобы привести другой пример, мы можем немедленно увидеть, что $(3,4)\oplus (4,3) = (4,2)$, не выполняя никакого сложного сложения на эллиптической кривой, поскольку мы знаем из логарифмического упорядочивания $E_{1,1}(\F_5)$, что $(3,4)= [4](0,1)$ и $(4,3)= [7](0,1)$. Поскольку $4+7 = 2$ в $\Z_9$, результат должен быть $[2](0,1)=(4,2)$.

Наконец, мы можем использовать $E_{1,1}(\F_5)$ как пример для понимания концепции очистки кофактора из \defname{} {\ref{def:cofactor_clearing}}. Поскольку порядок $E_{1,1}(\F_5)$ равен $9$, у нас есть только один множитель, который также оказывается кофактором. Очистка кофактора тогда означает, что мы можем отобразить любой элемент из $E_{1,1}(\F_5)$ в его простую фактор-группу $E_{1,1}(\F_5)[3]$ с помощью скалярного умножения на $3$. Например, взяв элемент $(3,4)$, который не принадлежит $E_{1,1}(\F_5)[3]$, и умножив его на $3$, мы получаем $[3](3,4)= (2,1)$, который, как и ожидалось, является элементом $E_{1,1}(\F_5)[3]$.
\end{example}



\begin{example}\label{ex:TJJ13-cofactor-clearing} Снова рассмотрим кривую \curvename{Tiny-jubjub} \TJJ{} из \examplename{} \ref{TJJ13}. В этом примере мы рассмотрим подгруппы кривой \curvename{Tiny-jubjub}, определим образующие и вычислим логарифмический порядок для вычислений вручную. Затем мы ещё раз взглянем на принцип очистки кофактора.

Поскольку порядок \TJJ{} равен $20$, а разложение на простые множители $20$ равно $2^2\cdot 5$, мы знаем, что \TJJ{} содержит ``большую'' подгруппу простого порядка размера $5$ и малую подгруппу простого порядка размера $2$. 

Чтобы вычислить эти группы, мы можем применить технику очистки кофактора \eqref{def:cofactor_clearing}\sme{check reference} в цикле проб и повторений. Мы начинаем цикл, произвольно выбирая элемент $P\in \mathit{TJJ\_13}$, затем умножая этот элемент на кофактор группы, которую мы хотим вычислить. Если результат равен $\Oinf$, мы пробуем другой элемент и повторяем процесс до тех пор, пока результат не станет отличным от точки на бесконечности $\Oinf$.  

Чтобы вычислить образующий для малой подгруппы простого порядка $(\mathit{TJJ\_13})[2]$, сначала заметим, что кофактор равен $10$, поскольку $20=2\cdot 10$. Затем мы произвольно выбираем точку кривой $(5,11)\in \mathit{TJJ\_13}$ и вычисляем $[10](5,11)=\Oinf$. Поскольку результатом является точка на бесконечности, мы должны попробовать другую точку кривой, скажем $(9,4)$. Мы получаем $[10](9,4)=(4,0)$ и можем сделать вывод, что $(4,0)$ является образующим $(\mathit{TJJ\_13})[2]$. Логарифмический порядок тогда даёт следующий порядок:
\begin{equation}
(\mathit{TJJ\_13})[2] = \{(4,0)\to \Oinf\}
\end{equation}
Это ожидаемо, поскольку мы знаем из \examplename{} \ref{ex:TJJ13-self-inverse}, что $(4,0)$ самообратна, причём $(4,0)\oplus (4,0)=\Oinf$. Мы перепроверяем вычисления с помощью Sage: 
\begin{sagecommandline}
sage: F13 = GF(13)
sage: TJJ = EllipticCurve(F13,[8,8])
sage: P = TJJ(5,11)
sage: INF = TJJ(0)
sage: 10*P == INF
sage: Q = TJJ(9,4)
sage: R = TJJ(4,0)
sage: 10*Q == R
\end{sagecommandline}
Мы можем применить тот же подход к ``большой'' подгруппе простого порядка $(\mathit{TJJ\_13})[5]$, которая содержит $5$ элементов. Чтобы вычислить образующий для этой группы, сначала заметим, что ассоциированный кофактор равен $4$, поскольку $20=5\cdot 4$. Снова выберем точку кривой $(9,4)\in \mathit{TJJ\_13}$ и вычислим $[4](9,4)=(7,11)$. Поскольку результат не является точкой на бесконечности, мы знаем, что $(7,11)$ — образующий $(\mathit{TJJ\_13})[5]$. Используя образующий $(7,11)$, мы вычисляем экспоненциальное отображение $[\cdot](7,11): \F_5 \to \mathit{TJJ\_13}[5]$ и получаем следующее:
\begin{align*}
[0](7,11) &= \Oinf\\
[1](7,11) &= (7,11)\\
[2](7,11) &= (8,5)\\
[3](7,11) &= (8,8)\\
[4](7,11) &= (7,2)
\end{align*}
Мы можем использовать это вычисление для записи большой простой группы $(\mathit{TJJ\_13})[5]$ кривой \curvename{Tiny-jubjub} в логарифмическом порядке, который мы будем довольно часто использовать в дальнейшем. Мы получаем следующее:
\begin{equation}\label{eq:TJJ13-logarithmic-order}
(\mathit{TJJ\_13})[5] = \{(7,11)\to(8,5)\to(8,8)\to(7,2)\to \Oinf\}
\end{equation}
Отсюда мы можем немедленно увидеть, например, что $(8,8)\oplus (7,2)= (8,5)$, поскольку $3+4=2$ в $\F_5$.
\end{example}
На основе двух предыдущих примеров у вас может сложиться впечатление, что вычисления на эллиптических кривых в значительной степени могут быть заменены модульной арифметикой. Однако это в общем случае не так, а является лишь особенностью малых групп, где возможно записать всю группу в логарифмическом порядке.
\begin{exercise} Рассмотрите \examplename{} \ref{ex:G1G2-subgroups} и вычислите множество 
$\{[1](0,1), [2](0,1),\ldots,[8](0,1,[9](0,1)\}$, используя только правило касательной.
\end{exercise}
\begin{exercise} Рассмотрите \examplename{} \ref{ex:TJJ13-cofactor-clearing} и вычислите скалярные умножения $[10](5,11)$, а также $[10](9,4)$ и $[4](9,4)$ вручную, используя алгоритм из упражнения \ref{alg_double-and-add}.
\end{exercise}
\subsection{Проективная \concept{короткая форма Вейерштрасса}}
% https://www.cosic.esat.kuleuven.be/bcrypt/lecture%20slides/wouter.pdf
Как мы видели в предыдущем разделе, описание эллиптических кривых как пар точек, удовлетворяющих определённому уравнению, относительно просто. Однако чтобы определить групповую структуру на множестве точек, мы должны были добавить специальную точку на бесконечности, выступающую нейтральным элементом. 

Вспомнив определение проективных плоскостей \ref{sec:planes},\sme{S: move that section here?} мы знаем, что точки на бесконечности в проективной геометрии рассматриваются как обычные точки. Поэтому имеет смысл рассмотреть определение \concept{кривой в короткой форме Вейерштрасса} в проективной геометрии.

Чтобы понять, что такое \concept{кривая в короткой форме Вейерштрасса} в проективных координатах, пусть $\F$ — конечное поле порядка $q$ и характеристики $>3$, пусть $a,b\in \F$ — два элемента поля такие, что $\Zmod{4a^3+ 27b^2}{q}\neq 0$, и пусть $\F\mathbb{P}^2$ — проективная плоскость над $\F$, как введено в \secname{} \ref{sec:planes}. Тогда \term{проективная \concept{эллиптическая кривая в короткой форме Вейерштрасса} (projective short Weierstrass elliptic curve)} над $\F$ — это множество всех точек $[X:Y:Z]\in \F\mathbb{P}^2$ из проективной плоскости, удовлетворяющих кубическому уравнению $Y^2\cdot Z = X^3+a\cdot X\cdot Z^2 + b\cdot Z^3$:

\begin{equation}
\label{def:projective_cubic_equation}
E(\F\mathbb{P}^2) = \{[X:Y:Z]\in \F\mathbb{P}^2\;|\; Y^2\cdot Z = X^3+a\cdot X\cdot Z^2 + b\cdot Z^3 \}
\end{equation}

Чтобы понять, как точка на бесконечности объединяется в этом определении, вспомним из \secname{} \ref{sec:planes}, что в проективной геометрии точки на бесконечности задаются проективными координатами $[X:Y:0]$. Подстановка представителей $(x_1,y_1,0)\in [X:Y:0]$ из этих координат в определяющее кубическое уравнение \ref{def:projective_cubic_equation} даёт следующее тождество:
\begin{align*}
y_1^2\cdot 0 & = x_1^3+a\cdot x_1\cdot 0^2 + b\cdot 0^3 & \Leftrightarrow \\
0 & = x_1^3
\end{align*} 

Это влечёт $X=0$ и показывает, что единственной проективной точкой на бесконечности, которая также является точкой на проективной \concept{короткой форме Вейерштрасса}, является класс $[0,1,0] = \{(0,y,0)\;|\; y\in \F\}$. Точка $[0:1:0]$ — проективное представление точки на бесконечности $\mathcal{O}$ в аффинной репрезентации. Проективное представление \concept{кривой в короткой форме Вейерштрасса} поэтому имеет то преимущество, что ему не нужен специальный символ для представления точки на бесконечности из аффинного определения.

\begin{example}
\label{ex:E1F5-projective}
 Чтобы получить интуицию о том, как выглядит эллиптическая кривая в проективной геометрии, рассмотрим кривую $E_{1,1}(\F_5)$ из \examplename{} \ref{E1F5}. Мы знаем, что в аффинной репрезентации множество точек на аффинной \concept{короткой форме Вейерштрасса} задаётся следующим образом:

\begin{equation}\label{eq:E1F5-affine}
E_{1,1}(\F_5) = \{\Oinf, (0,1),(2,1),(3,1),(4,2),(4,3),(0,4),(2,4),(3,4)\}
\end{equation}

Это определяется как множество всех пар $(x,y)\in \F_5\times \F_5$ таких, что выполняется уравнение аффинной \concept{короткой формы Вейерштрасса} $y^2 = x^3 + ax +b$ с $a=1$ и $b=1$.

Чтобы найти множество элементов $E_{1,1}(\F_5)$ в проективном представлении \concept{короткой формы Вейерштрасса} с теми же параметрами $a=1$ и $b=1$, мы должны вычислить множество проективных точек $[X:Y:Z]$ из проективной плоскости $\F_5\mathrm{P}^2$, удовлетворяющих следующему однородному кубическому уравнению для любого представителя $(x_1,y_1,z_1)\in [X:Y:Z]$:
\begin{equation}\label{eq:homogenous-cubic}
y_1^2z_1 = x_1^3 + 1\cdot x_1 z_1^2 + 1\cdot z_1^3
\end{equation}
Мы знаем из \secname{} \ref{sec:planes}, что проективная плоскость $\F_5\mathrm{P}^2$ содержит $5^2+5+1= 31$ элементов, поэтому мы можем подставить все элементы в уравнение \eqref{eq:homogenous-cubic} и проверить, совпадают ли обе стороны.

Например, рассмотрим проективную точку $[0:4:1]$. Мы знаем из \eqref{def:projective_coordinate}, что эта точка в проективной плоскости представляет следующую прямую в $\F_5^3\backslash\{(0,0,0)\}$:
\begin{equation}
\label{ex:projective_coordinate_1}
[0:4:1] = \{(0,4,1),(0,3,2),(0,2,3),(0,1,4)\}
\end{equation} 

Чтобы проверить, удовлетворяет ли $[0:4:1]$ уравнению \eqref{eq:homogenous-cubic}, мы можем подставить любого представителя, другими словами, любой элемент из \eqref{ex:projective_coordinate_1}. Каждый элемент удовлетворяет уравнению тогда и только тогда, когда все остальные элементы удовлетворяют уравнению. В качестве произвольного выбора подставим $(0,3,2)$ и получим следующий результат:
$$
3^2\cdot 2 = 0^3 + 1\cdot 0\cdot 2^2 + 1\cdot 2^3 \Leftrightarrow
3 = 3
$$
Это говорит нам, что аффинная точка $[0:4:1]$ действительно является решением уравнения \eqref{eq:homogenous-cubic}, но мы могли бы с таким же успехом подставить любого другого представителя элемента. Например, подстановка $(0,1,4)$ также удовлетворяет \eqref{eq:homogenous-cubic}: 
$$
1^2\cdot 4 = 0^3 + 1\cdot 0\cdot 4^2 + 1\cdot 4^3 \Leftrightarrow
4=4
$$
Чтобы найти проективное представление $E_{1,1}(\F_5)$, мы сначала замечаем, что проективная прямая на бесконечности $[1:0:0]$ не является точкой кривой ни на одной проективной \concept{короткой форме Вейерштрасса}, поскольку она не может удовлетворять определяющему уравнению в \eqref{def:projective_cubic_equation} ни при каких параметрах $a$ и $b$. Поэтому мы можем исключить её из рассмотрения. 

Более того, точка на бесконечности $[X:Y:0]$ может удовлетворять уравнению в \eqref{def:projective_cubic_equation} для любых $a$ и $b$ только при $X=0$, что означает, что единственной точкой на бесконечности, имеющей отношение к \concept{эллиптическим кривым в короткой форме Вейерштрасса}, является $[0:1:0]$, поскольку $[0:k:0]= [0:1:0]$ для всех $k\in\F^*$. Поэтому мы можем исключить все точки на бесконечности, кроме точки $[0:1:0]$.

Все оставшиеся точки — аффинные точки $[X:Y:1]$. Подставив все их в \eqref{eq:homogenous-cubic}, мы получаем множество всех проективных точек кривой следующим образом:

\begin{multline}
\label{E1F5_projective_points_set}
E_{1,1}(\F_5\mathrm{P}^2)=\{[0:1:0], [0:1:1], [2:1:1], [3:1:1], \\ [4:2:1], [4:3:1], [0:4:1], [2:4:1], [3:4:1]\}
\end{multline}

Если мы сравним это с аффинным представлением, мы увидим, что существует взаимно однозначное соответствие между точками в аффинном представлении в \eqref{eq:E1F5-affine} и аффинными точками в проективной геометрии, и что проективная точка $[0:1:0]$ представляет дополнительную точку $\Oinf$ в аффинном представлении.
\end{example} 

\begin{exercise}
Рассмотрите \examplename{} \ref{ex:E1F5-projective} и вычислите множество \eqref{E1F5_projective_points_set}, подставив все точки из проективной плоскости $\F_5\mathrm{P}^2$ в определяющее проективное уравнение \concept{короткой формы Вейерштрасса}.
\end{exercise}

\begin{exercise}
Вычислите проективное представление кривой \curvename{Tiny-jubjub} (\examplename{} \ref{TJJ13}) и логарифмический порядок её большой подгруппы простого порядка относительно образующего $[7:11:1]$ в проективных координатах.
\end{exercise}

\subsubsection{Проективный групповой закон}
\label{sec:projective_group_law}
Как мы видели в разделе \ref{sec:affine_group_law}, одно из ключевых свойств эллиптической кривой заключается в том, что она сопровождается определением группового закона на множестве её точек, описываемого геометрически правилом хорд и касательных (\defname{} \ref{def:chord-and-tangent}). Это правило было довольно интуитивным, за исключением выделенной точки на бесконечности, которая появлялась всякий раз, когда хорда или касательная не имели третьей точки пересечения с кривой.

Одна из ключевых особенностей проективных координат заключается в том, что в проективном пространстве гарантируется, что любая хорда всегда будет пересекать кривую в трёх точках, а любая касательная — в двух точках. Таким образом, геометрическая картина упрощается, поскольку нам не нужно рассматривать внешние символы и связанные с ними случаи. Цена, которую нужно заплатить за это математическое упрощение, — в том, что проективные координаты могут быть менее интуитивными для начинающих.

Можно показать, что точки эллиптической кривой в проективном пространстве образуют коммутативную группу относительно правила касательных и хорд так, что проективная точка $[0:1:0]$ является нейтральным элементом, а аддитивно обратной к точке $[X:Y:Z]$ является $[X:-Y:Z]$. Закон сложения обычно описывается \algname{} \ref{alg:projective_group_law}, минимизирующим число необходимых сложений и умножений в базовом поле. % https://www.hyperelliptic.org/EFD/precomp.pdf

\begin{algorithm}\caption{Закон сложения в проективной \concept{короткой форме Вейерштрасса}}
\label{alg:projective_group_law}
% https://en.wikibooks.org/wiki/Cryptography/Prime_Curve/Standard_Projective_Coordinates
\begin{algorithmic}[0]
\Require $[X_1:Y_1:Z_1],[X_2:Y_2:Z_2] \in E(\F\mathbb{P}^2)$
\Procedure{Add-Rule}{$[X_1:Y_1:Z_1],[X_2:Y_2:Z_2]$}
\If{$[X_1:Y_1:Z_1] == [0:1:0]$}
  \State $[X_3:Y_3:Z_3] \gets [X_2:Y_2:Z_2]$
\ElsIf{$[X_2:Y_2:Z_2] == [0:1:0]$}
  \State $[X_3:Y_3:Z_3] \gets [X_1:Y_1:Z_1]$
\Else
  \State $U_1 \gets Y_2\cdot Z_1$
  \State $U_2 \gets Y_1\cdot Z_2$
  \State $V_1 \gets X_2\cdot Z_1$
  \State $V_2 \gets X_1\cdot Z_2$
  \If{$V_1 == V_2$}
    \If{$U_1 \neq U_2$}
      $[X_3:Y_3:Z_3] \gets [0:1:0]$
    \Else
      \If{$Y_1 == 0$}
        $[X_3:Y_3:Z_3] \gets [0:1:0]$
      \Else
        \State $W \gets a\cdot Z_1^2 + 3\cdot X_1^2$
        \State $S \gets Y_1\cdot Z_1$
        \State $B \gets X_1\cdot Y_1\cdot S$
        \State $H \gets W^2 - 8\cdot B$
        \State $X' \gets 2\cdot H\cdot S$
        \State $Y' \gets W\cdot (4\cdot B - H) - 8\cdot Y_1^2\cdot S^2$
        \State $Z' \gets 8\cdot S^3$
        \State $[X_3:Y_3:Z_3] \gets [X':Y':Z']$
      \EndIf
    \EndIf
  \Else
    \State $U = U_1 - U_2$
    \State $V = V_1 - V_2$
    \State $W = Z_1\cdot Z_2$
    \State $A = U^2\cdot W - V^3 - 2\cdot V^2\cdot V_2$
    \State $X' = V\cdot A$
    \State $Y' = U\cdot(V^2\cdot V_2 - A) - V^3\cdot U_2$
    \State $Z' = V^3\cdot W$
    \State $[X_3:Y_3:Z_3]\gets [X':Y':Z']$
  \EndIf
\EndIf
\State \textbf{return} $[X_3:Y_3:Z_3]$
\EndProcedure
\Ensure $ [X_3:Y_3:Z_3] == [X_1:Y_1:Z_1] \oplus [X_2:Y_2:Z_2]$
\end{algorithmic}
\end{algorithm}
%\begin{example}[Polynomial evaluation on secret points]
%Since scalar multiplication is assumed to be a one way function, it can be used to encrypt computations. For example it can be used to proof identities of bounded degree polynomials (with some probability), without actually revealing the polynomials. To see what this means, consider the moon-jubjub curve $TJJ(\F_{13})$ from XXX\sme{add reference} and the set $\F_{13}[x]_{\leq 2}$ of all polynomials with coefficients in $\F_{13}$ and maximum degree $2$.

%Now assume that there are two parties $A$ and $B$ such that $A$ choose polynomial $P_A$ and $B$ chooses polynomial $P_B$ from $\F_{13}[x]_{\leq 2}$. The task is to check (with some probability) weather or not $P_A$ equals $P_B$ without actually revealing any information about the polynomials. 

%This task can be solved, by evaluating the polynomials at a secret point in the exponent of a (DFHM-PROPERTY) group and then compare the results.

%So we assume that there is some trusted third party $C$ that chooses a publicly known generator of a large prime-order subgroup of $TJJ(\F_{13})$, say  a secrete point $s\in\F_{13}$, say $s=2$. $C$ then c
%\end{example}
\begin{exercise} Снова рассмотрите \examplename{} \ref{ex:E1F5-projective}. Вычислите следующие выражения для проективных точек на $E_{1,1}(\F_5\mathrm{P}^2)$, используя \algname{} \ref{alg:projective_group_law}:
\begin{itemize}
\item $[0:1:0]\oplus [4:3:1]$
\item $[0:3:0]\oplus [3:1:2]$
\item $-[0:4:1]\oplus [3:4:1]$
\item $[4:3:1]\oplus [4:2:1]$
\end{itemize}
а затем решите уравнение $[X:Y:Z] \oplus [0:1:1]= [2:4:1]$ для некоторой точки $[X:Y:Z]$ из проективной \concept{короткой формы Вейерштрасса} $E_{1,1}(\F_5\mathrm{P}^2)$.
\end{exercise}

\begin{exercise}
Сравните аффинный закон сложения для \concept{короткой формы Вейерштрасса} с проективным правилом сложения. Какая ветвь в проективном правиле соответствует какому случаю в аффинном законе?
\end{exercise}

\subsubsection{Преобразования координат} Как мы видим, сравнивая примеры 
\ref{ex:E1F5-projective} и \ref{ex:E1F5-projective},\sme{same example twice} между аффинной и проективной репрезентациями \concept{кривой в короткой форме Вейерштрасса} существует тесная связь. Это не совпадение. На самом деле, с математической точки зрения, проективные и аффинные \concept{кривые в короткой форме Вейерштрасса} описывают одно и то же, поскольку для любых параметров $a$ и $b$ существует взаимно однозначное соответствие (изоморфизм) между обеими репрезентациями. 

Чтобы задать соответствие, пусть $E(\F)$ и $E(\F\mathbb{P}^2)$ — аффинная и проективная \concept{кривые в короткой форме Вейерштрасса}, определённые для одних и тех же параметров $a$ и $b$. Тогда функция в \eqref{eq:weierstrass-isomorphism-map} отображает точки из аффинной репрезентации в точки проективной репрезентации \concept{короткой формы Вейерштрасса}. Другими словами, если пара элементов поля $(x,y)$ удовлетворяет аффинному уравнению \concept{короткой формы Вейерштрасса} $y^2= x^3 + ax + b$, то все однородные координаты $(x_1,y_1,z_1)\in [x:y:1]$ удовлетворяют проективному уравнению \concept{короткой формы Вейерштрасса} $y_1^2\cdot z_1= x_1^3 + ay_1\cdot z_1^2 + b\cdot z_1^3$. 

\begin{equation}\label{eq:weierstrass-isomorphism-map}
I : E(\F) \to E(\F\mathbb{P}^2)\;:\;
\begin{array}{lcl}
(x,y)       &\mapsto & [x:y:1]\\
\mathcal{O} &\mapsto & [0:1:0]
\end{array}
\end{equation}
Это отображение является взаимно однозначным соответствием, что означает, что оно отображает ровно одну точку из аффинной репрезентации в одну точку проективной репрезентации. Поэтому можно обратить это отображение, чтобы отображать точки из проективной репрезентации в точки аффинной репрезентации \concept{короткой формы Вейерштрасса}. Обратное отображение задаётся следующим образом:
\begin{equation}
I^{-1} : E(\F\mathbb{P}^2)\to E(\F) \;:\; [X:Y:Z] \mapsto \begin{cases}
(\frac{X}{Z},\frac{Y}{Z}) & \text{ if } Z\neq 0\\
\mathcal{O} & \text{ if } Z=0
\end{cases}
\end{equation}
Заметьте, что единственной проективной точкой $[X:Y:Z]$ с $Z = 0$, удовлетворяющей уравнению в \ref{def:projective_cubic_equation}, является класс $[0:1:0]$. Ключевая особенность $I$ и его обратного отображения заключается в том, что оба отображения сохраняют групповую структуру, что означает, что нейтральный элемент отображается в нейтральный элемент $I(\mathcal{O})=[0:1:0]$, и что $I((x_1,y_1)\oplus (x_2,y_2))$ равно $I(x_1,y_1)\oplus I(x_2,y_2)$. То же самое верно для обратного отображения $I^{-1}$.

Отображения с такими свойствами называются \term{групповыми изоморфизмами (group isomorphisms)}, и, с математической точки зрения, существование функции $I$ означает, что аффинное и проективное определения \concept{эллиптических кривых в короткой форме Вейерштрасса} эквивалентны и представляют одну и ту же математическую сущность в двух различных представлениях. Реализации поэтому могут свободно выбирать между этими двумя репрезентациями. 


\section{Кривые Монтгомери}\label{sec:montgomery}
% https://eprint.iacr.org/2017/212.pdf
Аффинная и проективная \concept{короткие формы Вейерштрасса} — наиболее общие способы описания эллиптических кривых над полями характеристики больше $3$. Однако в определённых ситуациях может быть выгодно рассмотреть более специализированные представления эллиптических кривых, например, чтобы получить более быстрые алгоритмы для группового закона или скалярного умножения. 

Как мы увидим в этом разделе, так называемые \term{кривые Монтгомери (Montgomery curves)} — это подмножество всех эллиптических кривых, которые могут быть записаны в \term{форме Монтгомери (Montgomery form)}. Эти кривые допускают алгоритмы с постоянным временем для (специализаций) скалярного умножения на эллиптической кривой. 

Чтобы понять, что такое кривая Монтгомери в аффинной репрезентации, пусть $\F$ — простое поле порядка $p>3$, и пусть $A,B\in \F$ — два элемента поля такие, что $B\neq 0$ и $A^2 \neq \Zmod{4}{p}$. \term{Эллиптическая кривая Монтгомери (Montgomery elliptic curve)} $M(\F)$ над $\F$ в аффинной репрезентации — это множество всех пар элементов поля $(x,y)\in \F\times \F$, удовлетворяющих \term{кубическому уравнению Монтгомери (Montgomery cubic equation)} $B\cdot y^2 = x^3 + A\cdot x^2 + x$, вместе с выделенным символом $\Oinf$, называемым \term{точкой на бесконечности (point at infinity)}.

\begin{equation}
\label{eq:montgomery-curve}
M(\F) = \{(x,y)\in \F\times \F\;|\; B\cdot y^2 = x^3 + A\cdot x^2 + x  \} \bigcup \{\Oinf\}
\end{equation}

Несмотря на то что кривые Монтгомери выглядят иначе, чем \concept{кривые в короткой форме Вейерштрасса}, они являются лишь особым способом описания определённых \concept{кривых в короткой форме Вейерштрасса}. На самом деле, каждая кривая в аффинной форме Монтгомери может быть преобразована в эллиптическую кривую в \concept{короткой форме Вейерштрасса}. Чтобы увидеть это, предположим, что кривая задана в форме Монтгомери $B y^2 = x^3 + A x^2 + x$. Ассоциированная \concept{короткая форма Вейерштрасса} тогда определяется следующим образом:

\begin{equation}\label{eq:montgomery-to-weierstrass}
y^2 = x^3 + \frac{3-A^2}{3\cdot B^2}\cdot x + \frac{2\cdot A^3-(\Zmod{9}{p})\cdot A}{(\Zmod{27}{p})\cdot B^3}
\end{equation}

С другой стороны, не каждую эллиптическую кривую $E(\F)$ над простым полем $\F$ характеристики $p>3$, заданную в \concept{короткой форме Вейерштрасса} $y^2 = x^3 + a x + b$, можно преобразовать в форму Монтгомери. Чтобы \concept{кривая в короткой форме Вейерштрасса} была кривой Монтгомери, должны выполняться следующие условия:

\begin{definition}\label{def:montgomery}
\defvsep{}

\begin{itemize}
\item Число точек на $E(\F)$ делится на $4$.
\item Многочлен $z^3 + a z + b \in \F[z]$ имеет по крайней мере один корень $z_0\in\F$.
\item $3z_0^2 + a$ является квадратичным вычетом в $\F^*$.
\end{itemize}
\end{definition}

Когда эти условия выполнены, тогда для $s=({\sqrt{3z_0^{2}+a}})^{-1}$ кривая Монтгомери определяется следующим уравнением:
\begin{equation}\label{eq:montgomery-form}
sy^{2}=x^{3}+(3z_0 s)x^{2}+x
\end{equation}

В следующем примере мы снова рассмотрим кривую \curvename{Tiny-jubjub} и покажем, что она на самом деле является кривой Монтгомери.
\begin{example}\label{TJJ13-montgomery}
Рассмотрим простое поле $\F_{13}$ и кривую \curvename{Tiny-jubjub} \TJJ{} из \examplename{} \ref{TJJ13}. Чтобы увидеть, что она является кривой Монтгомери, мы должны проверить требования из \defname{} \ref{def:montgomery}: 

Поскольку порядок \TJJ{} равен $20$, что делится на $4$, первое требование выполнено.

Что касается второго критерия, поскольку $a=8$ и $b=8$, мы должны проверить, что многочлен $P(z) = z^3 + 8z + 8$ имеет корень в $\F_{13}$. Чтобы увидеть это, мы просто вычисляем $P$ во всех числах $z\in \F_{13}$ и находим, что $P(4)=0$, так что корень задаётся $z_0=4$. 

На последнем шаге мы должны проверить, что $3\cdot z_0^2 + a$ имеет корень в $\F_{13}^*$. Мы вычисляем следующим образом:
\begin{align*}
3z_0^2 + a & = 3\cdot 4^2 + 8 \\
           & = 3 \cdot 3 + 8 \\
           & = 9 + 8 \\
           & = 4
\end{align*}

Чтобы увидеть, что $4$ является квадратичным вычетом в $\F_{13}$, мы используем критерий Эйлера (\ref{eq: Euler_criterion}), чтобы вычислить символ Лежандра для $4$. Мы получаем следующее:

$$
\left(\frac{4}{13}\right) = 4^{\frac{13-1}{2}} = 4^6 = 1
$$ 
Это означает, что $4$ действительно имеет корень в $\F_{13}$. На самом деле, вычисление корня из $4$ в $\F_{13}$ легко, поскольку целочисленный корень $2$ из $4$ также является одним из его корней в $\F_{13}$. Другой корень задаётся $13-2=11$.

Поскольку все требования выполнены, мы показали, что \TJJ{} действительно является кривой Монтгомери, и мы можем использовать \eqref{eq:montgomery-form}, чтобы вычислить её ассоциированную форму Монтгомери следующим образом:
\begin{align*}
s & = \left(\sqrt{3\cdot z_0^2 +8}\right)^{-1} \\
  & = 2^{-1} & \text{\# Fermat's little theorem} \\
  & = 2^{13-2} & \text{\# }\Zmod{2048}{13} = 7\\
  & = 7
\end{align*}

Определяющее уравнение для формы Монтгомери кривой \curvename{Tiny-jubjub} задаётся следующим образом:
\begin{align*}
sy^{2} & =x^{3}+(3z_0 s)x^{2}+x  & \Rightarrow\\
7\cdot y^{2} & =x^{3}+(3\cdot 4 \cdot 7)x^{2}+x &\Leftrightarrow\\
7\cdot y^{2} & =x^{3}+6x^{2}+x
\end{align*}
Итак, мы получаем определяющие параметры $B= 7$ и $A=6$, и можем записать кривую \curvename{Tiny-jubjub} в аффинной репрезентации Монтгомери следующим образом:
\begin{equation}\label{eq:TJJ13-montgomery-representation}
\mathit{TJJ\_13} = \{(x,y)\in \F_{13}\times \F_{13}\;|\; 7\cdot y^{2} =x^{3}+6x^{2}+x \}\bigcup \{\Oinf\}
\end{equation}

Теперь, когда у нас есть абстрактное определение кривой \curvename{Tiny-jubjub} в форме Монтгомери, мы можем вычислить множество точек, подставив все пары $(x,y)\in\F_{13}\times \F_{13}$, аналогично тому, как мы вычисляли точки кривой в её \concept{короткой форме Вейерштрасса} \eqref{eq:TJJ13-weierstrass}\sme{check reference}. Мы получаем следующее:
\begin{multline}
\label{tjj_13_montgomery_points_set}
\mathit{M\_TJJ\_13} = \{\Oinf, (0, 0),(1, 4),(1, 9),(2, 4),(2, 9),(3, 5),(3, 8),(4, 4),(4, 9),\\ (5, 1),(5, 12),(7, 1),(7, 12),(8, 1),(8, 12),(9, 2),(9, 11),(10, 3),(10, 10)\}
\end{multline}

Мы можем проверить результаты с помощью Sage:

\begin{sagecommandline}
sage: F13 = GF(13)
sage: L_MTJJ = []
....: for x in F13:
....:     for y in F13:
....:         if F13(7)*y^2 == x^3 + F13(6)*x^2 +x:
....:             L_MTJJ.append((x,y))
sage: MTJJ = Set(L_MTJJ)
sage: # не вычисляет точку на бесконечности
\end{sagecommandline}
\end{example}
\begin{exercise}
Рассмотрите \examplename{} \ref{TJJ13-montgomery} и вычислите множество в \eqref{tjj_13_montgomery_points_set}, подставив каждую пару элементов поля $(x,y)\in \F_{13}\times \F_{13}$ в определяющее уравнение Монтгомери.
\end{exercise}
\begin{exercise}
Рассмотрите эллиптическую кривую $E_{1,1}(\F_5)$ из \examplename{} \ref{E1F5} и покажите, что $E_{1,1}(\F_5)$ не является кривой Монтгомери.
\end{exercise}
\begin{exercise}
Рассмотрите эллиптическую кривую \curvename{secp256k1} из \examplename{} \ref{secp256k1} и покажите, что \curvename{secp256k1} не является кривой Монтгомери.
\end{exercise}

\subsection{Аффинное преобразование координат Монтгомери} Сравнивая репрезентацию Монтгомери в \eqref{eq:TJJ13-montgomery-representation} с \concept{короткой формой Вейерштрасса} той же кривой в \eqref{eq:TJJ13-weierstrass}, мы видим, что между точками кривых в этих двух примерах существует взаимно однозначное соответствие. Это не случайность. На самом деле, если $M_{A,B}$ — кривая Монтгомери, а $E_{a,b}$ — \concept{кривая в короткой форме Вейерштрасса} с $a = \frac{3-A^2}{3B^2}$ и $b= \frac{2A^2 -9A}{27B^3}$, то следующая функция отображает все точки в репрезентации Монтгомери в точки \concept{короткой формы Вейерштрасса}:
\begin{equation}
I: M_{A,B} \to E_{a,b}\; : \; (x,y) \mapsto \left(\frac{3x + A}{3B}, \frac{y}{B}\right)
\end{equation}
Точка на бесконечности формы Монтгомери отображается в точку на бесконечности \concept{короткой формы Вейерштрасса}. Это отображение является взаимно однозначным соответствием (изоморфизмом), и его обратное отображение задаётся следующим уравнением (где $z_0$ — корень многочлена $z^3 + a z + b \in \F[z]$ и $s=({\sqrt{3z_0^{2}+a}})^{-1}$).
\begin{equation}
I^{-1}: E_{a,b} \to M_{A,B}\; : \; (x,y) \mapsto \left(s\cdot(x-z_0), s\cdot y\right)
\end{equation}
Точка на бесконечности \concept{короткой формы Вейерштрасса} отображается в точку на бесконечности формы Монтгомери. Используя это отображение, реализации кривых Монтгомери могут свободно переходить между \concept{короткой формой Вейерштрасса} и формой Монтгомери. 
 
\begin{example} Снова рассмотрим кривую \curvename{Tiny-jubjub}. В \ref{eq:TJJ13-weierstrass} мы определили её \concept{короткую форму Вейерштрасса}, а в \examplename{} \ref{eq:TJJ13-montgomery-representation} вывели её репрезентацию Монтгомери. 

Чтобы увидеть, как работает преобразование координат $I$ в этом примере, давайте отобразим точки из репрезентации Монтгомери в точки \concept{короткой формы Вейерштрасса}. Подстановка, например, точки $(0,0)$ из репрезентации Монтгомери \eqref{eq:TJJ13-montgomery-representation} в $I$ даёт следующее:
\begin{align*}
I(0,0) & = \left(\frac{3\cdot 0 + A}{3B}, \frac{0}{B}\right) \\
          & = \left(\frac{3\cdot 0 + 6}{3\cdot 7}, \frac{0}{7}\right) \\
          & = \left(\frac{6}{8}, 0\right) \\
          & = \left(4, 0\right) \\
\end{align*}

Как мы видим, точка Монтгомери $(0,0)$ отображается в самообратную точку $(4,0)$ \concept{короткой формы Вейерштрасса}. С другой стороны, мы можем использовать наши вычисления $s=7$ и $z_0=4$ из \examplename{} \ref{TJJ13-montgomery}, чтобы вычислить обратное отображение $I^{-1}$, которое отображает точки \concept{короткой формы Вейерштрасса} в точки формы Монтгомери. Подставив, например, $(4,0)$, мы получаем следующее:
\begin{align*}
I^{-1}(4,0) & = \left(s\cdot(4-z_0), s\cdot 0\right)\\
               & = \left(7\cdot(4-4), 0\right)\\
               & = (0,0)
\end{align*}

Как и ожидалось, обратное отображение переводит точку \concept{короткой формы Вейерштрасса} обратно туда, откуда она взялась в форме Монтгомери. Мы можем использовать Sage, чтобы проверить, что наше вычисление $I$ верно:
\begin{sagecommandline}
sage: # Вычислить I формы Монтгомери:
sage: L_I_MTJJ = []
sage: for (x,y) in L_MTJJ: # LMTJJ как определено ранее                                   
....:     v = (F13(3)*x + F13(6))/(F13(3)*F13(7))
....:     w = y/F13(7)
....:     L_I_MTJJ.append((v,w))
sage: I_MTJJ = Set(L_I_MTJJ)
sage: # Вычисление \concept{короткой формы Вейерштрасса}
sage: C_WTJJ = EllipticCurve(F13,[8,8]) 
sage: L_WTJJ = [P.xy() for P in C_WTJJ.points() if P.order() > 1]
sage: WTJJ = Set(L_WTJJ)
sage: # проверить I(Монтгомери) == Вейерштрасс
sage: WTJJ == I_MTJJ
sage: # проверить обратное отображение I^(-1)
sage: L_IINV_WTJJ = []
sage: for (v,w) in L_WTJJ:
....:     x = F13(7)*(v-F13(4))
....:     y = F13(7)*w
....:     L_IINV_WTJJ.append((x,y))
sage: IINV_WTJJ = Set(L_IINV_WTJJ)
sage: MTJJ == IINV_WTJJ
\end{sagecommandline}
\end{example}

\subsection{Групповой закон Монтгомери} Мы видели, что кривые Монтгомери — частные случаи \concept{кривых в короткой форме Вейерштрасса}. Как таковые, они имеют групповую структуру, определённую на множестве своих точек, которая также может быть выведена из правила хорд и касательных. В соответствии с \concept{кривыми в короткой форме Вейерштрасса}, можно показать, что тождество $x_1=x_2$ влечёт $y_2=\pm y_1$, что означает, что следующие правила являются полным описанием группового закона эллиптической кривой:

\begin{definition}[\deftitle{Групповой закон Монтгомери}]\label{def:montgomery-group-law}
\defvsep{}

\begin{itemize}
\item (Нейтральный элемент) Точка на бесконечности $\Oinf$ является нейтральным элементом.
\item (Обратный элемент) Обратным для $\mathcal{O}$ является $\mathcal{O}$. Для любой другой точки кривой $(x,y) \in M(\F_q)\backslash \{\mathcal{O}\}$ обратный элемент задаётся $(x,-y)$.
\item (Групповой закон) Для любых двух точек кривой $P, Q \in M(\F_q)$ групповой закон определяется одним из следующих случаев:
\begin{enumerate}
\item (Нейтральный элемент) Если $Q=\Oinf$, то сумма определяется как $P\oplus Q=P$.
\item (Обратные элементы) Если $P=(x,y)$ и $Q=(x,-y)$, групповой закон определяется как $P\oplus Q=\Oinf$.
\item (Правило касательной) Если $P=(x,y)$ с $y\neq 0$, групповой закон $P\oplus P=(x',y')$ определяется следующим образом:
$$
\begin{array}{llr}
x' = (\frac{3x^2 + 2A x +1}{2By})^2\cdot B - 2x - A &,&
y' = \frac{3x^2 + 2A x +1}{2By}(x-x') - y
\end{array} 
$$
\item (Правило хорды) Если $P=(x_1,y_1)$ и $Q=(x_2,y_2)$ такие, что $x_1 \neq x_2$, групповой закон $R=P\oplus Q$ с $R=(x',y')$ определяется следующим образом:
$$
\begin{array}{llr}
x' = (\frac{y_2-y_1}{x_2-x_1})^2B - (x_1 + x_2) - A &, &
y' = \frac{y_2-y_1}{x_2-x_1}(x_1-x') - y_1
\end{array} 
$$
\end{enumerate}
\end{itemize}
\end{definition}
\begin{exercise}
Рассмотрите коммутативную группу $(\mathit{M\_TJJ\_13},\oplus)$ кривой \curvename{Tiny-jubjub} в форме Монтгомери из \examplename{} \eqref{tjj_13_montgomery_points_set}. 
\begin{enumerate}
\item Вычислите обратные элементы для $(1,9)$, $\Oinf$, $(7,12)$ и $(4,9)$.
\item Решите уравнение $x \oplus (3,8) = (10,3) $ для некоторого $x\in \mathit{M\_TJJ\_13}$.
\end{enumerate}
Выберите некоторый элемент $x\in \mathit{M\_TJJ\_13}$ и проверьте, является ли $x$ образующим $\mathit{M\_TJJ\_13}$. Если $x$ не является образующим, повторяйте, пока не найдёте некоторый образующий $x$. Запишите $\mathit{M\_TJJ\_13}$ в логарифмическом порядке относительно $x$.
\end{exercise}
\begin{exercise}
Рассмотрите кривую \curvename{alt\_bn128} из примера \ref{BN128}. Покажите, что эта кривая не является кривой Монтгомери.
\end{exercise}
%\subsubsection{Projective Montgomery Form}
%As with more general curves in \concept{short Weierstrass} form, we can look at Montgomery curves in projective space. To see how such a curve looks in projective coordinates is, let $\F_{q}$ be a finite field and $A,B\in \F_q$ two field elements such that $B\neq 0$ and $A^2\neq 4$. Then a \term{Montgomery elliptic curve} $E/\F_q$ over $\F_q$ in its projective representation is the set
%\begin{equation}
%\label{def_short_weierstrass_curve}
%E/\F_q\mathbb{P}^2 = \{[X:Y:Z]\in \F_q\mathbb{P}^2\;|\; B\cdot Y^2 \cdot Z = X^3 + A\cdot X^2\cdot Z + X\cdot Z^2  \}
%\end{equation}
%of all points $[X:Y:Z]\in \F_q\mathbb{P}^2$ from the projective plane that satisfy the \term{homogenous} cubic equation $B\cdot Y^2 \cdot Z = X^3 + A\cdot X^2\cdot Z + X\cdot Z^2$.

% TODO: https://maths-people.anu.edu.au/~brent/pd/Subramanya-thesis.pdf
% x-coordinate only aka differential arithmetics on Montgomery curves
\section{Кривые в крученой форме Эдвардса}\label{sec:edwards}
Как мы видели в \defname{} \ref{def:chord-tangent-algebra} и \defname{} \ref{def:montgomery-group-law}, и \concept{кривые в короткой форме Вейерштрасса}, и кривые Монтгомери имеют довольно сложные групповые законы, поскольку необходимо различать множество случаев. Это может добавить сложности программной реализации, потому что каждый случай переводится в дополнительную ветвь компьютерной программы. Однако в контексте разработки SNARK используются вычислительные модели для ограниченных вычислений, в которых ветвления программы нежелательно дороги (подробнее об этом в \secname{} \ref{sec:R1CS} и \ref{sec:circuits}). Чтобы сделать эллиптические кривые ``SNARK-дружественными'', поэтому выгодно искать кривые с групповым законом, не требующим ветвлений и использующим как можно меньше операций в поле.

Так называемые \term{SNARK-дружественные \concept{кривые в крученой форме Эдвардса} (SNARK-friendly twisted Edwards curves)} особенно полезны здесь, поскольку эти кривые имеют компактный и легко реализуемый групповой закон, который работает для всех точек, включая точку на бесконечности. Реализация этого закона не требует ветвлений. 

% https://eprint.iacr.org/2008/013.pdf
Чтобы понять, как выглядит \term{\concept{кривая в крученой форме Эдвардса} (twisted Edwards curve)} в аффинной форме, пусть $\F$ — конечное поле характеристики $>3$, и пусть $a,d\in \F\backslash\{0\}$ — два ненулевых элемента поля такие, что $a\neq d$. \term{\concept{Эллиптическая кривая в крученой форме Эдвардса} (twisted Edwards elliptic curve)} в аффинной репрезентации — это множество всех пар $(x,y)$ из $\F\times \F$, удовлетворяющих уравнению \concept{крученой формы Эдвардса} $a\cdot x^2+y^2= 1+d\cdot x^2y^2$:
\begin{equation}\label{eq:twisted-edwards}
E(\F)=\{(x,y)\in\F\times\F\;|\; a\cdot x^2+y^2= 1+d\cdot x^2y^2\}
\end{equation} 
\concept{Кривая в крученой форме Эдвардса} называется \term{SNARK-дружественной \concept{кривой в крученой форме Эдвардса} (SNARK-friendly twisted Edwards curve)}, если параметр $a$ является квадратичным вычетом, а параметр $d$ — квадратичным невычетом.

Как мы видим из определения, аффинные \concept{кривые в крученой форме Эдвардса} выглядят несколько иначе, чем \concept{кривые в короткой форме Вейерштрасса}, поскольку их аффинная репрезентация не требует специального символа для представления точки на бесконечности. На самом деле, пара $(0,1)$ всегда является точкой на любой \concept{кривой в крученой форме Эдвардса}, и она играет роль точки на бесконечности.

Несмотря на различия во внешнем виде, \concept{кривые в крученой форме Эдвардса} эквивалентны кривым Монтгомери в том смысле, что для каждой \concept{кривой в крученой форме Эдвардса} существует кривая Монтгомери и способ отобразить точки одной кривой на другую (и наоборот) взаимно однозначно. 

Чтобы увидеть это, предположим, что дана кривая в \concept{крученой форме Эдвардса}. Ассоциированная кривая Монтгомери тогда определяется уравнением Монтгомери:
\begin{equation}
\frac{4}{a-d} y^2 = x^3 + \frac{2(a+d)}{a-d}\cdot x^2 + x 
\end{equation}

С другой стороны, кривая Монтгомери $By^{2}=x^{3}+Ax^{2}+x$ такая, что $B\neq 0$ и $A^2\neq 4$, порождает \concept{кривую в крученой форме Эдвардса}, определяемую следующим уравнением:
\begin{equation}\label{eq:montgomery-to-twisted-edwards}
(\frac{A+2}{B})x^2+y^2= 1+(\frac{A-2}{B})x^2y^2
\end{equation}

\begin{example}\label{ex:TJJ13-twisted-edwards}Снова рассмотрим кривую \curvename{Tiny-jubjub} из \examplename{} \ref{TJJ13}. Мы знаем из \examplename{} \ref{TJJ13-montgomery}, что она является кривой Монтгомери, и, поскольку кривые Монтгомери эквивалентны \concept{кривым в крученой форме Эдвардса}, мы хотим записать эту кривую в \concept{крученой форме Эдвардса}. Мы используем уравнение \eqref{eq:montgomery-to-twisted-edwards} и вычисляем параметры $a$ и $d$ следующим образом:
\begin{align*}
a & = \frac{A+2}{B} & \text{\# подставить A=6 и B=7}\\
  & = \frac{8}{7} = 3 & \text{\# } 7^{-1}= 2 \\
  \\
d & = \frac{A-2}{B} \\
  & = \frac{4}{7} = 8 
\end{align*}

Thus, we get the defining parameters  $a= 3$ and $d=8$. 

Поскольку наша цель — использовать эту кривую позже в реализациях SNARK для работы вручную, давайте покажем, что \curvename{Tiny-jubjub} также является SNARK-дружественной \concept{кривой в крученой форме Эдвардса}. Чтобы увидеть это, мы должны показать, что $a$ — квадратичный вычет, а $d$ — квадратичный невычет. Поэтому мы вычисляем символы Лежандра для $a$ и $d$ с помощью критерия Эйлера. Мы получаем следующее:
\begin{align*}
\left(\frac{3}{13}\right) &= 3^{\frac{13-1}{2}} \\
                          & = 3^6 
                            = 1\\
                          \\
\left(\frac{8}{13}\right) &= 8^{\frac{13-1}{2}} \\
                          & = 8^6 
                            = 12
                            = -1                     
\end{align*}

Это доказывает, что \curvename{Tiny-jubjub} является SNARK-дружественной. Мы можем записать кривую \curvename{Tiny-jubjub} в аффинной \concept{крученой форме Эдвардса} следующим образом:
\begin{equation}\label{eq:TJJ13-twisted-edwards}
\mathit{TJJ\_13} = \{(x,y)\in \F_{13}\times \F_{13}\;|\; 3\cdot x^{2} + y^2 =1+ 8\cdot x^{2}\cdot y^2 \}
\end{equation}

Теперь, когда у нас есть абстрактное определение нашей кривой \curvename{Tiny-jubjub} в \concept{крученой форме Эдвардса}, мы можем вычислить множество точек, подставив все пары $(x,y)\in\F_{13}\times \F_{13}$, аналогично тому, как мы вычисляли точки кривой в её \concept{короткой форме Вейерштрасса} или репрезентации Эдвардса. Мы получаем следующее:
\begin{equation}
\begin{split}
\mathit{TE\_TJJ\_13} = \{(0, 1),(0, 12),(1, 2),(1, 11),(2, 6),(2, 7),(3, 0),(5, 5),(5, 8),(6, 4),\\
(6, 9),(7, 4),(7, 9),(8, 5),(8, 8),(10, 0),(11, 6),(11, 7),(12, 2),(12, 11)\}
\end{split}
\end{equation}

Мы перепроверяем наши результаты с помощью Sage:

\begin{sagecommandline}
sage: F13 = GF(13)
sage: L_ETJJ = []
....: for x in F13:
....:     for y in F13:
....:         if F13(3)*x^2 + y^2 == 1+ F13(8)*x^2*y^2:
....:             L_ETJJ.append((x,y))
sage: ETJJ = Set(L_ETJJ)
\end{sagecommandline}
\end{example}
\subsection{Групповой закон крученой формы Эдвардса}
\label{sec:twisted_ed_group_law} Как мы видели, \concept{кривые в крученой форме Эдвардса} эквивалентны кривым Монтгомери, и, как таковые, они также имеют групповой закон. Однако, в отличие от кривых Монтгомери и \concept{короткой формы Вейерштрасса}, групповой закон SNARK-дружественных \concept{кривых в крученой форме Эдвардса} может быть описан единственным вычислением, которое работает во всех случаях, даже если мы складываем нейтральный элемент, обратный элемент, или если нам нужно удвоить точку. 

Чтобы увидеть, как выглядит групповой закон \concept{крученой формы Эдвардса}, пусть $(x_1, y_1)$, $(x_2, y_2)$ — две точки на кривой Эдвардса $E(\F)$. Сумма этих точек тогда задаётся следующим уравнением:

\begin{equation}\label{twisted-edwards-group-law}
(x_1, y_1) \oplus (x_2, y_2) =\left(\frac{x_1y_2+y_1x_2}{1 +dx_1x_2y_1y_2},\frac{y_1y_2-ax_1x_2}{1-dx_1x_2y_1y_2}\right)
\end{equation}

Чтобы увидеть, каков нейтральный элемент группового закона, сначала заметим, что точка $(0,1)$ является решением уравнения \concept{крученой формы Эдвардса} $a\cdot x^{2} + y^2 =1+ d\cdot x^{2}\cdot y^2$ для любых параметров $a$ и $d$, и, следовательно, $(0,1)$ является точкой на любой \concept{кривой в крученой форме Эдвардса}. Можно показать, что $(0,1)$ служит нейтральным элементом, а обратной к точке $(x_1, y_1)$ является точка $(-x_1, y_1)$.
\begin{example}
\label{example:TETJJ13}
 Снова рассмотрим кривую \curvename{Tiny-jubjub} в форме Эдвардса из \eqref{TJJ13-twisted-edwards}. Как мы видели, эта кривая задаётся следующим образом:
\begin{multline}
\mathit{TE\_TJJ\_13} = \{(0, 1),(0, 12),(1, 2),(1, 11),(2, 6),(2, 7),(3, 0),(5, 5),(5, 8),(6, 4),\\
(6, 9),(7, 4),(7, 9),(8, 5),(8, 8),(10, 0),(11, 6),(11, 7),(12, 2),(12, 11)\}
\end{multline}
Чтобы понять закон сложения \concept{крученой формы Эдвардса}, давайте сначала сложим нейтральный элемент $(0,1)$ с самим собой. Мы применяем групповой закон из \eqref{twisted-edwards-group-law} и получаем следующее:
\begin{align*}
(0, 1) \oplus (0, 1) &= \left(\frac{0\cdot 1+1 \cdot 0}{1 +8\cdot0\cdot 0\cdot 1\cdot 1},\frac{1\cdot 1-3\cdot 0\cdot 0}{1-8\cdot 0\cdot 0\cdot 1\cdot 1}\right)\\
                     & = (0,1)
\end{align*}
Итак, как и ожидалось, сложение нейтрального элемента с самим собой даёт нейтральный элемент. 

Теперь давайте сложим нейтральный элемент с какой-нибудь другой точкой кривой. Мы получаем следующее:
\begin{align*}
(0, 1) \oplus (8, 5) &= \left(\frac{0\cdot 5+1 \cdot 8}{1 +8\cdot0\cdot 8\cdot 1\cdot 5},\frac{1\cdot 5 - 3\cdot 0\cdot 8}{1-8\cdot 0\cdot 8\cdot 1\cdot 5}\right)\\
                     & = (8,5)
\end{align*}

Снова, как и ожидалось, сложение нейтрального элемента с любым элементом даёт этот элемент. 

Для любой точки кривой $(x,y)$ мы знаем, что её обратная задаётся $(-x,y)$. Чтобы увидеть, как работает сложение точки с её обратной, мы вычисляем следующим образом:
\begin{align*}
(5, 5) \oplus (8, 5) &= \left(\frac{5\cdot 5+5 \cdot 8}{1 +8\cdot 5\cdot 8\cdot 5\cdot 5},\frac{5\cdot 5 - 3\cdot 5\cdot 8}{1-8\cdot 5\cdot 8\cdot 5\cdot 5}\right)\\
                     &= \left(\frac{12+1}{1 +5},\frac{12 - 3}{1-5}\right)\\
                     &= \left(\frac{0}{6},\frac{12 + 10}{1+8}\right)\\
                     &= \left(0,\frac{9}{9}\right)\\
                     &=  (0,1)
\end{align*}

Сложение точки кривой с её обратной даёт нейтральный элемент, как и ожидалось. 


Как мы видели из этих примеров, закон сложения \concept{крученой формы Эдвардса} особенно хорошо и единообразно обрабатывает крайние случаи.
\end{example}
\begin{exercise} Рассмотрите коммутативную группу $(TE\_TJJ\_13,\oplus)$ из \examplename{} \ref{example:TETJJ13}.
\begin{enumerate}
\item Вычислите обратные элементы для $(1,11)$, $(0,1)$, $(3,0)$ и $(5,8)$.
\item Решите уравнение $x \oplus (5,8) = (1,11) $ для некоторого $x\in \mathit{TE\_TJJ\_13}$.
\end{enumerate}
Выберите некоторый элемент $x\in \mathit{TE\_TJJ\_13}$ и проверьте, является ли $x$ образующим $\mathit{TE\_TJJ\_13}$. Если $x$ не является образующим, повторяйте, пока не найдёте некоторый образующий $x$. Запишите $\mathit{TE\_TJJ\_13}$ в логарифмическом порядке относительно $x$.
\end{exercise}
%\begin{example}[Non \concept{twisted Edwards} curves have order 4 points]
%In this example, we will show that every Edwards curve has a point of order $4$. To see that let $E$ be an arbitrary Edwards curve ($a=1$). Then the point $(1,0)$ is on that curve, since $1^2+0^2= 1+d 1^2 0^2$. We compute 
%\begin{align*}
%[4](1,0) = \\
%[2]([2](1,0))=\\
%[2]\left(\frac{1 0 +1 0}{1 +d 1 1 0 0},\frac{00-1\cdot 1}{1-d1 1 00}\right)=\\
%[2](0,-1)=\\
%\left(\frac{0(-1)+(-1)0}{1 +d 0 0 (-1)(-1)},\frac{(-1)(-1)-00}{1-d00(-1)(-1)}\right) =\\
%(0,1)
%\end{align*} 

%Now having seen that every Edwards curve has point of order $4$, we can deduce that the order of every Edwards curve must contain $4$ as a factor. This is restrictive and in fact Edwards curves are rare. 
%\end{example}
\section{Спаривания на эллиптических кривых}
\label{sec:elliptic_curve_pairings}
% TODO: analyse the 2-torsion case
Как введено в \eqref{pairing-map}, некоторые группы сопровождаются понятием отображения спаривания. В этом разделе мы обсуждаем \term{спаривания на эллиптических кривых (pairings on elliptic curves)}, которые лежат в основе нескольких zk-SNARK и других схем доказательств с нулевым разглашением, по существу потому, что они позволяют разбить вычисления ``в показателе'' (см. \examplename{} \ref{ex:in-the-exponent}) на различные части, вычисляемые разными сторонами.\footnote{Более подробное введение в спаривания на эллиптических кривых можно найти, например, в \chaptname{} 6, \secname{} 6.8 и 6.9 в \cite{hoffstein-2008}.}

Мы начинаем с определения некоторых важных подгрупп так называемой полной группы кручения эллиптической кривой. Затем мы вводим спаривание Вейля эллиптической кривой и описываем алгоритм Миллера, который делает эти спаривания эффективно вычислимыми.

\subsection{Степени вложения} Как мы увидим далее, каждая эллиптическая кривая порождает отображение спаривания. Однако мы также увидим в \examplename{} \ref{ex:pairings_on_secp256k1}, что не каждое такое спаривание может быть эффективно вычислено. Чтобы отличить кривые с эффективно вычислимыми спариваниями от остальных, нам нужно начать с введения в так называемую \term{степень вложения (embedding degree)} кривой. 

Чтобы понять, что такое степень вложения эллиптической кривой, пусть $\F$ — конечное поле порядка $|\F|=q$, $E(\F)$ — эллиптическая кривая над $\F$, и пусть $r$ — простой множитель порядка $n$ кривой $E(\F)$. Степень вложения $E(\F)$ относительно $r$ — это наименьшее целое число $k$ такое, что выполняется следующее уравнение:
\begin{equation}
\label{def:embedding-degree}
r\,|\, q^k-1
\end{equation}
Малая теорема Ферма \eqref{fermats-little-theorem} влечёт, что для каждой эллиптической кривой всегда существует степень вложения $k(r)$ и для любого множителя $r$ порядка $n$ кривой, поскольку $k=r-1$ всегда является решением сравнения $\kongru{q^{k}}{1}{r}$. Это означает, что остаток от целочисленного деления $q^{r-1}-1$ на $r$ равен $0$.

\begin{notation} Пусть $\F$ — конечное поле порядка $q$, и пусть $E(\F)$ — эллиптическая кривая над $\F$ такая, что $r$ — простой множитель порядка $E(\F)$. Мы пишем $k(r)$ для степени вложения $E(\F)$ относительно $r$.
\end{notation}

\begin{example} Чтобы лучше понять степень вложения, рассмотрим эллиптическую кривую $E_{1,1}(\F_5)$ из \examplename{} \ref{E1F5}. Мы знаем, что порядок $E_{1,1}(\F_5)$ равен $9$, и, поскольку единственный простой множитель $9$ — это $3$, мы вычисляем степень вложения $E_{1,1}(\F_5)$ относительно $3$. 

Чтобы найти степень вложения, мы должны найти наименьшее целое $k$ такое, что $3$ делит $q^k-1= 5^k-1$. Мы пробуем и увеличиваем, пока не найдём подходящее $k$. 

\begin{align*}
k=1 &\text{ : } 5^1-1 = 4 & \text{ not divisible by } 3\\ 
k=2 &\text{ : } 5^2-1 = 24 & \text{ divisible by } 3
\end{align*} 

Это показывает, что степень вложения эллиптической кривой $E_{1,1}(\F_5)$ относительно простого множителя $3$ порядка $E_{1,1}(\F_5)$ равна $2$.
\end{example}

\begin{example}\label{ex:TJJ13-embedding-degree} Рассмотрим кривую \curvename{Tiny-jubjub} \TJJ{} из \examplename{} \ref{TJJ13}. Мы знаем, что порядок \TJJ{} равен $20$, и что порядок поэтому имеет два простых множителя: большой простой множитель $5$ и малый простой множитель $2$. 

Мы начинаем с вычисления степени вложения \TJJ{} относительно большого простого множителя $5$. Чтобы найти эту степень вложения, мы должны найти наименьшее целое $k$ такое, что $5$ делит $q^k-1= 13^k-1$. Мы пробуем и увеличиваем, пока не найдём подходящее $k$. 
\begin{align*}
k=1 &\text{: } 13^1-1 = 12 & \text{ not divisible by } 5\\ 
k=2 &\text{: } 13^2-1 = 168 & \text{ not divisible by } 5\\ 
k=3 &\text{: } 13^3-1 = 2196 & \text{ not divisible by } 5\\ 
k=4 &\text{: } 13^4-1 = 28560 & \text{ divisible by } 5
\end{align*} 
Теперь мы знаем, что степень вложения \TJJ{} относительно простого множителя $5$ равна $k(5)=4$.

В реальных приложениях, как в случае кривых, дружественных к спариванию, таких как \curvename{BLS\_12-381}\sme{add reference}, обычно имеет значение только степень вложения большого простого множителя. В случае нашей кривой \curvename{Tiny-jubjub} это представлено $5$. Следует отметить, однако, что каждый простой множитель порядка кривой имеет свою собственную степень вложения, несмотря на то что это в приложениях в основном не имеет значения.

Чтобы найти степень вложения малого простого множителя $2$, мы должны найти наименьшее целое $k$ такое, что $2$ делит $q^k-1= 13^k-1$. Мы пробуем и увеличиваем, пока не найдём подходящее $k$. 
\begin{align*}
k=1 &\text{: } 13^1-1 = 12 & \text{ divisible by } 2
\end{align*} 

Теперь мы знаем, что степень вложения \TJJ{} относительно простого множителя $2$ равна $1$. Как мы видели, различные простые множители в общем случае могут иметь разные степени вложения.

Мы проверяем наши вычисления с помощью Sage:

\begin{sagecommandline}
sage: p = ZZ(13)
sage: # большой простой множитель
sage: r = ZZ(5)
sage: k = ZZ(1)
sage: while k < r:  # малая теорема Ферма
....:     if (p^k-1)%r == 0:
....:         break
....:     k=k+1
sage: k
sage: # малый простой множитель
sage: r = ZZ(2)
sage: k = ZZ(1)
sage: while k < r:  # малая теорема Ферма
....:     if (p^k-1)%r == 0:
....:         break
....:     k=k+1
sage: k
\end{sagecommandline}
\end{example}

\begin{example}
\label{ex:pairings_on_secp256k1}
 Чтобы привести пример криптографически стойкой реальной эллиптической кривой, не имеющей малой степени вложения, снова рассмотрим кривую \curvename{secp256k1}. Мы знаем из \examplename{} \ref{secp256k1}, что порядок этой кривой является простым числом, что означает, что у нас есть только одна степень вложения.

Чтобы проверить возможные степени вложения $k$, скажем, в диапазоне $1\leq k < 1000$, мы можем использовать Sage и вычислить следующим образом:
\begin{sagecommandline}
sage: p = ZZ(115792089237316195423570985008687907853269984665640564039457584007908834671663)
sage: r = ZZ(115792089237316195423570985008687907852837564279074904382605163141518161494337)
sage: k = ZZ(1)
sage: while k < 1000:
....:     if (p^k-1)%r == 0:
....:         break
....:     k=k+1
sage: k
\end{sagecommandline}
Мы видим, что \curvename{secp256k1} не имеет степени вложения $k<1000$, что означает, что \curvename{secp256k1} — это кривая, не имеющая малой степени вложения. Это свойство будет важным позже.\sme{add reference}
\end{example}
\begin{example}
\label{ex:embedding_degre_BN128}
Чтобы привести пример криптографически стойкой реальной эллиптической кривой, имеющей малую степень вложения, снова рассмотрим кривую \curvename{alt\_bn128}. Мы знаем из \examplename{} \ref{BN128}, что порядок этой кривой является простым числом, что означает, что у нас есть только одна степень вложения.

Чтобы вычислить степени вложения $k$, мы можем использовать Sage и перебирать малые степени вложения, пока не найдём совпадение. Мы вычисляем следующим образом:
\begin{sagecommandline}
sage: p = ZZ(21888242871839275222246405745257275088696311157297823662689037894645226208583)
sage: r = ZZ(21888242871839275222246405745257275088548364400416034343698204186575808495617)
sage: k = ZZ(1)
sage: # степень предполагается малой
sage: while k < 50: 
....:     if (p^k-1)%r == 0:
....:         break
....:     k=k+1
sage: k
\end{sagecommandline}
\end{example}



\subsubsection{Эллиптические кривые над расширениями полей}
\label{sec:curve-extensions}
\sme{TODO:Rewrite intro together with Sven} Предположим, что $p$ — простое число, и $\F_p$ — соответствующее простое поле. Мы знаем из уравнения \eqref{eq:prime-extension-field}, что поля $\F_{p^m}$ являются расширениями $\F_p$ в том смысле, что $\F_{p}$ — подполе $\F_{p^m}$. Это означает, что мы можем расширить аффинную плоскость, на которой определена эллиптическая кривая, заменив базовое поле на любое расширенное поле. Более точно, пусть 
$E(\F)=\{ (x,y)\in \F \times \F \;|\; y^2 = x^3 +a\cdot x +b \}$ — аффинная \concept{кривая в короткой форме Вейерштрасса} с параметрами $a$ и $b$ из $\F$. Если $\F'$ — расширенное поле $\F$, то мы расширяем область определения кривой, определяя $E(\F')$ следующим образом:

\begin{equation}\label{elliptic-curve-extension}
E(\F')=\{ (x,y)\in \F' \times \F' \;|\; y^2 = x^3 +a\cdot x +b \}
\end{equation}   

Хотя мы не изменили определяющие параметры, теперь мы рассматриваем точки кривой из аффинной плоскости над расширенным полем. Поскольку $\F\subset \F'$, можно показать, что исходная эллиптическая кривая $E(\F)$ является подкривой расширенной кривой $E(\F')$.

\begin{example}\label{ex:EF52} Рассмотрим простое поле $\F_5$ из \examplename{} \ref{prime-field-F5} вместе с эллиптической кривой $E_{1,1}(\F_5)$ и её определением из \examplename{} \ref{E1F5}, а также построение расширенного поля $\F_{5^2}$ относительно многочлена $t^2+2 \in \F_5[t]$ из упражнения \ref{exercise:finite_fieldF5_2}. В этом примере мы расширяем определение $E_{1,1}(\F_5)$ до эллиптической кривой над $\F_{5^2}$ и вычисляем её множество точек:
$$
E_{1,1}(\F_{5^2}) = \{ (x,y)\in \F_{5^2}\times \F_{5^2} \;|\; y^2 = x^3 + x +1 \}
$$
Поскольку $\F_{5^2}$ содержит $25$ точек, чтобы вычислить множество $E_{1,1}(\F_{5^2})$, мы должны проверить $25\cdot 25 = 625$ пар, что, вероятно, немного утомительно. Вместо этого мы используем Sage для вычисления кривой. Для этого мы выбираем представление $\F_{5^2}$ из \ref{exercise:finite_fieldF5_2}. Мы получаем:
\begin{sagecommandline}
sage: F5= GF(5)
sage: F5t.<t> = F5[] 
sage: P_MOD_2 = F5t(t^2+2)
sage: P_MOD_2.is_irreducible()
sage: F5_2.<t> = GF(5^2, name='t', modulus=P_MOD_2)
sage: E1F5_2 = EllipticCurve(F5_2,[1,1])
sage: E1F5_2.order()
\end{sagecommandline}
Кривая $E_{1,1}(\F_{5^2})$ состоит из $27$ точек, в отличие от кривой $E_{1,1}(\F_{5})$, которая состоит из $9$ точек. Запись этих точек даёт следующее:
\begin{multline*}
E_{1,1}(\F_{5^2}) = \{\Oinf, (0, 4), (0, 1), (3, 4), (3, 1), (4, 3), (4, 2), (2, 4), (2, 1),\\ 
(4t + 3, 3t + 4), (4t + 3, 2t + 1),  (3t + 2, t), (3t + 2, 4t),\\ 
(2t + 2, t), (2t + 2, 4t), (2t + 1, 4t + 4), (2t + 1, t + 1),\\ 
(2t + 3, 3), (2t + 3, 2), (t + 3, 2t + 4), (t + 3, 3t + 1),\\ 
(3t + 1, t + 4), (3t + 1, 4t + 1), (3t + 3, 3), (3t + 3, 2), (1, 4t),  (1, t)
\}
\end{multline*}
Как мы видим, множество точек эллиптической кривой $E_{1,1}(\F_5)$ является подмножеством множества точек эллиптической кривой $E(\F_{5^2})$. Это ожидалось, поскольку простое поле $\F_5$ является подполем конечного поля $\F_{5^2}$.
\end{example}
\begin{exercise} Рассмотрите \concept{эллиптическую кривую в короткой форме Вейерштрасса} $E(\F_{5^2})$ из \examplename{} \ref{ex:EF52}, вычислите выражение $(4t+3,2t+1)\oplus(3t+3,2)$ вручную и перепроверьте вычисление с помощью Sage. Затем решите уравнение $x\oplus (3t+3,3)=(3,4)$ для некоторого $x\in E(\F_{5^2})$. После этого вычислите скалярное умножение $[5](2t + 1, 4t + 4)$, используя алгоритм удвоения-сложения из упражнения \ref{alg_double-and-add}.
\end{exercise}
\begin{exercise}
\label{exercise:TJJ134}
 Рассмотрите кривую \curvename{Tiny-jubjub} из \examplename{} \ref{TJJ13}. Покажите, что многочлен $t^4+2\in \F_{13}[t]$ неприводим. Затем напишите программу на Sage для реализации расширения конечного поля $\F_{13^4}$, реализуйте расширение кривой $TJJ\_13(\F_{13^4})$ и вычислите число точек кривой.
\end{exercise}
\begin{exercise}
\label{exercise:BN128-extension}
Рассмотрите кривую \curvename{alt\_bn128} и ассоциированное с ней базовое поле $\F_p$ из \examplename{} \ref{BN128}. Как мы знаем из примера \ref{ex:embedding_degre_BN128}, эта кривая имеет степень вложения $12$. Используйте Sage, чтобы найти неприводимый многочлен $P\in \F_p[t]$, и напишите программу на Sage для реализации расширения конечного поля $\F_{p^{12}}$ и для реализации расширения кривой $alt\_bn128(\F_{p^{12}})$, а также вычислите число точек кривой.
\end{exercise}
\subsection{Полные группы кручения}
\label{sec:full-torsion} Как мы увидим далее, криптографически интересные спаривания определены на так называемых подгруппах кручения эллиптических кривых. Чтобы определить \term{группы кручения (torsion groups)} эллиптической кривой, пусть $\F$ — конечное поле, $E(\F)$ — эллиптическая кривая порядка $n$ и $r$ — множитель $n$. Тогда \term{$r$-группа кручения ($r$-torsion group)} эллиптической кривой $E(\F)$ определяется как множество
\begin{equation}
\label{def:torsion_group}
E(\F)[r]:= \{P\in E(\F)\;|\; [r]P = \mathcal{O}\}
\end{equation} 
Основная теорема о конечных циклических группах \ref{def:fundamental_theorem_groups} утверждает, что каждый множитель $r$ порядка циклической группы однозначно определяет подгруппу размера этого множителя, и эти подгруппы являются важными примерами $r$-групп кручения. Мы видели примеры таких подгрупп в \ref{ex:G1G2-subgroups} и \ref{ex:TJJ13-cofactor-clearing}.

Когда мы рассматриваем расширения эллиптических кривых, определённые в \ref{elliptic-curve-extension}, мы можем спросить, что происходит с $r$-группами кручения в расширении. Может интуитивно казаться, что их расширение просто параллельно расширению кривой. Например, когда $E(\F_p)$ — кривая над простым полем $\F_p$ с некоторой $r$-группой кручения $E(\F_p)[r]$, и когда мы расширяем кривую до $E(\F_{p^m})$, тогда может существовать большая $r$-группа кручения $E(\F_{p^m})[r]$ такая, что $E(\F_p)[r]$ является подгруппой $E(\F_{p^m})[r]$. Это может казаться интуитивно понятным, поскольку $E(\F_p)$ — подмножество $E(\F_{p^m})$. 

Однако фактическая ситуация немного более удивительна. Чтобы увидеть это, пусть $\F_p$ — простое поле и пусть $E(\F_p)$ — эллиптическая кривая порядка $n$ такая, что $r$ — множитель $n$, со степенью вложения $k(r)$ и $r$-группой кручения $E(\F_p)[r]$. Тогда $r$-группа кручения $E(\F_{p^m})[r]$ расширения кривой равна $E(\F_p)[r]$ только до тех пор, пока степень $m$ меньше степени вложения $k(r)$ кривой $E(\F_p)$. 

Для простой степени $p^{k(r)}$ $r$-группа кручения $E(\F_{p^{k(r)}})[r]$ может оказаться больше, чем $E(\F_p)[r]$, и она содержит $E(\F_p)[r]$ как подгруппу. Мы называем её \term{полной $r$-группой кручения (full $r$-torsion group)} этой эллиптической кривой и записываем следующим образом
\begin{equation}
\label{def:full_torsion_group}
E[r] := E(\F_{p^{k(r)}})[r]
\end{equation}
$r$-группы кручения $E(\F_{p^m})[r]$ любых расширений кривой при $m>k(r)$ все равны $E[r]$. В этом смысле $E[r]$ уже является наибольшей $r$-группой кручения, что оправдывает название. Полная $r$-группа кручения содержит $r^2$ элементов и состоит из $r+1$ подгрупп, одна из которых — $E(\F_{p})[r]$. Следующая диаграмма суммирует ситуацию:
\begin{equation}
\label{def:full_torsion_group_tower}
\begin{array}{lccclclclcl}
E(\F_{p}) & \subset & \cdots  & \subset & E(\F_{p^{k(r)-1}}) & \subset & E(\F_{p^{k(r)}}) & \subset & E(\F_{p^{k(r)+1}}) & \subset & \ldots\\
E(\F_{p})[r] & = & \cdots  & = & E(\F_{p^{k(r)-1}})[r] & \subset & E(\F_{p^{k(r)}})[r] & = & 
E(\F_{p^{k(r)+1}})[r] & = & \ldots
\end{array}
\end{equation}

Итак, когда мы рассматриваем вложенные расширения эллиптических кривых, как в \ref{def:full_torsion_group_tower}, упорядоченные по простой степени $m$, тогда $r$-группа кручения остаётся постоянной для каждого уровня $m$, который меньше степени вложения $k(r)$, тогда как она внезапно расцветает в большую группу на уровне $k(r)$ с $r+1$ подгруппами, и затем все $r$-группы кручения на более высоких уровнях $m\geq k(r)$ остаются неизменными. Другими словами, как только расширенное поле достаточно велико, чтобы найти ещё одну точку кривой $P$ с $[r]P=\mathcal{O}$, которая не является элементом кривой над базовым полем, тогда мы фактически находим все точки в полной группе кручения.

\begin{example}
\label{example:E1_full_torsion}
 Снова рассмотрим кривую $E_{1,1}(\F_5)$. Мы знаем из \ref{ex:G1G2-subgroups}, что она содержит $3$-группу кручения и что степень вложения $3$ равна $k(3)=2$. Отсюда мы можем сделать вывод, что мы можем найти полную $3$-группу кручения $E_{1,1}[3]$ в расширении кривой $E_{1,1}(\F_{5^2})$, последнее из которых мы вычислили в \examplename{} \ref{ex:EF52}. 

Поскольку эта кривая мала, чтобы найти полную $3$-группу кручения, мы можем перебрать все элементы $E_{1,1}(\F_{5^2})$ и проверить определяющее уравнение $[3]P= \Oinf$. Обратившись к Sage и используя нашу реализацию $E_{1,1}(\F_{5^2})$ в Sage из \ref{ex:EF52}, мы вычисляем следующим образом:
\begin{sagecommandline}
sage: INF = E1F5_2(0) # Точка на бесконечности
sage: L_E1_3 = []
sage: for p in E1F5_2:
....:     if 3*p == INF:
....:         L_E1_3.append(p)
sage: E1_3 = Set(L_E1_3) # Полное множество 3-кручения
\end{sagecommandline}
$$
E_{1,1}[3] = \{\Oinf,(2,1),(2,4),(1,t),(1,4t), (2t + 1,t + 1), (2t + 1, 4t + 4), (3t + 1,t + 4), (3t + 1, 4t + 1) \}
$$
Как мы видим, группа $E_{1,1}[3]$ содержит $9=3^2$ элементов, и $3$-группа кручения $E_{1,1}(\F_5)[3]$ кривой над простым полем является подмножеством полной группы кручения. 
\end{example}
\begin{example}\label{ex:TJJ13-full-torsion} HHHH Рассмотрим кривую \curvename{Tiny-jubjub} из \examplename{} \ref{TJJ13}. Мы знаем из \examplename{} \ref{ex:TJJ13-embedding-degree}, что она содержит $5$-группу кручения и что степень вложения $5$ равна $4$. Это означает, что мы можем найти полную $5$-группу кручения $\mathit{TJJ\_13}[5]$ в расширении кривой $\mathit{TJJ\_13}(\F_{13^4})$. 

Чтобы вычислить полное кручение, сначала заметим, что, поскольку $\F_{13^4}$ содержит $28561$ элемент, вычисление $\mathit{TJJ\_13}(\F_{13^4})$ означает проверку $28561^2=815730721$ элементов. Для каждой из этих точек кривой $P$ мы затем должны проверить уравнение $[5]P=\Oinf$. Делать это для $815730721$ немного слишком медленно даже на компьютере.

К счастью, Sage имеет функцию, которая вычисляет все точки $P$ такие, что $[m]P=Q$ для данного целого $m$ и точки кривой $Q$. Используя расширение кривой из упражнения \ref{exercise:TJJ134}, следующий код на Sage предоставляет способ вычислить полную группу кручения:
\begin{sagecommandline}
sage: # определить расширенное поле
sage: F13= GF(13) # простое поле
sage: F13t.<t> = F13[] # многочлены от t
sage: P_MOD_4 = F13t(t^4+2) # неприводимый многочлен степени 4
sage: P_MOD_4.is_irreducible()
sage: F13_4.<t> = GF(13^4, name='t', modulus=P_MOD_4)
sage: TJJF13_4 = EllipticCurve(F13_4,[8,8]) # расширение TJJ
sage: # вычислить полное 5-кручение
sage: INF = TJJF13_4(0) # точка на бесконечности
sage: L_TJJF13_4_5 = INF.division_points(5) # [5]P == INF
sage: TJJF13_4_5 = Set(L_TJJF13_4_5)
sage: TJJF13_4_5.cardinality()	# число элементов
\end{sagecommandline}
Как и ожидалось, мы получаем группу, содержащую $5^2=25$ элементов. Чтобы увидеть, что степень вложения $4$ действительно является наименьшей простой степенью для нахождения полной $5$-группы кручения, давайте вычислим $5$-группу кручения кривой \curvename{Tiny-jubjub} над расширенным полем $\F_{13^3}$. Мы получаем следующее:
\begin{sagecommandline}
sage: # определить расширенное поле
sage: P_MOD_3 = F13t(t^3+2) # неприводимый многочлен степени 3
sage: P_MOD_3.is_irreducible()
sage: F13_3.<t> = GF(13^3, name='t', modulus=P_MOD_3)
sage: TJJF13_3 = EllipticCurve(F13_3,[8,8]) # расширение TJJ
sage: # вычислить 5-кручение
sage: INF = TJJF13_3(0)
sage: L_TJJF13_3_5 = INF.division_points(5) # [5]P == INF
sage: TJJF13_3_5 = Set(L_TJJF13_3_5) # $5$-кручение
sage: TJJF13_3_5.cardinality()	# число элементов
\end{sagecommandline}

Как мы видим, $5$-группа кручения \curvename{Tiny-jubjub} над $\F_{13^3}$ равна $5$-группе кручения \curvename{Tiny-jubjub} над самим $\F_{13}$. 
\end{example}

\begin{example} 
\label{example:secp256k1}
% https://www.sikoba.com/docs/SKOR_SV_Pairing_Based_Crypto.pdf
Рассмотрим кривую \curvename{secp256k1}. Мы знаем из \examplename{} \ref{secp256k1}, что кривая имеет некоторый простой порядок $r$. Из-за этого единственной группой кручения, которую нужно рассмотреть, является сама кривая, то есть группа кривой — это $r$-кручение. 

Чтобы найти полное $r$-кручение \curvename{secp256k1}, нам нужно вычислить степень вложения $k$. И, как мы видели в \ref{ex:pairings_on_secp256k1}, она по крайней мере не мала. Однако мы знаем из малой теоремы Ферма \ref{fermats-little-theorem}, что конечная степень вложения должна существовать. Можно показать, что она задаётся следующим 256-битным числом:
$$
k = \scriptstyle 192986815395526992372618308347813175472927379845817397100860523586360249056 
$$
Это означает, что степень вложения очень велика, что влечёт, что расширение поля $\F_{p^k}$ также очень велико. Чтобы понять, насколько велико $\F_{p^k}$, вспомним, что элемент $\F_{p^m}$ может быть представлен как строка $<x_0,\ldots,x_m>$ из $m$ элементов, каждый из которых содержит число из простого поля $\F_p$. Теперь, в случае \curvename{secp256k1}, такое представление имеет $k$ записей, каждая размером $256$ бит. Таким образом, без каких-либо оптимизаций представление такого элемента потребовало бы $k\cdot 256$ бит. Следовательно, не только невозможно вычислить полную $r$-группу кручения \curvename{secp256k1}, но более того, вообще невозможно даже записать отдельные элементы этой группы. 
\end{example}
\begin{exercise} Рассмотрите полную $5$-группу кручения $TJJ\_13[5]$ из \examplename{} \ref{ex:TJJ13-full-torsion}. Запишите множество всех элементов этой группы и определите подмножество всех элементов из $TJJ\_13(\F_{13})[5]$, а также $TJJ\_13(\F_{13^2})[5]$. Затем вычислите $5$-группу кручения $TJJ\_13(\F_{13^{8}})[5]$.
\end{exercise}
\begin{exercise} Рассмотрите кривую \curvename{secp256k1} из \examplename{} \ref{secp256k1} и её полную $r$-группу кручения, введённую в \examplename{} \ref{example:secp256k1}. Запишите один элемент из полной группы кручения кривой, не являющийся точкой на бесконечности.
\end{exercise}
\begin{exercise}
\label{exercise:BN128-full-torsion}
 Рассмотрите кривую \curvename{alt\_bn128} из \examplename{} \ref{BN128} и её расширение кривой из упражнения \ref{exercise:BN128-extension}. Напишите программу на Sage, которая вычисляет образующий из полной группы кручения кривой.
\end{exercise}

\subsection{Группы спаривания}
\label{sec:pairing_groups}
 Как мы утверждали выше, любая полная $r$-группа кручения содержит $r+1$ циклических подгрупп, две из которых представляют особый интерес в криптографии на эллиптических кривых на основе спариваний. Чтобы охарактеризовать эти группы, нам нужно рассмотреть так называемый \term{эндоморфизм Фробениуса (Frobenius endomorphism)} эллиптической кривой $E(\F)$ над некоторым конечным полем $\F$ характеристики $p$:
\begin{equation}\label{eq:frobenius-enomorphism}
\pi : E(\F) \to E(\F): \;\; 
\begin{array}{lcl}
(x,y)       &\mapsto & (x^p,y^p)\\
\Oinf &\mapsto & \Oinf
\end{array} 
\end{equation}
Можно показать, что $\pi$ отображает точки кривой в точки кривой. Первое, на что следует обратить внимание, — это то, что в случае, когда $\F$ — простое поле, эндоморфизм Фробениуса действует как тождественное отображение, поскольку $(x^p,y^p) = (x,y)$ на простых полях благодаря малой теореме Ферма \ref{fermats-little-theorem}. Это означает, что отображение Фробениуса более интересно на эллиптических кривых над расширениями простых полей.

Имея в распоряжении отображение Фробениуса, мы можем охарактеризовать две важные подгруппы полной $r$-группы кручения $E[r]$ эллиптической кривой. Первая подгруппа — это группа элементов из полной $r$-группы кручения, на которых отображение Фробениуса действует тривиально. Поскольку в криптографии на основе спариваний эта группа обычно записывается как $\G_1$, предполагая, что простой множитель $r$ в определении неявно задан, мы определяем $\G_1$ следующим образом:

\begin{equation}
\label{def:pairing_group_G1}
\G_1[r] := \{(x,y)\in E[r]\;|\; \pi(x,y) = (x,y)\;\}
\end{equation}

Можно показать, что $\G_1$ в точности является $r$-группой кручения $E(\F_p)[r]$ нерасширенной эллиптической кривой, определённой над простым полем. Существует ещё одна подгруппа полной $r$-группы кручения, которую можно охарактеризовать с помощью отображения Фробениуса, и в контексте криптографии на основе спариваний эта подгруппа часто называется $\G_2$. Эта группа определяется следующим образом:

\begin{equation}
\label{def:pairing_group_G2}
\G_2[r]:= \{(x,y)\in E[r]\;|\; \pi(x,y) = [p](x,y)\;\}
\end{equation}

\begin{notation} Если $E(\F)$ — эллиптическая кривая и $r$ — наибольший простой множитель порядка кривой, мы называем $\G_1[r]$ и $\G_2[r]$ \term{группами спаривания (pairing groups)}. Если простой множитель $r$ ясен из контекста, мы иногда просто пишем $\G_1$ и $\G_2$, имея в виду $\G_1[r]$ и $\G_2[r]$ соответственно.
\end{notation}

Следует отметить, что другие определения $\G_2$ также существуют в литературе. Однако в контексте криптографии на основе спариваний это распространённый выбор, поскольку он особенно полезен, потому что мы можем определить эффективные хеш-функции, отображающие в $\G_2$, что невозможно для всех подгрупп полной $r$-группы кручения.

\begin{example} Снова рассмотрим кривую $E_{1,1}(\F_5)$ из \examplename{} \ref{E1F5}. Как мы видели, эта кривая имеет степень вложения $k=2$, и полная $3$-группа кручения задаётся следующим образом:
\begin{multline}
E_{1,1}[3] = \{\Oinf,(2,1),(2,4), (1,t), (1,4t), (2t + 1,t + 1),\\ (2t + 1, 4t + 4),
(3t + 1,t + 4), (3t + 1, 4t + 1) \}
\end{multline}

Согласно общей теории, $E_{1,1}[3]$ содержит $4$ подгруппы, и мы можем охарактеризовать подгруппы $\G_1$ и $\G_2$ с помощью эндоморфизма Фробениуса. К сожалению, на момент написания Sage не имеет предопределённого эндоморфизма Фробениуса для эллиптических кривых, поэтому нам приходится использовать эндоморфизм Фробениуса базового поля как временное решение. Используя нашу реализацию $E_{1,1}[3]$ в Sage из \examplename{} \ref{example:E1_full_torsion}, мы вычисляем $\G_1$ следующим образом:
\begin{sagecommandline}
sage: L_G1 = []
sage: for P in E1_3: 
....:     PiP = E1F5_2([a.frobenius() for a in P]) # pi(P)
....:     if P == PiP:
....:         L_G1.append(P)
sage: G1 = Set(L_G1)
\end{sagecommandline}
Как и ожидалось, группа $\G_1=\{\Oinf, (2,4), (2,1)\}$ идентична $3$-группе кручения (нерасширенной) кривой над простым полем $E_{1,1}(\F_5)$. 

Чтобы вычислить группу $\G_2$ для кривой $E_{1,1}(\F_5)$, мы можем использовать почти тот же алгоритм, который использовали для вычисления $\G_1$. Поскольку $p=5$, мы получаем следующее:
\begin{sagecommandline}
sage: L_G2 = []
sage: for P in E1_3: 
....:     PiP = E1F5_2([a.frobenius() for a in P]) # pi(P)
....:     pP = 5*P # [5]P
....:     if pP == PiP:
....:         L_G2.append(P)
sage: G2 = Set(L_G2)
\end{sagecommandline}

Таким образом, мы вычислили группу спаривания $\G_2$ полной $3$-группы кручения кривой $E_{1,1}(\F_5)$ как множество $\G_2 = \{\Oinf, (1,t), (1,4t)\}$. 
\end{example}

\begin{example}
\label{example:TJJ_pairing_groups}
Рассмотрим кривую \curvename{Tiny-jubjub} \TJJ{} из \examplename{} \ref{TJJ13}. В \examplename{} \ref{ex:TJJ13-full-torsion} мы вычислили её полное $5$-кручение, которое является группой, имеющей $6$ подгрупп. Мы вычисляем $\G_1$ с помощью Sage следующим образом:
\begin{sagecommandline}
sage: L_TJJ_G1 = []
sage: for P in TJJF13_4_5: 
....:     PiP = TJJF13_4([a.frobenius() for a in P]) # pi(P)
....:     if P == PiP:
....:         L_TJJ_G1.append(P)
sage: TJJ_G1 = Set(L_TJJ_G1)
\end{sagecommandline}
Мы получаем $\G_1= \{\Oinf, (7,2), (8,8), (8,5), (7,11)\}$ и, как и ожидалось, $\G_1$ идентична $5$-группе кручения (нерасширенной) кривой над простым полем $TJJ\_13$, вычисленной в \examplename{} \ref{eq:TJJ13-logarithmic-order}.

Чтобы вычислить группу $\G_2$ для кривой Tiny-jubjub, мы можем использовать почти тот же алгоритм, который использовали для вычисления $\G_1$. Поскольку $p=13$, мы получаем следующее:
\begin{sagecommandline}
sage: L_TJJ_G2 = []
sage: for P in TJJF13_4_5: 
....:     PiP = TJJF13_4([a.frobenius() for a in P]) # pi(P)
....:     pP = 13*P # [13]P
....:     if pP == PiP:	# pi(P) ==[13]P
....:         L_TJJ_G2.append(P)
sage: TJJ_G2 = Set(L_TJJ_G2)
\end{sagecommandline}
$\G_2 = \{\Oinf, (9t^2 + 7,t^3 + 11t), (9t^2 + 7, 12t^3 + 2t), (4t^2 + 7,5t^3 + 10t),(4t^2 + 7,8t^3 + 3t)\}$
\end{example}

\begin{example}
\label{example:secp256k1_pairing_groups}
Снова рассмотрим кривую Bitcoin \curvename{secp256k1}. Поскольку группа $\G_1$ идентична группе кручения нерасширенной кривой, и поскольку \curvename{secp256k1} имеет простой порядок, мы знаем, что в этом случае $\G_1$ идентична самой \curvename{secp256k1}. Однако невозможно вычислить элементы из $\G_2$, поскольку, согласно \examplename{} \ref{example:secp256k1}, мы не можем хранить типичные точки кривой из расширенной кривой $secp256k1(\F_{p^k})$ на каком-либо компьютере, не говоря уже о вычислении их образов под отображением Фробениуса.
\end{example}
\begin{exercise}
Рассмотрите малый простой множитель $2$ кривой \curvename{Tiny-jubjub}. Вычислите полную $2$-группу кручения $TJJ\_13$, а затем вычислите группы $\G_1[2]$ и $\G_2[2]$.
\end{exercise}
\begin{exercise}
\label{exercise:BN128-pairing-groups}
 Рассмотрите кривую \curvename{alt\_bn128} из \examplename{} \ref{BN128} и её расширение кривой из упражнения \ref{exercise:BN128-extension}. Напишите программу на Sage, которая вычисляет образующий для каждой из групп кручения $\G_1[p]$ и $\G_2[p]$.
\end{exercise}


\subsection{Спаривание Вейля}
\label{sec:weil-pairing} Вспомним определение невырожденного группового спаривания из \ref{pairing-map}. В этой части мы рассматриваем функцию спаривания, определённую на
подгруппах $\G_1[r]$ и $\G_2[r]$ полного $r$-кручения $E[r]$ \concept{эллиптической кривой в короткой форме Вейерштрасса}. Более точно, пусть $E(\F_p)$ — эллиптическая кривая степени вложения $k$ такая, что $r$ — простой множитель её порядка. Тогда \term{спаривание Вейля (Weil pairing)} определяется как следующее билинейное невырожденное отображение:

\begin{equation}\label{eq:weil-pairing}
e(\cdot,\cdot) : \G_1[r] \times \G_2[r] \to \F^*_{p^k}\; ;\; 
(P,Q)\mapsto (-1)^r \cdot \frac{f_{r,P}(Q)}{f_{r,Q}(P)}
\end{equation} 

Элементы расширенного поля $f_{r,P}(Q), f_{r,Q}(P)\in \F_{p^k}$ в определении спаривания Вейля вычисляются \term{алгоритмом Миллера (Miller's algorithm)}, приведённым ниже.

\begin{algorithm}\caption{Алгоритм Миллера для \concept{кривых в короткой форме Вейерштрасса} $y^2 = x^3 +ax +b$}\label{alg:millersalgo}
% https://www.math.u-bordeaux.fr/~damienrobert/csi2018/pairings.pdf
\begin{algorithmic}[0]
\Require $r>3$, $P \in E[r]$, $Q\in E[r]$ and
\State $b_0,\ldots, b_t\in \{0,1\}$ with $r= b_0\cdot 2^0 + b_1\cdot 2^1 + \ldots + b_t\cdot 2^t$ and $b_t=1$
\Procedure{Miller's Algorithm}{$P,Q$}
\If{$P = \Oinf$ or $Q = \Oinf$ or $P = Q$}
	\State \textbf{return} $f_{r,P}(Q) \gets (-1)^r$
\EndIf
\State $(x_T,y_T) \gets (x_P,y_P)$
\State $f_1\gets 1$
\State $f_2\gets 1$
\For{$j\gets t-1,\ldots, 0$}
	\State $m \gets \frac{3\cdot x_T^2+a}{2\cdot y_T}$	
    \State $f_1 \gets f_1^2\cdot (y_Q - y_T - m\cdot(x_Q-x_T))$
	\State $f_2 \gets f_2^2\cdot (x_Q + 2x_T -m^2)$
	\State $x_{2T} \gets m^2 - 2 x_T$
	\State $y_{2T} \gets -y_T - m\cdot (x_{2T}-x_T)$
	\State $(x_T,y_T)\gets (x_{2T},y_{2T})$ 
	\If{$b_j = 1$}
		\State $m \gets \frac{y_T -y_P}{x_T - x_P}$
		\State $f_1 \gets f_1\cdot (y_Q -y_T -m\cdot (x_Q - x_T))$
		\State $f_2 \gets f_2\cdot (x_Q + (x_P+x_T) - m^2)$
		\State $x_{T+P} \gets m^2 -x_T -x_P$
		\State $y_{T+P}\gets -y_T - m\cdot (x_{T+P}-x_T)$
		\State $(x_T,y_T)\gets (x_{T+P},y_{T+P})$
	\EndIf
\EndFor
\State $f_1 \gets f_1\cdot (x_Q - x_T)$
\State \textbf{return} $f_{r,P}(Q) \gets \frac{f_1}{f_2}$
\EndProcedure
\end{algorithmic}
\end{algorithm}

Понимание деталей того, как и почему работает этот алгоритм, требует понятия \term{дивизоров (divisors)}, которое выходит за рамки этой книги. Заинтересованный читатель может обратиться к \chaptname{} 6, \secname{} 6.8.3 в \cite{hoffstein-2008}, или к \chaptname {} 3 в \cite{costello-pairings}. Как мы видим, алгоритм более эффективен для простых чисел $r$, имеющих малую хэммингову вес \ref{def:binary_representation_integer}.

Мы называем эллиптическую кривую $E(\F_p)$ \term{дружественной к спариванию (pairing-friendly)}, если существует простой множитель порядка группы такой, что спаривание Вейля эффективно вычислимо относительно этого простого множителя. В реальных приложениях кривых, дружественных к спариванию, степень вложения обычно является малым числом, таким как $2$, $4$, $6$ или $12$, а число $r$ — наибольший простой множитель порядка кривой. 

\begin{example}Рассмотрим кривую $E_{1,1}(\F_5)$ из \examplename{} \ref{E1F5}. Поскольку единственный простой множитель порядка группы — $3$, мы не можем вычислить спаривание Вейля в этой группе, используя наше определение алгоритма Миллера. На самом деле, поскольку $\G_1$ имеет порядок $3$, выполнение алгоритма приведёт к ``делению на нуль''.
\end{example}

\begin{example} Рассмотрим кривую \curvename{Tiny-jubjub} $\mathit{TJJ\_13}(\F_{13})$ из \examplename{} \ref{TJJ13} и ассоциированные с ней группы спаривания из \examplename{} \ref{example:TJJ_pairing_groups}:
\begin{align*}
\G_1[5] & = \{\Oinf, (7,2), (8,8), (8,5), (7,11)\}\\
\G_2[5] & = \{\Oinf, (9t^2 + 7,t^3 + 11t), (9t^2 + 7, 12t^3 + 2t), 
(4t^2 + 7, 5t^3 + 10t), (4t^2 + 7, 8t^33 + 3t)\}
\end{align*}

Поскольку мы знаем из \examplename{} \ref{ex:TJJ13-embedding-degree}, что степень вложения $5$ равна $4$, мы можем конкретизировать общее определение спаривания Вейля для этого примера следующим образом:
$$
e(\cdot,\cdot): \G_1[5] \times \G_2[5] \to \F_{13^4}
$$ 

Первое условие if в алгоритме Миллера влечёт, что $e(\Oinf,Q)=1$, а также $e(P,\Oinf)=1$ для всех аргументов $P\in\G_1[5]$ и $Q\in \G_2[5]$. Чтобы вычислить нетривиальное спаривание Вейля, мы выбираем аргументы $P=(7,2)\in \G_1$ и $Q=(9t^2 + 7, 12t^3 + 2t)\in\G_2$. Обратившись к Sage, мы получаем следующее вычисление спаривания Вейля:  
\begin{sagecommandline}
sage: F13 = GF(13)
sage: F13t.<t> = F13[]
sage: P_MOD_4 = F13t(t^4+2)
sage: F13_4.<t> = GF(13^4, name='t', modulus=P_MOD_4)
sage: TJJF13_4 = EllipticCurve(F13_4,[8,8])
sage: P=TJJF13_4([7,2])
sage: Q=TJJF13_4([9*t^2+7,12*t^3+2*t])
sage: P.weil_pairing(Q,5)
\end{sagecommandline}
\end{example}
\begin{example}
Снова рассмотрим кривую Bitcoin \curvename{secp256k1}. Как мы видели в \examplename{} \ref{example:secp256k1_pairing_groups}, невозможно вычислить элементы из группы спаривания $\G_2$, и, как мы знаем из \examplename{} \ref{example:secp256k1}, более того, невозможно выполнять вычисления в расширенном поле $\F_{p^k}$. Следовательно, спаривание Вейля не является эффективно вычислимым, и \curvename{secp256k1} не является дружественной к спариванию.
\end{example}
\begin{exercise}
\label{exercise:BN128-pairing}
Рассмотрите кривую \curvename{alt\_bn128} из \examplename{} \ref{BN128} и образующие $g_1$ и $g_2$ групп $\G_1[p]$ и $\G_2[p]$ из упражнения \ref{exercise:BN128-pairing-groups}. Напишите программу на Sage, которая вычисляет спаривание Вейля $e(g_1,g_2)$.
\end{exercise}

% http://www.pdmi.ras.ru/~lowdimma/BSD/Silverman-Arithmetic_of_EC.pdf
% p. 396ff

\section{Хеширование на кривые} Криптография на эллиптических кривых часто требует возможности хешировать данные на эллиптические кривые. Если порядок кривой не является простым числом, хеширование в подгруппы простого порядка также важно, и в контексте кривых, дружественных к спариванию, иногда необходимо хешировать специально в группу спаривания $\G_1$ или $\G_2$, как введено в \ref{sec:pairing_groups}.

Как мы видели в разделе \ref{sec:hashing-to-groups}, известны некоторые общие методы хеширования в конечные циклические группы. Поскольку эллиптические кривые над конечными полями являются конечными коммутативными группами, эти методы могут быть использованы путём хеширования в выбранную циклическую подгруппу (часто с последующей очисткой кофактора). Однако в дальнейшем мы хотим описать некоторые методы, специфичные для эллиптических кривых, которые часто используются в реальных приложениях. 

\subsection{Хеш-функции методом проб и приращений}
Один из наиболее прямых способов безопасного хеширования на точку эллиптической кривой — использовать криптографическую хеш-функцию вместе с одним из методов хеширования в модульную арифметику, как описано в разделе \ref{hash-to-modular-arithmetics}.

Обе конструкции могут быть объединены таким образом, что образ предоставляет элемент базового поля эллиптической кривой вместе с одним вспомогательным битом. Элемент базового поля затем может быть интерпретирован как $x$-координата потенциальной точки кривой, а вспомогательный бит может быть использован для определения одной из двух возможных $y$-координат этой точки кривой, как объяснено в \ref{sec:affine_point_compression}.

Такой подход был бы детерминированным и лёгким в реализации, и он сохранял бы криптографические свойства исходной хеш-функции. Однако не все $x$-координаты, сгенерированные таким образом, дадут квадратичные вычеты при подстановке в определяющее уравнение. Следовательно, не все элементы поля порождают реальные точки кривой. 

In fact,
% https://www.cs.umd.edu/users/gasarch/TOPICS/res/burgess.pdf
на простом поле только половина элементов поля являются квадратичными вычетами. Следовательно, предполагая равномерное распределение хеш-значений в поле, этот метод не смог бы сгенерировать точку кривой примерно в половине попыток. 

Один из способов учесть эту проблему — следующий так называемый \term{метод проб и приращений (try-and-increment)}. Вместо простого хеширования двоичной строки $s$ в поле этот метод использует хеширование методом проб и приращений в базовое поле, как описано в \ref{def:try_and_increment_hash}, в сочетании с одним вспомогательным битом, полученным из базовой криптографической хеш-функции.

Если любая попытка хеширования в поле не даёт элемента поля или действительной точки кривой, счётчик увеличивается, и хеширование повторяется. Это делается до тех пор, пока не будет найдена действительная точка кривой (см. алгоритм ниже).

\begin{algorithm}\label{alg:hash-to-e}\caption{Хеширование в $E(\F_p)$}
\begin{algorithmic}[0]
\Require $p \in \Z$ с $|p|=k$ и $s\in\{0,1\}^*$
\Require Уравнение кривой $y^2 = x^3 + ax +b$ над $\F_p$
\Procedure{Try-and-Increment}{$r,k,s$}
\State $c \gets 0$	\Comment{Счётчик проб и приращений}
\Repeat
\State $s' \gets s||Bits(c)$
\State $x \gets H(s')_0\cdot 2^0 + H(s')_1\cdot 2^1 + \ldots + H(s')_{k}\cdot 2^{k}$ \Comment{потенциальное $x$}
\State $y^2 \gets z^3 + a\cdot z + b$ \Comment{потенциальное $y^2$}
\State $c\gets c+1$
\Until{$x<p$ и $\Zmod{(y^2)^{\frac{p-1}{2}}}{r}=1$ } \Comment{Проверить $x$ в поле и $y^2$ имеет корень}
\If {$H(s')_{k+1} == 0$} \Comment{вспомогательный бит определяет корень}
\State $y \gets y'\in \sqrt{y^2}$ с $0\leq y' \leq (p-1)/2$
\Else 
\State $y \gets y'\in \sqrt{y^2}$ с $(p-1)/2 < y' < p$
\EndIf
\State \textbf{return} $(x,y)$
\EndProcedure
\Ensure $(x,y)\in E(\F_r)$
\end{algorithmic}
\end{algorithm}

Метод проб и приращений относительно легко реализовать, и он сохраняет криптографические свойства исходной хеш-функции. Следует отметить, что если кривая не имеет простого порядка, образ хеша методом проб и приращений будет общей точкой кривой, которая может не быть элементом большой подгруппы простого порядка. Чтобы отобразить в большую подгруппу простого порядка, поэтому необходимо применить технику очистки кофактора, как объяснено в \ref{def:cofactor_clearing}.

\begin{example} Рассмотрим кривую \curvename{Tiny-jubjub} из \examplename{} \ref{TJJ13}. Мы хотим построить хеш-функцию методом проб и приращений, которая отображает двоичную строку $s$ произвольной длины в большую подгруппу простого порядка размера $5$ из \examplename{} \ref{eq:TJJ13-logarithmic-order}. 

Поскольку кривая $TJJ\_13$ определена над полем $\F_{13}$, а двоичное представление $13$ равно $Bits(13)=<1,1,0,1>$, один из способов реализовать функцию проб и приращений — применить SHA256 из библиотеки hashlib Sage к конкатенации $s||c$ для некоторой двоичной строки-счётчика $c$, и использовать первые $4$ бита образа, чтобы попытаться хешировать в $\F_{13}$. В случае, если мы можем хешировать в значение $x$ такое, что $x^3 +8\cdot x + 8$ является квадратичным вычетом в $\F_{13}$, мы используем пятый бит, чтобы решить, какой из двух возможных корней $x^3 + 8\cdot x + 8$ мы выберем в качестве $y$-координаты. Результат — точка кривой, отличная от точки на бесконечности. Чтобы спроецировать её в большую подгруппу простого порядка $TJJ\_13[5]$, мы умножаем её на кофактор $4$. Если результат не является точкой на бесконечности, это результат хеша.

Чтобы конкретизировать это, пусть $s=<1,1,1,0,0,1,0,0,0,0>$ — наша двоичная строка, которую мы хотим хешировать в $TJJ\_13[5]$. Мы используем двоичную строку-счётчик, начиная с нуля, то есть выбираем $c=<0>$. Обратившись к Sage, мы определяем функцию try-hash следующим образом:
\begin{sagecommandline}
sage: import hashlib
sage: def try_hash(s,c):
....:     s_1 = s+c # string concatenation
....:     hasher = hashlib.sha256(s_1.encode('utf-8')) # compute SHA256
....:     digest = hasher.hexdigest()
....:     z = ZZ(digest, 16) # cast into integer
....:     z_bin = z.digits(base=2, padto=256) # cast to 256 bits
....:     x = z_bin[0]*2^0 + z_bin[1]*2^1 + z_bin[2]*2^2+z_bin[3]*2^3
....:     return (x,z_bin[4])
sage: try_hash('1110010000','0')
\end{sagecommandline}

Как мы видим, наша первая попытка хешировать в $\F_{13}$ не была успешной, поскольку $15$ не является элементом $\F_{13}$, поэтому мы увеличиваем двоичный счётчик на $1$ и пробуем снова: 
\begin{sagecommandline}
sage: try_hash('1110010000','1')
\end{sagecommandline}

С этой попыткой мы нашли хеш в $\F_{13}$. Однако эта точка не гарантированно определяет точку кривой. Чтобы увидеть это, мы подставляем $x=3$ в правую часть уравнения \concept{короткой формы Вейерштрасса} кривой \curvename{Tiny-jubjub} и вычисляем $3^3 + 8\cdot 3 + 8 = 7$. Однако $7$ не является квадратичным вычетом в $\F_{13}$, поскольку $7^{\frac{13-1}{2}}=7^6=12=-1$. Это означает, что элемент поля $7$ не подходит в качестве $x$-координаты какой-либо точки кривой. Поэтому мы должны увеличить счётчик ещё раз: 
\begin{sagecommandline}
sage: try_hash('1110010000','10')
\end{sagecommandline}
Поскольку $12^3 + 8\cdot 12 + 8 = 12$, и мы имеем $\sqrt{12} = \{5, 8\}$, мы наконец нашли действительную $x$-координату $x=12$ для хеша точки кривой. Теперь, поскольку вспомогательный бит этого хеша равен $1$, мы выбираем больший корень $y=8$ в качестве $y$-координаты и получаем следующий хеш, который является действительной точкой кривой на \curvename{Tiny-jubjub}:
$$
H_{TJJ\_13}(<1,1,1,0,0,0,0,0>) = (12,8)
$$

Чтобы спроецировать это в ``большую'' подгруппу простого порядка, мы должны выполнить очистку кофактора, то есть мы должны умножить точку на кофактор $4$. Используя Sage, мы получаем
\begin{sagecommandline}
sage: P = TJJ_13(12,8)
sage: (4*P).xy()
\end{sagecommandline}

Это означает, что хеширование двоичной строки $<1,1,1,0,0,0,0,0>$ в большую подгруппу простого порядка $TJJ\_13[5]$ даёт в результате хеш-значение $(8,8)$. 
$$
H_{TJJ\_13[5]}(<1,1,1,0,0,0,0,0>) = (8,8)
$$
\end{example}
\begin{exercise}
Используйте наше определение алгоритма $try\_hash$, чтобы реализовать хеш-функцию $H_{TJJ\_13[5]} : \{0,1\}^*\to TJJ\_13(\F_{13})[5]$, которая отображает двоичные строки произвольной длины в $5$-группу кручения $TJJ\_13(\F_{13})$.
\end{exercise}
\begin{exercise}
Реализуйте криптографическую хеш-функцию $H_{secp256k1} : \{0,1\}^*\to secp256k1$, которая отображает двоичные строки произвольной длины на эллиптическую кривую \curvename{secp256k1}.
\end{exercise}
\section{Построение эллиптических кривых} Криптографически стойкие эллиптические кривые, такие как \curvename{secp256k1} \ref{secp256k1}, известны уже довольно давно. Однако, учитывая последние достижения в криптографии, часто необходимо проектировать и создавать эллиптические кривые с нуля, удовлетворяющие определённым очень специфическим свойствам. 

Например, в контексте разработки SNARK стало необходимо проектировать эллиптические кривые, которые могут быть эффективно реализованы внутри так называемой \uterm{алгебраической схемы}, чтобы обеспечить примитивы, такие как \uterm{схемы подписей} на эллиптических кривых, в доказательстве с нулевым разглашением. Такой кривой является кривая Baby-jubjub, определённая в \cite{whitehat-21}, и мы параллельно ввели кривую \curvename{Tiny-jubjub} из \examplename{} \ref{TJJ13}. Как мы видели, эти кривые являются примерами так называемых \concept{кривых в крученой форме Эдвардса}, и как таковые имеют легко реализуемые законы сложения, работающие без ветвления. Однако мы ввели кривую \curvename{Tiny-jubjub} из ниоткуда, просто дав параметры кривой, не объясняя, как мы их получили.

Ещё одно требование в контексте многих так называемых \term{систем доказательств с нулевым разглашением на основе спариваний (pairing-based zero-knowledge proofing systems)} — существование подходящей кривой, дружественной к спариванию, с заданным уровнем безопасности и малой степенью вложения, как определено в \ref{def:embedding-degree}. Известные примеры — кривые BLS\_12 и NMT.\sme{add references}

Основная цель этого раздела — объяснить наиболее важный метод проектирования эллиптических кривых с заранее заданными свойствами с нуля, называемый \term{\concept{методом комплексного умножения} (complex multiplication method)} (см. \chaptname{} 6 в \cite{silverman-1994}). Мы применим этот метод в разделе \ref{BLS6}, чтобы синтезировать конкретную кривую BLS6, которая является одной из наиболее нестойких кривых, но особенно хорошо подходит в качестве основной кривой для построения наших SNARK для работы вручную. Как мы увидим, эта кривая имеет ``большую'' подгруппу простого множителя порядка $13$, что означает, что мы можем использовать нашу кривую \curvename{Tiny-jubjub} для реализации определённых криптографических примитивов на эллиптических кривых в схемах над этой кривой BLS6. 
 
Прежде чем ввести \concept{метод комплексного умножения}, мы должны объяснить несколько свойств эллиптических кривых, которые имеют ключевое значение для понимания этого метода. 

\subsection{След Фробениуса} Чтобы понять \concept{метод комплексного умножения} эллиптических кривых, мы сначала должны определить так называемый \term{след (trace)} эллиптической кривой.

Мы знаем, что группы эллиптических кривых над конечными полями являются конечными коммутативными группами. Поэтому интересный вопрос заключается в том, возможно ли оценить число элементов, которые содержит эта кривая. Поскольку аффинная \concept{кривая в короткой форме Вейерштрасса} состоит из пар $(x,y)$ элементов из конечного поля $\F_q$ плюс точка на бесконечности, а поле $\F_q$ содержит $q$ элементов, число точек кривой не может быть произвольно большим, поскольку оно может содержать не более $q^2+1$ элементов. 

Однако существует более точная оценка, обычно называемая \term{границей Хассе (Hasse bound)}. Чтобы понять её, пусть $E(\F_q)$ — аффинная \concept{кривая в короткой форме Вейерштрасса} над конечным полем $\F_q$ порядка $q$, и пусть $|E(\F_q)|$ — порядок кривой. Тогда существует целое число $t\in \Z$, называемое \term{следом Фробениуса (trace of Frobenius)} кривой, такое что $|t| \leq 2\sqrt{q}$ и выполняется следующее уравнение:
\begin{equation}\label{hasse-bound}
|E(\F)| = q +1 -t
\end{equation}

Положительный след, следовательно, означает, что кривая содержит не больше точек, чем базовое поле, тогда как неотрицательный след означает, что кривая содержит больше точек. Однако оценка $|t| \leq 2\sqrt{q}$ означает, что разница не очень велика ни в одном направлении, и число элементов в эллиптической кривой всегда примерно того же порядка величины, что и размер базового поля кривой.

\begin{example}\label{ex:E1F5-frobenius} Рассмотрим эллиптическую кривую $E_{1,1}(\F_5)$ из \examplename{} \ref{E1F5}. Мы знаем, что она содержит $9$ точек кривой. Поскольку порядок $\F_5$ равен $5$, мы вычисляем след $E_{1,1}(\F_5)$ как $t=-3$, поскольку граница Хассе задаётся следующим уравнением:
$$
9 = 5 + 1 - (-3)
$$
Действительно, мы имеем $|t| \leq 2\sqrt{q}$, поскольку $\sqrt{5}> 2$ и 
$|-3|= 3 \leq 4 = 2\cdot 2< 2\cdot \sqrt{5}$.
\end{example}

\begin{example}\label{ex:TJJ13-frobenius} Чтобы вычислить след кривой \curvename{Tiny-jubjub}, вспомним из \examplename{} \ref{TJJ13}, что порядок \TJJ{} равен $20$. Поскольку порядок $\F_{13}$ равен $13$, мы можем использовать границу Хассе и вычислить след как $t=-6$:
\begin{equation}
20 = 13 + 1 - (-6)
\end{equation}

Снова мы имеем $|t| \leq 2\sqrt{q}$, поскольку $\sqrt{13}> 3$ и 
$|-6|= 6 = 2\cdot 3< 2\cdot \sqrt{13}$.
\end{example}

\begin{example}\label{ex:secp256k1-trace}Чтобы вычислить след \curvename{secp256k1}, вспомним из \examplename{} \ref{secp256k1}, что эта кривая определена над простым полем с $p$ элементами, и что порядок этой группы задаётся $r$:  
\begin{align*}
p &= \scriptstyle 115792089237316195423570985008687907853269984665640564039457584007908834671663\\
r &= \scriptstyle 115792089237316195423570985008687907852837564279074904382605163141518161494337
\end{align*}

Используя границу Хассе $r = p + 1 -t$, мы поэтому вычисляем $t= p+1 -r$, что даёт след кривой \curvename{secp256k1} следующим образом:
$$
t = \scriptstyle 432420386565659656852420866390673177327
$$

Как мы видим, \curvename{secp256k1} содержит меньше элементов, чем её базовое поле. Однако разница крошечна, поскольку порядок \curvename{secp256k1} того же порядка величины, что и порядок базового поля. По сравнению с $p$ и $r$, целое число $t$ крошечно.

\begin{sagecommandline}
sage: p = 115792089237316195423570985008687907853269984665640564039457584007908834671663
sage: r = 115792089237316195423570985008687907852837564279074904382605163141518161494337
sage: t = p + 1 -r
sage: t.nbits()
sage: abs(RR(t)) <= 2*sqrt(RR(p))
\end{sagecommandline}
\end{example} 
\begin{exercise}
\label{exercise:BN128-trace}
Рассмотрите кривую \curvename{alt\_bn128} из \examplename{} \ref{BN128}. Напишите программу на Sage, которая вычисляет след Фробениуса для \curvename{alt\_bn128}. Содержит ли кривая больше или меньше элементов, чем её базовое поле $\F_p$?
\end{exercise}

\subsection{$j$-инвариант}
\label{sec:j-inv} Как мы видели в \ref{sec:isomorphic_curves}, две эллиптические кривые $E_1(\F)$, определённые уравнением $y^2 = x^3 + ax +b$, и $E_2(\F)$, определённые уравнением $y^2 + a'x + b'$, строго изоморфны тогда и только тогда, когда существует квадратичный вычет $d\in \F$ такой, что $a' = a d^2$ и $b' = b d^3$. 

Однако существует более общий способ классифицировать эллиптические кривые над конечными полями $\F_q$, основанный на так называемом \term{$j$-инварианте ($j$-invariant)} эллиптической кривой с $j(E(\F_q))\in\F_q$, как определено ниже:
\begin{equation}\label{eq:j-invariant1}
j(E(\F_q)) = \Zmod{1728\cdot \frac{4\cdot a^3}{4\cdot a^3+ 27\cdot b^2}}{q}
\end{equation}

Подробное описание $j$-инварианта выходит за рамки этой книги. Для наших текущих целей достаточно отметить, что $j$-инвариант является важным инструментом для классификации эллиптических кривых и он необходим в \concept{методе комплексного умножения}, чтобы определить фактическую реализацию кривой, реализующую абстрактно выбранные свойства.

\begin{example} Рассмотрим эллиптическую кривую $E_{1,1}(\F_5)$ из \examplename{} \ref{E1F5}\sme{check reference}. Мы вычисляем её $j$-инвариант следующим образом:
\begin{align*}
j(E_{1,1}(\F_5)) &= \Zmod{1728\cdot \frac{4\cdot 1^3}{4\cdot 1^3+ 27\cdot 1^2}}{5}\\
             &= 3 \frac{4}{4+ 2}\\
             &= 3\cdot 4\\
             & = 2
\end{align*}
\end{example}
\begin{example} Рассмотрим эллиптическую кривую \TJJ{} из \examplename{} \ref{TJJ13}\sme{check reference}. Мы вычисляем её $j$-инвариант следующим образом:
\begin{align*}
j(TJJ\_13) &= \Zmod{1728\cdot \frac{4\cdot 8^3}{4\cdot 8^3+ 27\cdot 8^2}}{13}\\
             &= 12\cdot \frac{4\cdot 5}{4\cdot 5+ 1\cdot 12}\\
             &= 12\cdot \frac{7}{7+ 12}\\
             &= 12\cdot 7\cdot 6^{-1}\\
             &= 2\cdot 7\\
             &= 1 
\end{align*}
\end{example}
\begin{example}Рассмотрим \curvename{secp256k1} из \examplename{} \ref{secp256k1}. Мы вычисляем её $j$-инвариант с помощью Sage: 
\begin{sagecommandline}
sage: p = 115792089237316195423570985008687907853269984665640564039457584007908834671663
sage: F = GF(p)
sage: j = F(1728)*((F(4)*F(0)^3)/(F(4)*F(0)^3+F(27)*F(7)^2))
sage: j == F(0)
\end{sagecommandline}
\end{example} 
\begin{exercise}
\label{exercise:BN128-j-inv}
Рассмотрите кривую \curvename{alt\_bn128} из \examplename{} \ref{BN128}. Напишите программу на Sage, которая вычисляет $j$-инвариант для \curvename{alt\_bn128}.
\end{exercise}
\subsection{\concept{Метод комплексного умножения}}\label{complex-multiplication-method}
Как мы видели в предыдущих разделах, эллиптические кривые имеют различные определяющие свойства, такие как их порядок, простые множители, степень вложения или порядок базового поля. \term{\concept{Метод комплексного умножения (complex multiplication method)}} (CM) предоставляет практический способ построения эллиптических кривых с заранее заданными ограничениями на порядок кривой и базовое поле.\footnote{Подробное объяснение метода комплексного умножения
и его вывода можно найти, например, в \cite{grech-2012}.}

% the detailed method is here https://arxiv.org/pdf/1207.6983.pdf
% http://users.uoa.gr/~kontogar/files/ElisavetDaras.pdf
% https://hal.inria.fr/inria-00075302/PDF/RR-1256.pdf
% https://www.ams.org/journals/mcom/2007-76-260/S0025-5718-07-01980-1/S0025-5718-07-01980-1.pdf
% https://graui.de/code/elliptic2/ //draw curves
% https://hal.inria.fr/inria-00075302/PDF/RR-1256.pdf proposition 2.1

\concept{Метод комплексного умножения} начинается с выбора базового поля $\F_{q}$ кривой $E(\F_q)$, которую мы хотим построить, такого что $q = p^m$ для некоторого простого числа $p$ и $m\in \N$. Мы предполагаем $p>3$, чтобы упростить дальнейшее изложение. 

Затем выбирается след Фробениуса $t\in \Z$ кривой такой, что $q$ и $t$ взаимно просты, то есть выполняется $gcd(q,t)=1$ и $|t|\leq 2\sqrt{q}$. Выбор $t$ также определяет порядок кривой $r$, поскольку $r=p+1-t$ по границе Хассе \eqref{hasse-bound}, так что выбор $t$ определит большую подгруппу порядка, а также все малые кофакторы. Полученный $r$ должен быть таким, чтобы кривая удовлетворяла требованиям безопасности приложения. 

Заметьте, что выбор $p$ и $t$ также определяет степень вложения $k$ любой подгруппы простого порядка кривой, поскольку $k$ определяется как наименьшее число такое, что простой порядок $n$ делит число $p^k-1$.

Чтобы \concept{метод комплексного умножения} работал, ни $q$, ни $t$ не могут быть произвольными, но должны быть выбраны таким образом, чтобы существовали два дополнительных целых числа $D\in \Z$ и $v\in \Z$ и выполнялись следующие условия:

\begin{equation}\label{eq:D-criteria}
\begin{split}
D<0\\
\Zmod{\left( D =0  \text{ or } D=1\right)}{4}\\
4q  = t^2 + |D|v^2 
\end{split}
\end{equation}

Если такие числа существуют, мы называем $D$ \term{CM-дискриминантом (CM-discriminant)}, и мы знаем, что можем построить кривую $E(\F_q)$ над конечным полем $\F_q$ такую, что порядок кривой равен $|E(\F_q)|= q+1-t$. В этом случае цель \concept{метода комплексного умножения} — фактически построить такую кривую, то есть найти параметры $a$ и $b$ из $\F_q$ в определяющем уравнении Вейерштрасса такие, что кривая имеет желаемый порядок $r$. 

Уравнение \ref{eq:D-criteria} имеет бесконечное число решений, и проведено много исследований для вычисления множеств решений, которые приводят к эллиптическим кривым с желаемыми свойствами. Распространённый подход — зафиксировать CM-дискриминант и определить $q$, $t$ и $v$ как функции параметра $x$ таким образом, что значения $q(x)$, $t(x)$ и $v(x)$ дают решение \ref{eq:D-criteria} для каждого параметра $x$. На самом деле многие из таких семейств известны под названиями кривых BLS \cite{bls-02} или NMT \cite{mnt-84}, что указывает на то, что порядок базового поля $q$, а также след Фробениуса $t$ каждого члена в таком семействе кривых вычисляются схожими способами. Это подход, используемый, например, в \cite{freeman-06}, для вычисления кривых, дружественных к спариванию.
\begin{example}[Кривые BLS6]
\label{ex:bls6-params} Чтобы лучше понять, как параметризованные множества решений уравнения \ref{eq:D-criteria} порождают семейства эллиптических кривых, мы рассмотрим параметризацию, найденную авторами Баррето, Линном и Скоттом в 2002 году \cite{bls-02}. Их подход порождает кривые, дружественные к спариванию, с различными степенями вложения, и в частности кривые со степенью вложения $6$ с CM-дискриминантом $D=-3$. Члены этих семейств обычно кратко называются кривыми BLS6. 

Более точно, пусть многочлены $t,q\in \Q[x]$ определены как $t(x)=x+1$ и $q(x) = \frac{1}{3}(x-1)^2(x^2-x+1) + x$. Тогда следующее множество определяет множество решений уравнения \ref{eq:D-criteria} для $D=-3$ и $v= \sqrt{(4 q - t^2)/3}$: 
\begin{equation}
PARAM(BLS6) = \{t(x), q(x) \;|\; x\in \NN \text{ and } q(x)\in\NN\}
\end{equation}  
\end{example}

Предполагая, что найдены подходящие параметры $q$, $t$, $D$ и $v$, мы должны вычислить так называемый \term{многочлен класса Гильберта (Hilbert class polynomial)} $H_D\in \Z[x]$ CM-дискриминанта $D$, который является многочленом с целочисленными коэффициентами. Для этого мы сначала должны вычислить следующее множество:
\begin{equation}
\begin{aligned}
\label{def:hilbert_class_domain}
S(D)=\{(A,B,C)\;|\; A,B,C\in\Z,\; &D = B^2-4AC,\; gcd(A,B,C)=1, \\
&|B|\leq A \leq \sqrt{\frac{|D|}{3}},\; A\leq C, 
\text{ if } B< 0 \text{ then } |B| < A < C\}
\end{aligned}
\end{equation}

Один из способов вычислить это множество — сначала вычислить целое число $A_{max}= Floor(\sqrt{\frac{|D|}{3}})$, затем перебрать все целые числа $0\leq A\leq A_{max}$, а также все целые числа $-A\leq B \leq A$ и проверить, существует ли целое число $C$, удовлетворяющее уравнению $D = B^2-4AC$ и остальным требованиям из \ref{def:hilbert_class_domain}.

Для вычисления многочлена класса Гильберта необходима так называемая \term{$j$-функция ($j$-function)}, которая является комплексной функцией, определённой на верхней полуплоскости $\mathbb{H}$ комплексной плоскости $\mathbb{C}$, обычно записываемой следующим образом:

\begin{equation}\label{eq:j-invariant2}
j: \mathbb{H} \to \mathbb{C}
\end{equation}

Это означает, что $j$-функция принимает комплексные числа $(x +y\cdot i)$ с положительной мнимой частью $y>0$ в качестве входных данных и возвращает комплексное число $j(x+i\cdot y)$ в качестве результата.

$j$-функция тесно связана с $j$-инвариантом \ref{sec:j-inv} эллиптической кривой. Однако для целей этой книги не важно понимать $j$-функцию подробно. Мы можем использовать Sage для её вычисления аналогично тому, как мы использовали бы Sage для вычисления любой другой хорошо известной функции, например квадратного корня. Следует отметить, однако, что вычисление $j$-функции в Sage иногда склонно к ошибкам точности. Например, $j$-функция имеет корень в $\frac{-1+i\sqrt{3}}{2}$, который Sage аппроксимирует только приближённо. Поэтому, используя Sage для вычисления $j$-функции, нам нужно учитывать потерю точности и, возможно, округлять до ближайшего целого.

\begin{sagecommandline}
sage: z = ComplexField(100)(0,1)
sage: z # (0+1i)
sage: elliptic_j(z)
sage: # j-функция определена только для аргументов с положительной мнимой частью
sage: z = ComplexField(100)(1,-1)
sage: try:
....:     elliptic_j(z)
....: except PariError:
....:     pass
sage: # корень в (-1+i sqrt(3))/2
sage: z = ComplexField(100)(-1,sqrt(3))/2
sage: elliptic_j(z)
sage: elliptic_j(z).imag().round()
sage: elliptic_j(z).real().round()
\end{sagecommandline}

Имея способ вычислить $j$-функцию и предвычисленное множество $S(D)$, мы теперь можем вычислить многочлен класса Гильберта следующим образом:
\begin{equation}
H_D(x) = \Pi_{(A,B,C)\in S(D)} \left(x - j\left(\frac{-B + \sqrt{D}}{2A}\right)\right)
\end{equation}

Другими словами, мы перебираем все элементы $(A,B,C)$ из множества $S(D)$ и вычисляем $j$-функцию в точке $\frac{-B + \sqrt{D}}{2A}$, где $D$ — CM-дискриминант, который мы определили на предыдущем шаге. Результат определяет множитель многочлена класса Гильберта, и все множители перемножаются вместе.

Можно показать, что многочлен класса Гильберта является целочисленным многочленом, но фактические вычисления требуют арифметики высокой точности, чтобы избежать ошибок аппроксимации, которые обычно возникают в компьютерных аппроксимациях $j$-функции (как показано выше). Таким образом, в случае если вычисленный многочлен класса Гильберта не имеет целочисленных коэффициентов, нам нужно округлить результат до ближайшего целого. При условии, что использованная нами точность была достаточно высокой, результат будет верен.

На следующем шаге мы используем многочлен класса Гильберта $H_D\in \Z[x]$ и проецируем его на многочлен $H_{D,q}\in\F_q[x]$ с коэффициентами в базовом поле $\F_q$, выбранном на первом шаге. Мы делаем это, просто приводя коэффициенты по модулю $p$, то есть если $H_D(x)= a_mx^m +a_{m-1}x^{m-1}+\ldots + a_1 x + a_0$, мы вычисляем $q$-модуль каждого коэффициента
$\tilde{a}_j = \Zmod{a_j}{p}$, что даёт \term{спроецированный многочлен класса Гильберта (projected Hilbert class polynomial)} следующим образом:
$$
H_{D,p}(x)=\tilde{a}_mx^m +\tilde{a}_{m-1}x^{m-1}+\ldots + \tilde{a}_1 x + \tilde{a}_0
$$
Затем мы ищем корни $H_{D,p}$, поскольку каждый корень $j_0$ многочлена $H_{D,p}$ определяет семейство эллиптических кривых над $\F_q$, которые все имеют $j$-инвариант \ref{eq:j-invariant2}, равный $j_0$. Мы можем выбрать любой корень, поскольку все они определяют эллиптическую кривую. Однако некоторые кривые с правильным $j$-инвариантом могут иметь порядок, отличный от того, который мы изначально выбрали. Поэтому нам нужен способ определить кривую с правильным порядком. 

Чтобы вычислить кривую с правильным порядком, мы должны различать несколько различных случаев в зависимости от нашего выбора корня $j_0\in \F_q$ и CM-дискриминанта $D\in \Z$. Если $j_0\neq 0$ или $j_0\neq \Zmod{1728}{q}$, мы вычисляем $c_1=\frac{j_0}{(\Zmod{1728}{q}) -j_0}\in \F_q$, затем выбираем некоторый произвольный квадратичный невычет $c_2\in \F_q$ и некоторый произвольный кубический невычет $c_3\in \F_q$. 

Следующий список гарантированно определяет кривую с правильным порядком $r= q+1 -t$ для порядка поля $q$ и следа Фробениуса $t$, который мы изначально выбрали:

\begin{itemize}
\label{def:curve-order-frobenius}
\item Случай $j_0 \neq 0 $ и $j_0\neq \Zmod{1728}{q}$. Кривая с правильным порядком определяется одним из следующих уравнений:
\begin{equation}
y^2 = x^3 + 3c_1x + 2c_1 \text{\;\; or \;\; } y^2 = x^3 + 3c_1c_2^2x + 2c_1c_2^3
\end{equation}
\item Случай $j_0 = 0 $ и $D\neq -3$. Кривая с правильным порядком определяется одним из следующих уравнений:
\begin{equation}
y^2 = x^3 + 1 \text{\;\; or \;\; } y^2 = x^3 + c_2^3
\end{equation}
\item Случай $j_0 = 0 $ и $D= -3$. Кривая с правильным порядком определяется одним из следующих уравнений:
\begin{align*}
y^2 = x^3 +1 & \text{\;\; or \;\; } y^2 = x^3 + c_2^3 \text{ \;\; or}\\  
y^2 = x^3 + c_3^2 & \text{\;\; or \;\; } y^2 = c_3^2 c_2^3 \text{\;\; or}\\
y^2 = x^3 + c_3^{-2} & \text{\;\; or \;\; }  y^2 = x^3 + c_3^{-2}c_2^3 
\end{align*}
\item Случай $j_0 = \Zmod{1728}{q} $ и $D\neq -4$. Кривая с правильным порядком определяется одним из следующих уравнений:
\begin{equation}
y^2 = x^3 + x \text{\;\; or \;\; } y^2 = x^3 + c_2^2x
\end{equation}
\item Случай $j_0 = \Zmod{1728}{q} $ и $D= -4$. Кривая с правильным порядком определяется одним из следующих уравнений:
\begin{align*}
y^2 = x^3 +x & \text{\;\; or \;\; } y^2 = x^3 + c_2x \text{ \;\; or}\\  
y^2 = x^3 + c_2^2x & \text{\;\; or \;\; } y^2 = x^3 + c_2^3x
\end{align*}
\end{itemize} 

Чтобы определить правильное определяющее уравнение \concept{короткой формы Вейерштрасса}, мы поэтому должны вычислить порядок любой из потенциальных кривых выше, а затем выбрать ту, которая соответствует нашим изначальным требованиям. 

В заключение, используя \concept{метод комплексного умножения}, возможно синтезировать эллиптические кривые с заранее заданным порядком над заранее заданными базовыми полями с нуля. Однако кривые, построенные таким образом, являются лишь некоторыми представителями более широкого класса кривых, все из которых имеют один и тот же порядок. Поэтому в реальных приложениях иногда более выгодно выбрать другого представителя из этого класса. Для этого вспомним из \ref{sec:isomorphic_curves}, что любая кривая, определённая уравнением \concept{короткой формы Вейерштрасса} $y^2 = x^3 + ax + b$, изоморфна кривой вида $y^2 = x^3 + ad^4 x + bd^6$ для некоторого обратимого элемента поля $d\in \F^*_q$. 

Чтобы найти подходящего представителя (например, с малыми параметрами $a$ и $b$), разработчик кривой может выбрать обратимый элемент поля $d$ так, чтобы преобразованная кривая имела малые параметры.

\begin{example} Рассмотрим кривую $E_{1,1}(\F_5)$ из \examplename{} \ref{E1F5}. Мы хотим использовать \concept{метод комплексного умножения}, чтобы вывести эту кривую с нуля. Поскольку $E_{1,1}(\F_5)$ — кривая порядка $r=9$ над простым полем порядка $q=5$, мы знаем из \examplename{} \ref{ex:E1F5-frobenius}, что её след Фробениуса равен $t=-3$, что также показывает, что $q$ и $t$ взаимно просты. 

Затем мы должны найти параметры $D,v\in\Z$ такие, что выполняются критерии в \ref{eq:D-criteria}. Мы получаем следующее:
\begin{align*}
4q & = t^2+ |D|v^2 & \Rightarrow \\
20 & = (-3)^2 + |D|v^2 & \Leftrightarrow \\
11 & = |D|v^2
\end{align*}
Теперь, поскольку $11$ — простое число, единственным решением здесь является $|D|=11$ и $v=1$. При $D=-11$ и евклидовом делении $-11$ на $4$, равном $-11 = -3\cdot 4 +1$, мы имеем $\Zmod{-11}{4}=1$, что показывает, что $D=-11$ — правильный выбор.

На следующем шаге мы должны вычислить многочлен класса Гильберта $H_{-11}$. Для этого мы сначала должны найти множество $S(D)$. Чтобы вычислить это множество, заметим, что, поскольку $\sqrt{\frac{|D|}{3}}\approx 1.915<2$, мы знаем из $A\leq \sqrt{\frac{|D|}{3}}$ и $A\in\Z$, а также $0\leq |B|\leq A$, что $A$ должен быть либо $0$, либо $1$. 

Для $A=0$ мы знаем $B=0$ из ограничения $|B|\leq A$. Однако в этом случае не может существовать $C$, удовлетворяющего $-11= B^2 -4AC$. Поэтому мы пробуем $A=1$ и выводим $B\in\{-1,0,1\}$ из ограничения $|B|\leq A$. Случай $B=-1$ можно исключить, поскольку тогда $B<0$ должен влечь $|B|<A$. Случай $B=0$ также можно исключить, поскольку не может существовать целого $C$ с $-11 = -4C$, поскольку $11$ — простое число. 

Остаётся случай $B=1$, и мы вычисляем $C=3$ из уравнения $-11 = 1^2 -4C$, что даёт решение $(A,B,C)=(1,1,3)$. Следовательно:
$$
S(D)=\{(1,1,3)\}
$$

Имея в распоряжении множество $S(D)$, мы можем вычислить многочлен класса Гильберта для $D=-11$. Для этого мы должны подставить выражение $\frac{-1+\sqrt{-11}}{2\cdot1}$ в $j$-функцию. Для этого сначала заметим, что $\sqrt{-11}=i\sqrt{11}$, где $i$ — мнимая единица, определяемая $i^2=-1$. Используя это, мы используем Sage для вычисления $j$-инварианта и получаем следующее:
$$
H_{-11}(x) = x - j\left(\frac{-1+i\sqrt{11}}{2}\right) = x + 32768
$$

Как мы видим, в этом конкретном случае многочлен класса Гильберта является линейной функцией с единственным целочисленным коэффициентом. На следующем шаге мы должны спроецировать его на многочлен из $\F_5[x]$, приводя коэффициенты $1$ и $32768$ по модулю $5$. Мы получаем $\Zmod{32768}{5}=3$, так что спроецированный многочлен класса Гильберта равен
$$
H_{-11,5}(x)=x+3
$$ 
Как мы видим, единственным корнем этого многочлена является $j=2$, поскольку $H_{-11,5}(2)=2+3=0$. Поэтому мы имеем ситуацию с $j\neq 0$ и $j\neq 3 = \Zmod{1728}{5}$, что говорит нам, что мы должны рассмотреть первый случай в \ref{def:curve-order-frobenius} и вычислить параметр $c_1$, где в нашем случае деление выполняется в арифметике по модулю $5$:
\begin{equation}
c_1 = \frac{2}{\Zmod{1728}{5}-2} = \frac{2}{3-2}= 2
\end{equation}
Чтобы определить правильное уравнение из первого случая в \ref{def:curve-order-frobenius}, мы должны проверить, имеет ли кривая $E(\F_5)$, определённая уравнением \concept{короткой формы Вейерштрасса} $y^2 =  x^3 + 3\cdot c_1 x + 2\cdot c_1 = x^3 + x + 4$, правильный порядок. Мы используем Sage и находим, что порядок действительно равен $9$, так что это кривая с требуемыми параметрами. Таким образом, мы успешно построили кривую с желаемыми свойствами.

Сравнивая нашу построенную кривую $y^2 =  x^3 + x + 4$ с определением $E_{1,1}(\F_5)$ из \examplename{} \ref{E1F5}, мы видим, что определяющие уравнения различны. Однако, поскольку обе кривые имеют одинаковый порядок, мы знаем из \ref{sec:isomorphic_curves}, что они изоморфны. На самом деле мы можем использовать \ref{sec:isomorphic_curves} и квадратичный вычет $4\in \F_5$, чтобы преобразовать кривую, определённую уравнением $y^2 = x^3 +x+4$, в кривую $y^2 = x^3 + 4^2\cdot x + 4\cdot 4^3$, что даёт определяющее уравнение $E_{1,1}(\F_5)$:
$$
y^2 = x^3 + x +1
$$
Таким образом, используя \concept{метод комплексного умножения}, мы смогли вывести кривую с конкретными свойствами с нуля.
\end{example}

\begin{example}
 Рассмотрим кривую \curvename{Tiny-jubjub} \TJJ{} из \examplename{} \ref{TJJ13}. Мы хотим использовать \concept{метод комплексного умножения}, чтобы вывести эту кривую с нуля. Мы знаем из \examplename{} \ref{TJJ13}, что \TJJ{} — кривая порядка $r=20$ над простым полем порядка $q=13$, и мы знаем из \examplename{} \ref{ex:TJJ13-frobenius}, что её след Фробениуса равен $t=-6$. Это показывает, что $q$ и $t$ взаимно просты. 

Чтобы применить \concept{метод комплексного умножения}, мы должны найти параметры $D,v\in\Z$ такие, что выполняется \ref{eq:D-criteria}. Мы получаем следующее:
\begin{align*}
4q & = t^2+ |D|v^2 & \Rightarrow \\
4\cdot 13 & = (-6)^2+ |D|v^2 & \Rightarrow \\
52 & = 36 + |D|v^2 & \Leftrightarrow \\
16 & = |D|v^2
\end{align*}

Это уравнение имеет четыре решения для $(D,v)$, а именно $(-4,\pm 2)$ и $(-16,\pm 1)$. Рассматривая первые два решения, мы показываем в упражнении \ref{ex:Low_order_Hilbert_polys}, что $D=-4$ влечёт $j=1728$, и из \ref{def:curve-order-frobenius} мы знаем, что в этом случае построенная кривая определяется уравнением \concept{короткой формы Вейерштрасса} \ref{def_short_weierstrass_curve}, имеющим обращающийся в нуль параметр $b=0$. Мы поэтому можем заключить, что $D=-4$ не поможет нам реконструировать \TJJ{}, поскольку $D=-4$ будет порождать кривые порядка $20$, но все эти кривые имеют $b=0$.

Поэтому мы рассматриваем третье и четвёртое решение для $D=-16$. На следующем шаге мы должны вычислить многочлен класса Гильберта $H_{-16}$. Для этого мы сначала должны найти множество $S(D)$. Чтобы вычислить это множество, заметим, что, поскольку $\sqrt{\frac{|-16|}{3}}\approx 2.31<3$, мы знаем из $A\leq \sqrt{\frac{|-16|}{3}}$ и $A\in\Z$ с $0\leq |B|\leq A$, что $A$ должен быть в диапазоне $0..2$. Поэтому мы перебираем все возможные значения $A$ и все возможные значения $B$ при ограничениях $|B|\leq A$, и если $B<0$, то $|B|<A$.
Затем мы вычисляем потенциальные $C$ из уравнения $-16 = B^2 -4AC$. Мы получаем следующие два решения для $S(D)$:
% sage has precomputed Hilbert class polynomials 
% https://doc.sagemath.org/html/en/reference/databases/sage/databases/db_class_polynomials.html
$$
S(D)=\{(1,0,4),(2,0,2)\}
$$
Имея в распоряжении множество $S(D)$, мы можем вычислить многочлен класса Гильберта для $D=-16$. Мы можем использовать Sage для вычисления $j$-инварианта и получаем следующее:
\begin{align*}
H_{-16}(x) &= \left(x - j\left(\frac{i\sqrt{16}}{2}\right)\right)
 \left(x - j\left(\frac{i\sqrt{16}}{4}\right)\right) \\
           &= (x- 287496)(x-1728)
\end{align*}

Как мы видим, в этом конкретном случае многочлен класса Гильберта является квадратичной функцией с двумя целочисленными коэффициентами. На следующем шаге мы должны спроецировать его на многочлен над $\F_{13}$, вычисляя остаток по модулю $13$ коэффициентов $287496$ и $1728$. Мы получаем $\Zmod{287496}{13}=1$ и $\Zmod{1728}{13}=12$, что означает, что спроецированный многочлен класса Гильберта равен следующему:
$$
H_{-11,5}(x)=(x-1)(x-12)= (x+12)(x+1)
$$ 
Это считается многочленом из $\F_{13}[x]$. Таким образом, мы имеем два корня, а именно $j=1$ и $j=12$. Мы уже знаем, что $j=12$ — неправильный корень для построения кривой \curvename{Tiny-jubjub}, поскольку $\Zmod{1728}{13}=12$, и этот случай не совместим с кривой с $b\neq 0$. Поэтому мы выбираем $j=1$.

Другой способ определить правильный корень — вычислить $j$-инвариант кривой \curvename{Tiny-jubjub}. Мы получаем следующее:
\begin{align*}
j(\mathit{TJJ\_13}) & = 12\frac{4\cdot 8^3}{4\cdot 8^3+ 1\cdot 8^2}\\
                    & = 12\frac{4\cdot 5}{4\cdot 5+ 12}\\
                    & = 12\frac{7}{7+ 12}\\
                    & = 12\frac{7}{7-1}\\
                    & = 1
\end{align*}

Это равно корню $j=1$ многочлена класса Гильберта $H_{-16,13}$, как и ожидалось. Поэтому мы имеем ситуацию с $j\neq 0$ и $j\neq 1728$, что говорит нам, что мы должны рассмотреть первый случай в \ref{def:curve-order-frobenius} и вычислить параметр $c_1$, где в нашем случае деление выполняется в арифметике по модулю $13$:
$$
c_1=\frac{1}{12-1} = \frac{1}{11} =6
$$
Чтобы определить правильное уравнение из первого случая в \ref{def:curve-order-frobenius}, мы должны проверить, имеет ли кривая $E(\F_{13})$, определённая уравнением \concept{короткой формы Вейерштрасса} $y^2 = x^3 + 3\cdot 6 x + 2\cdot 6 = x^3 + 5x +12$, правильный порядок. Мы используем Sage и находим, что порядок равен $8$, а не $20$, как ожидалось, что означает, что след этой кривой равен $6$, а не $-6$. Поэтому мы должны рассмотреть второе уравнение из первого случая в \ref{def:curve-order-frobenius} и выбрать некоторый квадратичный невычет $c_2\in\F_{13}$. Мы выбираем $c_2=5$ и вычисляем уравнение \concept{короткой формы Вейерштрасса} $y^2 = x^3 + 5 c_2^2 x + 12 c_2^3$ следующим образом:
$$
y^2 = x^3 + 8 x + 5
$$
Мы используем Sage и находим, что порядок равен $20$, что действительно является правильным. Сравнивая нашу построенную кривую $y^2 =  x^3 + 8x + 5$ с определением \TJJ{} из \examplename{} \ref{TJJ13}, мы видим, что определяющие уравнения различны. Однако, поскольку обе кривые имеют одинаковый порядок, мы знаем из \ref{sec:isomorphic_curves}, что они изоморфны.

На самом деле мы можем использовать \ref{sec:isomorphic_curves} и квадратичный вычет $12\in \F_{13}$, чтобы преобразовать кривую, определённую уравнением $y^2 =  x^3 + 8x + 5$, в кривую $y^2 = x^3 + 12^2\cdot 8x + 5\cdot 12^3$, что даёт следующее:
$$
y^2 = x^3 + 8x +8
$$

Это кривая \curvename{Tiny-jubjub}, которую мы широко использовали на протяжении всей книги. Таким образом, используя \concept{метод комплексного умножения}, мы смогли вывести кривую с конкретными свойствами с нуля.
\end{example}

\begin{example} Чтобы рассмотреть реальный пример, мы хотим использовать \concept{метод комплексного умножения} в сочетании с Sage для вычисления \curvename{secp256k1} с нуля. Исходя из \examplename{} \ref{secp256k1}, мы решаем вычислить эллиптическую кривую над простым полем $\F_p$ порядка $r$ для следующих параметров безопасности:
\begin{align*}
p &= \scriptstyle 115792089237316195423570985008687907853269984665640564039457584007908834671663\\
r &= \scriptstyle 115792089237316195423570985008687907852837564279074904382605163141518161494337
\end{align*}
Согласно \examplename{} \ref{ex:secp256k1-trace}, это даёт следующий след Фробениуса для любой кривой, изоморфной \curvename{secp256k1}:
$$t = \scriptstyle 432420386565659656852420866390673177327$$ 

Мы также решаем, что хотим кривую вида $y^2 = x^3 + b$, то есть хотим, чтобы параметр $a$ был равен нулю. Тогда таблица \ref{def:curve-order-frobenius} влечёт, что $j$-инвариант нашей кривой должен быть равен нулю.

На первом шаге мы должны найти CM-дискриминант $D$ и некоторое целое число $v$ такие, что уравнение 
$
4p = t^2 +|D|v^2
$
выполняется. Поскольку мы стремимся к обращающемуся в нуль $j$-инварианту, первое, что стоит попробовать, — это $D=-3$. В этом случае мы можем вычислить $v^2 = (4p -t^2)/|D|$, и если $v^2$ окажется целым числом, имеющим квадратный корень $v$, мы закончили. Обратившись к Sage, мы вычисляем следующим образом:
\begin{sagecommandline}
sage: D = -3
sage: p = 115792089237316195423570985008687907853269984665640564039457584007908834671663
sage: r = 115792089237316195423570985008687907852837564279074904382605163141518161494337
sage: t = p+1-r
sage: v_sqr = (4*p - t^2)/abs(D)
sage: v_sqr.is_integer()
sage: v = sqrt(v_sqr)
sage: v.is_integer()
sage: 4*p == t^2 + abs(D)*v^2
sage: v
\end{sagecommandline}
Пара $(D,v)=(-3, 303414439467246543595250775667605759171)$ действительно решает уравнение, что говорит нам, что существует кривая порядка $r$ над простым полем порядка $p$, определённая уравнением \concept{короткой формы Вейерштрасса} $y^2 = x^3 + b$ для некоторого $b\in \F_p$. Теперь нам нужно вычислить $b$.

Для $D=-3$ мы показываем в упражнении \ref{ex:Low_order_Hilbert_polys}, что ассоциированный многочлен класса Гильберта задаётся как $H_{-3}(x)=x$, что даёт спроецированный многочлен класса Гильберта как 
$H_{-3,p}=x$, и $j$-инвариант нашей кривой гарантированно равен $j=0$. Теперь, глядя на \ref{def:curve-order-frobenius}\sme{check reference}, мы видим, что существует $6$ возможных случаев для построения кривой с правильным порядком $r$. Чтобы построить кривые вопроса, мы должны выбрать некоторые произвольные квадратичный и кубический невычеты. Поэтому мы перебираем $\F_p$, чтобы найти их, обратившись к Sage:

\begin{sagecommandline}
sage: F = GF(p)
sage: for c2 in F:
....:     try: # quadratic residue
....:         _ = c2.nth_root(2)
....:     except ValueError: # quadratic non-residue
....:         break
sage: c2
sage: for c3 in F:
....:     try:
....:         _ = c3.nth_root(3)
....:     except ValueError:
....:         break
sage: c3
\end{sagecommandline}

Мы нашли квадратичный невычет $c_2=3$ и кубический невычет $c_3=2$. Используя эти числа, мы проверяем шесть случаев против ожидаемого порядка $r$ кривой, которую мы хотим синтезировать:
\begin{sagecommandline}
sage: C1 = EllipticCurve(F,[0,1])
sage: C1.order() == r
sage: C2 = EllipticCurve(F,[0,c2^3])
sage: C2.order() == r
sage: C3 = EllipticCurve(F,[0,c3^2])
sage: C3.order() == r
sage: C4 = EllipticCurve(F,[0,c3^2*c2^3])
sage: C4.order() == r
sage: C5 = EllipticCurve(F,[0,c3^(-2)])
sage: C5.order() == r
sage: C6 = EllipticCurve(F,[0,c3^(-2)*c2^3])
sage: C6.order() == r
\end{sagecommandline}

Как и ожидалось, мы нашли эллиптическую кривую правильного порядка $r$ над простым полем размера $p$. В принципе, мы закончили, поскольку нашли кривую с теми же базовыми свойствами, что и \curvename{secp256k1}. Однако кривая определена следующим уравнением, которое использует очень большой параметр $b_1$, и поэтому она может работать слишком медленно в определённых алгоритмах и не очень удобна для работы людей.
$$
\scriptstyle y^2 = x^3 + 86844066927987146567678238756515930889952488499230423029593188005931626003754
$$
Поэтому может быть выгодно найти изоморфную кривую с наименьшим возможным параметром $b_2$. Чтобы найти такой $b_2$, мы можем использовать \ref{sec:isomorphic_curves} и выбрать обратимый квадратичный вычет $d$ такой, что $b_2 = b_1\cdot d^3$ является наименьшим возможным. Для этого мы переписываем последнее уравнение в следующую форму:
$$
d = \sqrt[3]{\frac{b_2}{b_1}}
$$ 

Затем мы используем Sage для перебора значений $b_2\in \F_p$, пока он не найдёт некоторое число такое, что частное $\frac{b_2}{b_1}$ имеет кубический корень $d$, и этот кубический корень сам является квадратичным вычетом. 
\begin{sagecommandline}
sage: b1=F(86844066927987146567678238756515930889952488499230423029593188005931626003754)
sage: for b2 in F:
....:     if b2 == 0: continue
....:     try:
....:         d = (b2/b1).nth_root(3)
....:         _ = d.nth_root(2) # test
....:     except ValueError: continue
....:     break # found it
sage: b2
\end{sagecommandline}
Действительно, наименьшее возможное значение равно $b_2=7$, и определяющее уравнение \concept{короткой формы Вейерштрасса} кривой над $\F_p$ с простым порядком $r$ равно 
$
y^2 = x^3 + 7
$,
которую мы можем назвать \curvename{secp256k1}. Как мы только что видели, \concept{метод комплексного умножения} достаточно мощен, чтобы выводить криптографически стойкие кривые, такие как \curvename{secp256k1}, с нуля.
\end{example}
%\section{Twists}
%I think this has to wait for Volume 2, due to timing constraints...
\begin{exercise}
\label{ex:Low_order_Hilbert_polys} Покажите, что многочлены класса Гильберта для CM-дискриминантов $D=-3$ и $D=-4$ задаются как $H_{-3,q}(x)=x$ и $H_{-4,q}= x-(\Zmod{1728}{q})$.
\end{exercise}
\begin{exercise}
Используйте метод комплексного умножения, чтобы построить эллиптическую кривую порядка $7$ над простым полем $\F_{13}$.
\end{exercise}
\begin{exercise}
Используйте метод комплексного умножения, чтобы вычислить все классы изоморфизма всех эллиптических кривых порядка $7$ над простым полем $\F_{13}$.
\end{exercise}
\begin{exercise}
\label{exercise:BN128-cm-method}
Рассмотрите простой модуль $p$ кривой \curvename{alt\_bn128} из \examplename{} \ref{BN128} и её след $t$ из упражнения \ref{exercise:BN128-cm-method}. Используйте метод комплексного умножения, чтобы синтезировать эллиптическую кривую над $F_p$, изоморфную \curvename{alt\_bn128}, и вычислите явный изоморфизм между этими двумя кривыми.
\end{exercise}
\subsection{Кривая $BLS6\_6$ для работы вручную}\label{BLS6}

% https://arxiv.org/pdf/1207.6983.pdf
% construction 6.6 in https://eprint.iacr.org/2006/372.pdf
В этом параграфе мы обобщаем наше понимание эллиптических кривых, чтобы вычислить основной пример для работы вручную на оставшуюся часть книги. Для этого мы хотим использовать \concept{метод комплексного умножения}, чтобы вывести кривую, дружественную к спариванию, которая имеет свойства, аналогичные кривым, используемым в реальных криптографических протоколах. Однако мы проектируем кривую специально так, чтобы она была полезна в примерах для работы вручную, что в основном означает, что кривая должна содержать лишь несколько точек, чтобы мы могли вывести исчерпывающие таблицы сложения и спаривания. Конкретно мы используем конструкцию 6.6 в \cite{freeman-2020}.

Хорошо изученное семейство кривых, дружественных к спариванию, задаётся множеством кривых BLS. Как объясняется в примере \ref{ex:bls6-params}, в \cite{bls-02} авторы Баррето, Линн и Скотт нашли параметризованное множество решений уравнения \ref{eq:D-criteria}, которое порождает кривые, дружественные к спариванию, с различными степенями вложения и CM-дискриминантом $D = -3$. 

Большинство реальных кривых BLS имеют степень вложения $12$, однако эта степень слишком велика для удобной кривой для работы вручную. К счастью, кривые BLS со степенями вложения $k$, удовлетворяющими $\kongru{k}{0}{6}$, вычисляются аналогичным образом, и поскольку наименьшая степень вложения $k$, удовлетворяющая этому сравнению, равна $k=6$, мы стремимся к кривой BLS со степенью вложения $6$ как к нашему основному примеру для работы вручную. Мы называем такую кривую кривой $BLS6$, поскольку это соглашение — записывать степень вложения сразу после дескриптора, дающего подсказку о том, как была построена кривая.  

\subsubsection{Конструкция}
Чтобы применить \concept{метод комплексного умножения} \ref{complex-multiplication-method}, вспомним, что он начинается с выбора базового поля $\F_{p^m}$, а также следа Фробениуса $t$ и CM-дискриминанта $D$ кривой. В случае кривых $BLS$ параметр $m$ выбирается равным $1$, что означает, что кривые определены над простыми полями. Кроме того, CM-дискриминант выбирается равным $D=-3$, что означает, что кривая определяется уравнением $y^2 = x^3 +b$ для некоторого $b\in \F_p$. 

Как показано в примере \ref{ex:bls6-params}, для кривых $BLS6$ соответствующие параметры $p$ и $t$ сами параметризованы следующими функциями:
\begin{equation}
\begin{split}
t(x) &= x+1\\
p(x) &= \frac{1}{3}(x-1)^2(x^{2}-x+1) +x
\end{split}
\end{equation}
Здесь $x\in\NN$ — параметр, который разработчик должен выбрать таким образом, чтобы вычисление $p$ и $t$ в точке $x$ давало целые числа надлежащего размера, чтобы удовлетворять требованиям безопасности синтезируемой кривой. Гарантируется, что существует целое число $v\in\Z$ такое, что уравнение \ref{eq:D-criteria} выполняется для CM-дискриминанта $D=-3$, что означает, что \concept{метод комплексного умножения} может быть использован для вычисления элемента поля $b\in\F_p$ такого, что эллиптическая кривая $y^2 = x^3 +b$ имеет порядок $r=p+1-t$, CM-дискриминант $D=-3$ и степень вложения $k=6$. 

Чтобы спроектировать наименьшую кривую BLS6, мы поэтому должны найти параметр $x$ такой, что $t(x)$ и $p(x)$ — наименьшие натуральные числа, удовлетворяющие $p(x)>3$.\footnote{Наименьшая кривая BLS будет также наиболее нестойкой кривой BLS. Однако, поскольку наша цель с этой кривой — удобство вычислений вручную, а не безопасность, она соответствует целям этой книги.}

Мы поэтому начинаем процесс проектирования нашей кривой $BLS6$, подставляя малые значения $x$ в определяющие многочлены $t$ и $q$. Мы получаем следующие результаты:
$$
\begin{array}{lcr}
x=1 & (t(x),p(x)) & (2,1)\\
x=2 & (t(x),p(x)) & (3,3)\\
x=3 & (t(x),p(x)) & (4,\frac{37}{3})\\
x=4 & (t(x),p(x)) & (5,43)\\
\end{array}
$$
Поскольку $p(1)=1$ не является простым числом, первое $x$, дающее правильную кривую, — это $x=2$. Однако такая кривая была бы определена над базовым полем характеристики $3$, и мы бы предпочли избежать этого, поскольку в этой книге мы определяли эллиптические кривые только над полями характеристики больше $3$. Поэтому мы находим $x=4$, которое определяет кривую над простым полем характеристики $43$, имеющую след Фробениуса $t=5$. 

Поскольку простое поле $\F_{43}$ имеет $43$ элемента, а двоичное представление $43$ равно $43_2= 101011$, которое состоит из $6$ цифр, название нашей кривой для работы вручную должно быть $BLS6\_6$, поскольку обычно кривую BLS называют по её степени вложения и битовой длине модуля в базовом поле. Мы называем $BLS6\_6$ \term{кривой moon-math}.

Исходя из \ref{hasse-bound}\sme{check reference}, мы знаем, что граница Хассе влечёт, что $BLS6\_6$ будет содержать ровно $39$ элементов. Поскольку разложение на простые множители $39$ равно $39=3\cdot 13$, мы имеем ``большую'' группу простого множителя размера $13$ и малый кофактор $3$. К счастью, подгруппа порядка $13$ хорошо подходит для наших целей, поскольку $13$ элементов можно легко обрабатывать в связанных таблицах сложения, скалярного умножения и спаривания в стиле работы вручную. 

Мы можем проверить, что степень вложения действительно равна $6$, как и ожидалось, поскольку $k = 6$ — наименьшее число $k$ такое, что $r=13$ делит $43^k-1$. 
\begin{sagecommandline}
sage: k= 0
sage: for k in range(1,42): # малая теорема Ферма
....:     if (43^k-1)%13 == 0:
....:         break
sage: k
\end{sagecommandline}

Чтобы увидеть, что уравнение \ref{eq:D-criteria} действительно имеет решение для параметров $D=-3$, $p=43$ и $t=5$, как и ожидалось, мы вычисляем следующим образом:  
\begin{align*}
4p & = t^2 + |D|v^2 & \Rightarrow \\ 
4\cdot 43 & = 5^2 + 3\cdot v^2 & \Leftrightarrow \\ 
172 & = 25 + 3 v^2 & \Leftrightarrow \\ 
49 & = v^2 & \Leftarrow \\
v & = \pm 7
\end{align*}
Это означает, что мы можем использовать \concept{метод комплексного умножения}, как описано в \ref{complex-multiplication-method}, чтобы вычислить определяющее уравнение $y^2=x^3 + ax + b$ кривой BLS6\_6. 

Поскольку $D=-3$, мы знаем из упражнения \ref{ex:Low_order_Hilbert_polys}, что ассоциированный многочлен класса Гильберта задаётся как $H_{-3,43}(x) = x$, что означает, что $j$-инвариант $BLS6\_6$ равен $j(BLS6\_6)=0$. Поэтому мы должны рассмотреть третий случай в таблице \ref{def:curve-order-frobenius}, чтобы вывести $a=0$ и получить параметр $b$. Чтобы применить \ref{def:curve-order-frobenius} и определить правильное уравнение для $j_0=0$ и $D=-3$, мы должны выбрать некоторые произвольные квадратичный невычет $c_2$ и кубический невычет $c_3$ в $\F_{43}$. Мы выбираем $c_2 =5$ и $c_3=36$. Мы проверяем их с помощью Sage:
\begin{sagecommandline}
sage: F43 = GF(43)
sage: c2 = F43(5)
....: try: # квадратичный вычет
....:     c2.nth_root(2)
....: except ValueError: # квадратичный невычет
....:     print("OK") 
sage: c3 =F43(36)
....: try:
....:     c3.nth_root(3)
....: except ValueError:
....:     print("OK") 
\end{sagecommandline} 

Используя эти числа, мы проверяем шесть возможных случаев из \ref{def:curve-order-frobenius}\sme{check reference} против ожидаемого порядка $39$ кривой, которую мы хотим синтезировать:

\begin{sagecommandline}
sage: BLS61 = EllipticCurve(F43,[0,1])
sage: BLS61.order() == 39
sage: BLS62 = EllipticCurve(F43,[0,c2^3])
sage: BLS62.order() == 39
sage: BLS63 = EllipticCurve(F43,[0,c3^2])
sage: BLS63.order() == 39
sage: BLS64 = EllipticCurve(F43,[0,c3^2*c2^3])
sage: BLS64.order() == 39
sage: BLS65 = EllipticCurve(F43,[0,c3^(-2)])
sage: BLS65.order() == 39
sage: BLS66 = EllipticCurve(F43,[0,c3^(-2)*c2^3])
sage: BLS66.order() == 39
sage: BLS6 = BLS63 # our BLS6 curve in the book
\end{sagecommandline}
Как и ожидалось, мы нашли эллиптическую кривую правильного порядка $39$ над простым полем размера $43$. Поскольку $c_3^2=36^2=6$ в $\F_{43}$, кривая $BLS6\_6$ определяется следующим уравнением:

\begin{equation}
BLS6\_6 := \{(x,y)\;|\; y^2 = x^3 + 6 \text{ for all } x,y \in \F_{43}\}
\end{equation}

Поскольку существуют другие выборы для $c_3$, существуют и другие выборы для $b$, такие как $b=10$ или $b=23$. Однако благодаря \ref{sec:isomorphic_curves} все эти кривые изоморфны и поэтому представляют одну и ту же кривую различными способами. Мы выбрали $b=6$ без какой-либо особой причины.

Поскольку BLS6\_6 содержит всего $39$ точек, возможно использовать Sage, чтобы дать визуальное представление кривой:

\begin{sagesilent}
BLS63p = BLS63.plot()
\end{sagesilent}
\begin{center} 
\sageplot[scale=.5]{BLS63p} 
\end{center}

Как мы видим, наша кривая имеет некоторые желательные свойства: она не содержит самообратных точек, то есть точек с $y=0$. Следовательно, закон сложения может быть оптимизирован, поскольку ветвь для этих случаев может быть устранена. 

\subsubsection{Большая подгруппа простого порядка}
Обобщая предыдущую процедуру, мы использовали метод Баррето, Линна и Скотта\sme{add reference}, чтобы построить кривую, дружественную к спариванию, со степенью вложения $6$. Однако, чтобы выполнять криптографию на эллиптических кривых на этой кривой, заметим, что, поскольку порядок $BLS6\_6$ равен $39$, её группа не имеет простого порядка. Поэтому мы должны найти подходящую подгруппу в качестве нашей основной цели. Поскольку $39=13\cdot 3$, мы знаем, что кривая должна содержать ``большую'' группу простого порядка размера $13$, называемую $BLS6\_6(\F_{43})[13]$ согласно определению \ref{def:torsion_group}, и малую кофактор-группу порядка $3$. Мы используем следующее обозначение для большой группы простого порядка
\begin{equation}
\label{def:BLS6613}
\G_1[13]:= BLS6\_6(\F_{43})[13]
\end{equation}

Один из способов вычислить эту группу — найти образующий. Мы можем достичь этого, выбрав произвольный элемент $BLS6\_6$, не являющийся точкой на бесконечности, а затем умножив эту точку на кофактор $3$ числа $13$. Если результат не является точкой на бесконечности, это будет образующий $\G_1[13]$. Если это точка на бесконечности, мы должны выбрать другой элемент. 

Чтобы построить такой элемент из $BLS6\_6$ в стиле работы вручную, мы можем выбрать некоторое $x\in\F_{43}$ и проверить, существует ли некоторое $y\in\F_{43}$, удовлетворяющее определяющему уравнению \concept{короткой формы Вейерштрасса} $y^2 = x^3 + 6$. Мы выбираем $x=9$ и проверяем, что $y=2$ удовлетворяет уравнению кривой для $x$:
\begin{align*}
y^2 & = x^3 + 6 & \Rightarrow \\
2^2 & = 9^3 + 6 & \Leftrightarrow \\
4 & = 4
\end{align*}   

Это означает, что $P=(9,2)$ является точкой на $BLS6\_6$. Чтобы увидеть, можем ли мы спроецировать эту точку на образующего большой группы простого порядка $\G_1[13]$, мы должны умножить $P$ на кофактор $3$, то есть мы должны вычислить $[3](9,2)$. Мы получаем $[3](9,2) = (13,15)$ (см. упражнение \ref{ex:BLS66-scalar-product}). Поскольку это не точка на бесконечности, мы знаем, что $(13,15)$ — образующий $\G_1[13]$, которого мы будем использовать на протяжении всей книги: 

\begin{equation}\label{gBLS6-6-13}
g_{1} = (13,15)
\end{equation}

Поскольку $g_{1}$ — образующий, вспомним из \ref{exponentialmap}, что существует экспоненциальное отображение из поля $\F_{13}$ в $\G_1[13]$ относительно этого образующего, которое порождает группу в логарифмическом порядке:
$$
[\cdot]_{(13,15)}: \F_{13} \to \G_1[13]\;;\; x \mapsto [x](13,15)
$$ 
Мы можем использовать эту функцию для построения подгруппы $\G_1[13]$, многократно добавляя образующего к самому себе. Используя Sage, мы получаем следующее:
\begin{sagecommandline}
sage: P1 = BLS6(9,2)
sage: g1 = 3*P1 # generator
sage: g1.xy()
sage: G1_13 = [ x*g1 for x in range(0,13) ]
\end{sagecommandline}
Многократное сложение образующего с самим собой порождает малые группы в логарифмическом порядке относительно образующего, как объяснено в \ref{def:logarithmic_ordering}. Это даёт следующее представление:
\begin{multline}
\label{BLS6-G1-log}
\G_1[13]=
\{(13,15) \rightarrow (33,34) \rightarrow  (38,15) \rightarrow  (35,28) \rightarrow (26,34) \rightarrow  (27,34) \rightarrow  \\ 
(27,9)  \rightarrow  (26,9) \rightarrow  (35,15) \rightarrow  (38,28) \rightarrow  (33,9) \rightarrow (13,28) \rightarrow  \mathcal{O}\}$$
\end{multline}
Наличие логарифмического порядка этой группы полезно при вычислениях вручную. Чтобы увидеть это, рассмотрим следующий пример:
\begin{align*}
(27,34)\oplus (33,9)  & = [6](13,15)\oplus [11](13,15)\\
                      & = [6+11](13,15)\\
                      & = [4](13,15)\\
                      & = (35,28)\\
\end{align*}
Как показывает это вычисление, \ref{BLS6-G1-log} — это действительно всё, что нам нужно для эффективных вычислений в $\G_1[13]$. $\G_1[13]$ поэтому подходит в качестве криптографической группы для работы вручную. Однако, для удобства, следующая таблица перечисляет полную таблицу сложения группы $BLS6\_6$, поскольку она может быть полезна в некоторых вычислениях вручную, которые не ограничены подгруппой $\G_1$:
\begingroup
    \fontsize{7pt}{7pt}\selectfont
$$
\begin{array}{c|ccccccccccccc}
\oplus & \mathcal{O}  & (13,15) & (33,34) & (38,15) & (35,28) & (26,34) & (27,34) & (27,9) & (26,9) & (35,15) & (38,28) & (33,9) & (13,28)\\
\hline
\\
\mathcal{O} & \mathcal{O}  & (13,15) & (33,34) & (38,15) & (35,28) & (26,34) & (27,34) & (27,9) & (26,9) & (35,15) & (38,28) & (33,9) & (13,28)\\
\\
(13,15) & (13,15) & (33,34) & (38,15) & (35,28) & (26,34) & (27,34) & (27,9) & (26,9) & (35,15) & (38,28) & (33,9) & (13,28) & \mathcal{O}\\
\\
(33,34) & (33,34) & (38,15) & (35,28) & (26,34) & (27,34) & (27,9) & (26,9) & (35,15) & (38,28) & (33,9) & (13,28) & \mathcal{O} & (13,15)\\
\\
(38,15) & (38,15) & (35,28) & (26,34) & (27,34) & (27,9) & (26,9) & (35,15) & (38,28) & (33,9) & (13,28) & \mathcal{O} & (13,15) & (33,34)\\
\\
(35,28) & (35,28) & (26,34) & (27,34) & (27,9) & (26,9) & (35,15) & (38,28) & (33,9) & (13,28) & \mathcal{O} & (13,15) & (33,34) & (38,15)\\
\\
(26,34) & (26,34) & (27,34) & (27,9) & (26,9) & (35,15) & (38,28) & (33,9) & (13,28) & \mathcal{O} & (13,15) & (33,34) & (38,15) & (35,28)\\
\\
(27,34) & (27,34) & (27,9) & (26,9) & (35,15) & (38,28) & (33,9) & (13,28) & \mathcal{O} & (13,15) & (33,34) & (38,15) & (35,28) & (26,34)\\
\\
(27,9) & (27,9) & (26,9) & (35,15) & (38,28) & (33,9) & (13,28) & \mathcal{O} & (13,15) & (33,34) & (38,15) & (35,28) & (26,34) & (27,34)\\
\\
(26,9) & (26,9) & (35,15) & (38,28) & (33,9) & (13,28) & \mathcal{O} & (13,15) & (33,34) & (38,15) & (35,28) & (26,34) & (27,34) & (27,9)\\
\\
(35,15) & (35,15) & (38,28) & (33,9) & (13,28) & \mathcal{O} & (13,15) & (33,34) & (38,15) & (35,28) & (26,34) & (27,34) & (27,9) & (26,9)\\
\\
(38,28) & (38,28) & (33,9) & (13,28) & \mathcal{O} & (13,15) & (33,34) & (38,15) & (35,28) & (26,34) & (27,34) & (27,9) & (26,9) & (35,15)\\
\\
(33,9) & (33,9) & (13,28) & \mathcal{O} & (13,15) & (33,34) & (38,15) & (35,28) & (26,34) & (27,34) & (27,9) & (26,9) & (35,15) & (38,28)\\
\\
(13,28) & (13,28) & \mathcal{O} & (13,15) & (33,34) & (38,15) & (35,28) & (26,34) & (27,34) & (27,9) & (26,9) & (35,15) & (38,28) & (33,9)\\
\end{array}
$$
\endgroup
\begin{exercise}
\label{ex:BLS66-scalar-product}
Рассмотрите точку $P=(9,2)$. Покажите, что $P$ является точкой на кривой $BLS6\_6$, и вычислите скалярное произведение $[3]P$.
\end{exercise}
\begin{exercise}
Вычислите следующие выражения: $-(26,34)$, $(26,9)\oplus(13,28)$, $(35,15)\oplus \mathcal{O}$ и $(27,9)\ominus(33,9)$.
\end{exercise}
%To see how the small prime-order group of $BLS6\_6$ looks like we can apply the same approach but for the cofactor $13$ instead (EXERCISE). We get 
%$$BLS6\_6[3]=\{\mathcal{O},(0,7),(0,36)\}$$
%Now that we ha
\subsubsection{Группы спаривания}
\label{sec:bls66-pairing-groups}
Мы знаем, что $BLS6\_6$ является кривой, дружественной к спариванию, по конструкции, поскольку она имеет малую степень вложения $k=6$. Поэтому возможно эффективно вычислять спаривания Вейля. Однако, чтобы сделать это, мы должны определить группы спаривания $\G_1$ и $\G_2$, как определено в разделе \ref{sec:pairing_groups}. 

Поскольку $BLS6\_6$ имеет две нетривиальные подгруппы, можно было бы использовать любую из них в качестве $n$-группы кручения. Однако в криптографии единственным безопасным выбором является использование большой подгруппы простого порядка, которая в нашем случае является $\G_1[13]$, как представлено в \ref{BLS6-G1-log}. Поэтому мы решаем рассмотреть $13$-кручение для спаривания Вейля и использовать $\G_1[13]$ в качестве его левого аргумента.

Чтобы построить область определения для правого аргумента, нам нужно построить $\G_2[13]$, которая, согласно общей теории \ref{sec:pairing_groups}, должна быть определена теми элементами $P$ полной $13$-группы кручения $BLS6\_6[13]$, которые отображаются в $43\cdot P$ под действием эндоморфизма Фробениуса \ref{eq:frobenius-enomorphism}. 

Чтобы вычислить $\G_2[13]$, мы поэтому сначала должны найти полную $13$-группу кручения. Для этого мы используем технику из \ref{sec:full-torsion}, которая говорит нам, что полное $13$-кручение может быть найдено в расширении кривой $BLS6\_6(\F_{43^6})$ \ref{sec:curve-extensions} кривой $BLS6\_6$ над расширенным полем $\F_{43^6}$ \ref{field-extension}, поскольку степень вложения $BLS6\_6$ равна $6$. Поэтому мы должны рассмотреть следующую эллиптическую кривую:

\begin{equation}
\label{def:extended-bls6}
BLS6\_6(\F_{43^6}) := \{(x,y)\;|\; y^2 = x^3 + 6 \text{ for all } x,y \in \F_{43^6}\}
\end{equation}

Чтобы вычислить эту кривую, мы должны построить $\F_{43^6}$ — поле, содержащее более 6 миллиардов элементов. Мы используем общую конструкцию расширений простых полей из \ref{field-extension} и начинаем с выбора неприводимого многочлена степени $6$ из кольца многочленов $\F_{43}[t]$. Мы выбираем $p(t) = t^6+6$. Чтобы визуально отличать многочлены с коэффициентами в $\F_{43}$ от элементов в $\F_{43^6}$, мы используем символ $v$ исключительно для представления неопределённой в $\F_{43^6}$. Используя Sage, мы получаем следующее:
\begin{sagecommandline}
sage: F43 = GF(43)
sage: F43t.<t> = F43[]
sage: p = F43t(t^6+6)
sage: p.is_irreducible()
sage: F43_6.<v> = GF(43^6, name='v', modulus=p)
sage: F43_6.order()
\end{sagecommandline}

Вспомним из \ref{eq:prime-extension-field}, что элементы $x\in\F_{43^6}$ могут быть представлены как многочлены $a_0+a_1v + a_2v^2+\ldots + a_5 v^5$ с обычным сложением многочленов и умножением по модулю $v^6+6$. 

Чтобы вычислить $\G_2[13]$, мы сначала должны расширить $BLS6\_6$ до $\F_{43^6}$, то есть мы сохраняем определяющее уравнение, но расширяем область определения из $\F_{43}$ в $\F_{43^6}$. После этого мы должны найти по крайней мере один элемент $P$ из этой кривой, который не является точкой на бесконечности, принадлежит полному $13$-кручению и удовлетворяет тождеству $\pi(P) = [43]P$, где $\pi$ — эндоморфизм Фробениуса \ref{eq:frobenius-enomorphism}. Затем мы можем использовать этот элемент как наш образующий $\G_2[13]$ и построить все остальные элементы, многократно добавляя образующего к самому себе.

Поскольку $BLS6(\F_{43^6})$ содержит примерно столько же элементов, сколько $\F_{43^6}$, по границе Хассе \ref{hasse-bound}, не является хорошей стратегией просто перебирать все элементы. К счастью, Sage имеет способ перебирать элементы непосредственно из группы кручения:

\begin{sagecommandline}
sage: ExtBLS6 = EllipticCurve(F43_6,[0 ,6]) # расширение кривой
sage: INF = ExtBLS6(0) # точка на бесконечности
sage: for P in INF.division_points(13): # полное 13-кручение
....:     # pi(P) == [q]P
....:     if P.order() == 13: # исключить точку на бесконечности
....:         piP = ExtBLS6([a.frobenius() for a in P])
....:         qP = 43*P
....:         if piP == qP:
....:             break
sage: P.xy()
\end{sagecommandline}

Мы нашли элемент $P$ из полного $13$-кручения со свойством $\pi(P) = [43]P$, что означает, что он является элементом $\G_2[13]$. Поскольку $\G_2[13]$ циклическая и имеет простой порядок, этот элемент должен быть образующим:
\begin{equation}
g_{2} = (7v^2, 16v^3)
\end{equation}

Мы можем использовать этот образующий, чтобы вычислить $\G_2[13]$ в логарифмическом порядке относительно $g_{2}$. Используя Sage, мы получаем следующее:
\begin{sagecommandline}
sage: g2 = ExtBLS6(7*v^2,16*v^3)
sage: G2_13 = [ x*g2 for x in range(0,13) ]
\end{sagecommandline}
Многократное сложение образующего с самим собой порождает малые группы в логарифмическом порядке относительно образующего, как объяснено в \ref{def:logarithmic_ordering}. Поэтому мы получаем следующее представление:
\begin{multline}
\label{BLS6-G2-log}
\mathbb{G}_2=\{
(7v^2, 16v^3) \to
(10v^2, 28v^3)\to
(42v^2, 16v^3)\to
(37v^2, 27v^3)\to\\
(16v^2, 28v^3)\to
(17v^2, 28v^3)\to
(17v^2, 15v^3)\to
(16v^2, 15v^3)\to\\
(37v^2, 16v^3)\to
(42v^2, 27v^3)\to
(10v^2, 15v^3)\to
(7v^2, 27v^3)\to
\mathcal{O}\}
\end{multline}

Снова, наличие логарифмического описания $\G_2[13]$ полезно при вычислениях вручную, поскольку оно сводит сложные вычисления на расширенной кривой к арифметике по модулю $13$, как в следующем примере:
\begin{align*}
(17v^2,28v^3)\oplus (10v^2,15v^2)  & = [6](7v^2,16v^3)\oplus [11](7v^2,16v^3)\\
                      & = [6+11](7v^2,16v^3)\\
                      & = [4](7v^2,16v^3)\\
                      & = (37v^2,27v^3)\\
\end{align*}
Как показывает это вычисление, \ref{BLS6-G2-log} — это действительно всё, что нам нужно для эффективных вычислений в $\G_2[13]$. Группы $\G_1[13]$ и $\G_2[13]$ поэтому подходят в качестве криптографических групп спаривания для работы вручную. 

\subsubsection{Спаривание Вейля}
В \ref{sec:bls66-pairing-groups} мы вычислили две различные группы, $\G_1[13]$ и $\G_2[13]$, которые являются подгруппами полной $13$-группы кручения расширенной кривой moon-math $BLS6\_6(\F_{43^6})$. Как объяснено в \ref{sec:weil-pairing}, это означает, что существует спаривание Вейля 
\begin{equation}
\label{BLS6-weil-pairing}
e(\cdot,\cdot): \G_1[13]\times\G_2[13] \to \F_{43^6}
\end{equation}

Поскольку мы знаем логарифмический порядок \ref{BLS6-G1-log} для $\G_1[13]$ и логарифмический порядок \ref{BLS6-G2-log} для $\G_2[13]$, это спаривание Вейля эффективно вычислимо в вычислениях вручную, используя следующее тождество, выведенное в упражнении \ref{ex:pairing-arithmetics}:  
\begin{equation}\label{eq:bilinearityBLS6}
e([m]g_{1},[n]g_{2}) = 
e(g_{1},g_{2})^{\Zmod{m\cdot n}{13}}
\end{equation}
Во многих системах доказательств с нулевым разглашением на основе спариваний, таких как \ref{sec:gorth_16}, необходимо лишь показать равенство различных спариваний, и в вычислениях вручную точное значение $e(g_{1},g_{2})$ поэтому не важно. Чтобы увидеть, как это упрощает вычисление, предположим, что мы хотим доказать тождество $e((27,34),(16v^2,28v^3))=e((26,9),(17v^2,15v^3))$. В этом случае мы можем вычислить
\begin{align*}
e((27,34),(16v^2,28v^3)) 
               & = e([6](13,15),[5](7v^2,16v^3))\\
               & = e((13,15),(7v^2,16v^3))^{6\cdot 5}\\
               & = e((13,15),(7v^2,16v^3))^{4}\\
               & = e((13,15),(7v^2,16v^3))^{8\cdot 7}\\
               & = e([8](13,15),[7](7v^2,16v^3))\\
               & = e((26,9),(17v^2,15v^3))\\
\end{align*}
Если нужно фактическое значение $e(g_{1},g_{2})$, спаривание Вейля может быть вычислено либо с помощью уравнения \ref{eq:weil-pairing} и выполнения алгоритма Миллера \ref{alg:millersalgo} (Упражнение \ref{ex:generator-pairing}), либо можно обратиться к Sage:  
\begin{sagecommandline}
sage: g1 = ExtBLS6([13,15])
sage: g2 = ExtBLS6([7*v^2, 16*v^3])
sage: g1.weil_pairing(g2,13)
\end{sagecommandline}

\begin{exercise}
\label{ex:generator-pairing}
Рассмотрите расширенную кривую $BLS6\_6$, определённую в \ref{def:extended-bls6}, и две точки кривой $g1=(13,15)$ и $g2=(7v^2,16v^3)$. Вычислите спаривание Вейля 
$e(g1,g2)$, используя определение \ref{eq:weil-pairing} и алгоритм Миллера \ref{alg:millersalgo}.
\end{exercise}
\begin{comment}
\subsection{Хеширование в группы спаривания}
We give various constructions to hash into $\mathbb{G}_1$ and $\mathbb{G}_2$. 

We start with hashing to the scalar field... \smelong{TO APPEAR}\sme{finish writing this up}

None of these techniques work for hashing into $\mathbb{G}_2$. We therefore implement Pederson's Hash for BLS6. 

We start with $\mathbb{G}_1$. Our goal is to define an $12$-bit bounded hash function:
$$
H_{1}: \{0,1\}^{12} \to \mathbb{G}_1 
$$
Since $12= 3\cdot 4$ we ``randomly'' select $4$ uniformly distributed generators $\{(38, 15), (35,28),\\ (27, 34), (38, 28)\}$ from $\mathbb{G}_1$ and use the pseudo-random function from XXX\sme{add reference}. 
Therefore, we have to choose a set of $4$ randomly generated invertible elements from $\F_{13}$ for every generator. We choose the following:
$$
\begin{array}{lcl}
(38,15) &:& \{2,7,5,9\}\\
(35,28) &:& \{11,4,7,7\}\\
(27,34) &:& \{5,3,7,12\}\\
(38,28) &:& \{6,5,1,8\}
\end{array}
$$
Our hash function is then computed as follows:

\begin{multline*}
H_1(x_{11},x_1,\ldots, x_{0})=
[2\cdot 7^{x_{11}}\cdot 5^{x_{10}}\cdot 9^{x_9}](38,15)+
[11\cdot 4^{x_8}\cdot 7^{x_7}\cdot 7^{x_6}](35,28)+\\
[5\cdot 3^{x_5}\cdot 7^{x_4}\cdot 12^{x_3}](27,34) +
[6\cdot 5^{x_2}\cdot 1^{x_{1}}\cdot 8^{x_{0}}](38,28)
\end{multline*}

Note that $a^x=1$ when $x=0$. Hence, those terms can be omitted in the computation. 
In particular, the hash of the $12$-bit zero string is given as follows:\sme{correct computations}

\begin{multline*}\smelong{WRONG-ORDERING-REDO}\\
H_1(0)= [2](38,15)+[11](35,28)+[5](27,34)+[6](38,28)= \\
(27,34)+(26,34)+(35,28)+(26,9)= (33,9) + (13,28) = (38,28)
\end{multline*}

The hash of $011010101100$ is given as follows:\sme{fill in missing parts}
\begin{multline*}
H_1(011010101100)=\smelong{WRONG-ORDERING-REDO}\\
[2\cdot 7^{0}\cdot 5^{1}\cdot 9^{1}](38,15)+
[11\cdot 4^{0}\cdot 7^{1}\cdot 7^{0}](35,28)+
[5\cdot 3^{1}\cdot 7^{0}\cdot 12^{1}](27,34) +
[6\cdot 5^{1}\cdot 1^{0}\cdot 8^{0}](38,28)=\\
[2\cdot 5\cdot 9](38,15)+
[11\cdot 7](35,28)+
[5\cdot 3\cdot 12](27,34) +
[6\cdot 5](38,28)=\\
[12](38,15)+
[12](35,28)+
[11](27,34) +
[4](38,28)=\\ 
\smelong{TO APPEAR}
\end{multline*}
We can use the same technique to define a $12$-bit bounded hash function in $\mathbb{G}_2$:  
$$
H_{2}: \{0,1\}^{12} \to \mathbb{G}_2 
$$
Again, we ``randomly'' select $4$ uniformly distributed generators $\{(7v^2 , 16v^3 ), (42v^2 , 16v^3 ), \\(17v^2 , 15v^3 ), (10v^2 , 15v^3 )\}$ from $\mathbb{G}_2$, and use the pseudo-random function from XXX\sme{add reference}. Therefore, we have to choose a set of $4$ randomly generated invertible elements from $\F_{13}$ for every generator:
$$
\begin{array}{lcl}
(7v^2 , 16v^3 ) &:& \{8,4,5,7\}\\
(42v^2 , 16v^3 ) &:& \{12,1,3,8\}\\
(17v^2 , 15v^3 ) &:& \{2,3,9,11\}\\
(10v^2 , 15v^3 ) &:& \{3,6,9,10\}
\end{array}
$$
Our hash function is then computed like this:
\begin{multline*}
H_1(x_{11},x_{10},\ldots, x_{0})=
[8\cdot 4^{x_{11}}\cdot 5^{x_{10}}\cdot 7^{x_9}](7v^2 , 16v^3)+
[12\cdot 1^{x_8}\cdot 3^{x_7}\cdot 8^{x_6}](42v^2 , 16v^3 )+\\
[2\cdot 3^{x_5}\cdot 9^{x_4}\cdot 11^{x_3}](17v^2 , 15v^3 ) +
[3\cdot 6^{x_2}\cdot 9^{x_{1}}\cdot 10^{x_{0}}](10v^2 , 15v^3 )
\end{multline*}

We extend this to a hash function that maps unbounded bitstrings to $\mathbb{G}_2$ by precomposing with an actual hash function like $MD5$, and feed the first 12 bits of its outcome into our previously defined hash function, with $TinyMD5_{\mathbb{G}_2}(s)= H_2(MD5(s)_0,\ldots MD5(s)_{11})$:
$$
TinyMD5_{\mathbb{G}_2}: \{0,1\}^* \to \mathbb{G}_2
$$
For example, since 
$MD5("")=\\ 0xd41d8cd98f00b204e9800998ecf8427e$, and the binary representation of the hexadecimal number $0x27e$ is $001001111110$, we compute $TinyMD5_{\mathbb{G}_2}$ of the empty string as follows:
$$TinyMD5_{\mathbb{G}_2}("")= H_2(MD5(s)_{11},\ldots MD5(s)_{0}) = H_2(001001111110)=$$\sme{check equation}
\end{comment}




% FOR THE SECOND VERSION OF THE BOOK
%\subsection{MNT4 MNT6 Cycles}
% https://eprint.iacr.org/2006/372.pdf theorem 5.2
%\begin{theorem}
%Let $q$ be a prime and $E/\F_q$ be an ordinary elliptic curve such that $r= |E(Fq)|$ is a prime greater than $3$.  
%\begin{itemize}
%\item $E$ has embedding  degree $k= 4$ if and only if there  exists $x\in \mathbb{Z}$ such  that $t=-x$ or $t=x+1$, and $q=x^2+x+1$.\item $E$ has  embedding  degree $k= 6$ if and only if there  exists $x\in \Z$ such that $t= 1\pm 2x$ and $q=4x^2+1$.
%\item There is an elliptic curve $E/\F_q$ with embedding degree $6$, discriminant $D$, and $|E(Fq)| = r$ if and only if there is an elliptic curve $E'/\F_r$ with embedding degree $4$, discriminant $D$, and $|E'(\F_r)| =q$.
%\end{itemize}
%\end{theorem}

%We can use this theorem to find an MNT6-MNT4 cycle over very small prime fields with characteristics $>3$: 
%\subsubsection{MNT4}
%For our MNT4 curve, we can choose $x=2$. Then $q=7$ and if we choose $t= x+1 $ then $r = q + 1 - t = 7 + 1 -3 = 5$. Therefore our MNT4 curve is a curve $y^2=x^3+ax+b$ defined over $\F_7$ that consists of $5$ points. 

%To construct the actual curve we could use the \concept{complex multiplication method} again, but since the parameters $a$ and $b$ are from $\F_7$ there are only $48$ possibilities so we simply loop through all possible $a$'s and $b$'s and count the curve points until we find a curve that has $5$ rational points. We get
%$$
%y^2 = x^3 + 4x + 1
%$$
%defined over $\F_7$, with scalar field $\F_5$. Since $7= 2^2+2+1$, we know from theorem XXX that this curve has embedding degree $4$ and hence qualifies as a pen-and-paper pairing-friendly elliptic curve. Since the curve's order is a prime and therefore has no non trivial factors, it has no non trivial subgroups. The curve has the following set of elements
%$$MNT4=\{(0,1)\to (0,6)\to (4,2)\to (4,5) \to \mathcal{O}\}$$ 
%\begin{sagecommandline}
%sage: F7 = GF(7)
%sage: MNT4 = EllipticCurve (F7,[4 ,1])
%sage: [P.xy() for P in MNT4.points() if P.order() > 1]
%\end{sagecommandline}
%The multiplication table is
%\begingroup
%    \fontsize{10pt}{10pt}\selectfont
%$$
%\begin{array}{c|ccccc}
%\cdot & \mathcal{O} & (0,1) & (4,5) & (4,2) & (0,6)\\
%\hline
%\\
%\mathcal{O} & \mathcal{O} & (0,1) & (4,5) & (4,2) & (0,6)\\
%\\
%(0,1) & (0,1) & (4,5) & (4,2) & (0,6) & \mathcal{O}\\
%\\
%(4,5) & (4,5) & (4,2) & (0,6) & \mathcal{O} & (0,1)\\
%\\
%(4,2) & (4,2) & (0,6) & \mathcal{O} & (0,1) & (4,5)\\
%\\
%(0,6) & (0,6) & \mathcal{O} & (0,1) & (4,5) & (4,2)\\
%\end{array}
%$$
%\endgroup
%In what follows we choose our generator to be $g_{MNT4}=(0,1)$.

%In what follows we want to compute type 2 pairings on our MNT4 curve. We therefore need to extract the subgroup $\mathbb{G}_1$ as well as $\mathbb{G}_2$ from the full $5$-torsion group. Since the order of MNT4 is a prime number, we already know from XXX\sme{add reference} that $\mathbb{G}_1$ is given by  $$\mathbb{G}_1=\{(0,1)\to (0,6)\to (4,2)\to (4,5) \to \mathcal{O}\}$$ 

%In type 2 pairings, the group $\mathbb{G}_2$ is defined by those elements $P$ of the full $5$-torsion group that are mapped to $7\cdot P$ under the Frobenius endomorphism XXX\sme{add reference}. Since $MNT4/\F_{7^4}$ only contains $2475$ elements, we can  loop through all elements, to find the full $5$-torsion group and extract all elements from $\mathbb{G}_2$:
%\begin{sagecommandline}
%sage: F7t.<t> = F7[]
%sage: F7_4.<u> = GF(7^4, name='u', modulus=t^4+t+1) # embedding degree is 4
%sage: MNT4 = EllipticCurve (F7_4,[4 ,1])
%sage: INF = MNT4(0) # point at infinity
%sage: for P in INF.division_points(5): # PI(P) == [q]P
%....:     if P.order() == 5: # exclude point at infinity
%....:         PiP = MNT4([a.frobenius() for a in P])
%....:         qP = 7*P
%....:         if PiP == qP:
%....:             print(P.xy())
%\end{sagecommandline}

%Choose $g_2=(2u^3 + 5u^2 + 4u + 2, 2u^3 + 3u + 5)$ as generator of $\mathbb{G}_2$, we get
%\begin{multline*}
%\mathbb{G}_2=\{ 
%(2u^3 + 5u^2 + 4u + 2, 2u^3 + 3u + 5) \to
%(5u^3 + 2u^2 + 3u + 6, 2u^2 + 3u) \to \\
%(5u^3 + 2u^2 + 3u + 6, 5u^2 + 4u) \to
%(2u^3 + 5u^2 + 4u + 2, 5u^3 + 4u + 2)\to
%\mathcal{O}\}
%\end{multline*}
%e.g. $[3]g_2= (5u^3 + 2u^2 + 3u + 6, 5u^2 + 4u)$.

%Having those groups we can do pairings. We choose the Weil pairing and use Sagemath. For example the Weil pairing between our two generators is
%$$
%e(g_1,g_2)= 5u^3 + 2u^2 + 6u
%$$
%\begin{sagecommandline}
%sage: g1 = MNT4([0,1])
%sage: g2 = MNT4(2*u^3 + 5*u^2 + 4*u + 2, 2*u^3 + 3*u + 5)
%sage: g1.weil_pairing(g2,5)
%\end{sagecommandline}
%The full pairing table can the be written as
%\begingroup
%    \fontsize{10pt}{10pt}\selectfont
    
% generate the table entries as:
% sage: for i in range(5):
% ....:     for j in range(5):
% ....:         p = (i*g1).weil_pairing((j*g2),5)
% ....:         print('e([',i,']g1,[',j,']g2=',p)         
    
    
%$$
%\begin{array}{c|lllll}
%e(\cdot,\cdot)    & \mathcal{O} & g_1            & [2]g_1         & [3]g_1         %& [4]g_1\\
%\hline
%\\
%      \mathcal{O} & 1           & 1              & 1              & 1              %& 1\\
%\\
%              g_2 & 1           & 5u^3+2u^2+6u   & 6u^3+5u^2+6    & 2u^3+u^2+2u+3  %& u^3+6u^2+6u+4\\
%\\
%\left[2\right]g_2 & 1           & 6u^3+5u^2+6    & u^3+6u^2+6u+4  & 5u^3+2u^2+6u   %& 2u^3+u^2+2u+3\\
%\\
%\left[3\right]g_2 & 1           & 2u^3+u^2+2u+3  & 5u^3+2u^2+6u   & u^3+6u^2+6u+4  & 6u^3+5u^2+6\\
%\\
%\left[4\right]g_2 & 1           & u^3+6u^2+6u+4  & 2u^3+u^2+2u+3  & 6u^3+5u^2+6    & 5u^3+2u^2+6u\\
%\end{array}
%$$
%\endgroup

%\subsubsection{MNT6}
%For our MNT6 curve, we can choose $x=1$. Then $q=5$ and if we choose $t= 1 + 2x $ then $r= q + 1 - t = 5 + 1 + 1 = 7$. Therefore our MNT6 curve is a curve $y^2=x^3+ax+b$ defined over $\F_5$ that consists of $7$ points. 

%To construct the actual curve we could use the \concept{complex multiplication method} again, but since the parameters $a$ and $b$ are from $\F_5$ there are only $24$ possibilities, we simply loop through all possible $a$'s and $b$'s and count the curve points until we find a curve that has $7$ rational points. We get
%$$
%y^2 = x^3 + 2x + 1
%$$
%defined over $\F_5$. Since $5= 4\cdot 1 + 1$, we know from theorem XXX that this curve has embedding degree $6$ and hence qualifies as a pen-and-paper pairing-friendly elliptic curve. 

%The curve has the following set of elements
%$$MNT6=\{(1,2)\to (3,3)\to (0,1)\to (0,4)\to (3,2)\to (1,3)\to \mathcal{O}\}$$
%The multiplication table is
%\begingroup
%    \fontsize{10pt}{10pt}\selectfont
%$$
%\begin{array}{c|ccccccc}
%\cdot & \mathcal{O} & (1,2) & (3,3) & (0,1) & (0,4) & (3,2) & (1,3)\\
%\hline
%\\
%\mathcal{O} & \mathcal{O} & (1,2) & (3,3) & (0,1) & (0,4) & (3,2) & (1,3)\\
%\\
%(1,2) & (1,2) & (3,3) & (0,1) & (0,4) & (3,2) & (1,3) & \mathcal{O}\\
%\\
%(3,3) & (3,3) & (0,1) & (0,4) & (3,2) & (1,3) & \mathcal{O} & (1,2)\\
%\\
%(0,1) & (0,1) & (0,4) & (3,2) & (1,3) & \mathcal{O} & (1,2) & (3,3)\\
%\\
%(0,4) & (0,4) & (3,2) & (1,3) & \mathcal{O} & (1,2) & (3,3) & (0,1)\\
%\\
%(3,2) & (3,2) & (1,3) & \mathcal{O} & (1,2) & (3,3) & (0,1) & (0,4)\\
%\\
%(1,3) & (1,3) & \mathcal{O} & (1,2) & (3,3) & (0,1) & (0,4) & (3,2)\\
%\end{array}
%$$
%\endgroup

%In what follows we choose our generator to be $g_{MNT6}=(1,2)$.

%In what follows we want to compute type 2 pairings on our MNT6 curve. We therefore need to extract the subgroup $\mathbb{G}_1$ as well as $\mathbb{G}_2$ from the full $7$-torsion group. Since the order of MNT6 is a prime number, we already know from XXX that $\mathbb{G}_1$ is given by
%$$\mathbb{G}_1=\{(1,2)\to (3,3)\to (0,1)\to (0,4)\to (3,2)\to (1,3)\to \mathcal{O}\}$$
%In type 2 pairings, the group $\mathbb{G}_2$ is defined by those elements $P$ of the full $7$-torsion group that are mapped to $5\cdot P$ under the Frobenius endomorphism XXX. Since $MNT6/\F_{5^6}$ contains $15680$ elements, we can still loop through all elements, to find the full $7$-torsion group and extract all elements from $\mathbb{G}_2$

%\begin{sagecommandline}
%sage: G.<x> = GF(5^6) # embedding degree is 6
%sage: MNT6 = EllipticCurve (G,[2 ,1])
%sage: INF = MNT6(0) # point at infinity
%sage: for P in INF.division_points(7): # PI(P) == [q]P
%....:     if P.order() == 7: # exclude point at infinity
%....:         PiP = MNT6([a.frobenius() for a in P])
%....:         qP = 5*P
%....:         if PiP == qP:
%....:             print(P.xy())
%\end{sagecommandline}

%\begin{multline*}
%\mathbb{G}_2=\{ 
%(x^3+2x^2+4x,x^5+2x^4+4x^3+3x^2+3)\to
%(x^5+4x^4+2x^3+3x^2+x+2,3x^4+2x^3+x)\to\\
%(4x^5+x^4+2x^3,3x^5+x^4+x^3+4x+4)\to
%(4x^5+x^4+2x^3,2x^5+4x^4+4x^3+x+1) \to\\
%(x^5+4x^4+2x^3+3x^2+x+2,2x^4+3x^3+4x)\to
%(x^3+2x^2+4x,4x^5+3x^4+x^3+2x^2+2)\to
%\mathcal{O}\}
%\end{multline*}
%We choose the generator $g_2 = (x^3+2x^2+4x,x^5+2x^4+4x^3+3x^2+3)$

%\begin{remark}
%Note however that our MNT6 curve discriminant $D=-16(4a^3 + 27 b^2)= -16(4\cdot 2^3 + 27\cdot 1^2)=-944$, while our MNT4 curve has discriminate XXX. Hence our example curves are not those guaranteed by theorem XXX. Those curve are both given by $y^2= x^3 + 2x +1$ over $\F_5$ and $\F_7$, respectively. However as both curves have the same defining equation, we rather choose examples that are visually distinguishable by their defining equations.




%\end{remark}
